\chapter{Crystalline Jet Calculus and Formal Differential Equations}
\label{chap:crystalline-jets}

\section{Orientation: from de Rham descent to crystalline jets}
\label{sec:jets-purpose}

Chapter~\ref{chap:spectral-de-Rham} turned finite infinitesimal variation
into an effective descent relation while keeping that relation distinct from the
universal de Rham reflection. The present chapter asks what this infinitesimal
sameness does to objects varying over a context. For every relative context
\(X\to S\), the preceding chapter constructed the relative de Rham unit and
its effective-image factorization
\[
  X\xrightarrow{e_{X/S}}Q_{X/S}
   \xrightarrow{m_{X/S}}X_{\mathrm{dR}/S},
\]
together with the effective relation groupoid
\[
  \mathcal R^{\mathrm{dR}}_{X/S,\bullet}
  \simeq
  \check C_\bullet
  \bigl(X\xrightarrow{e_{X/S}}Q_{X/S}\bigr).
\]
These two targets continue to play different roles. The universal object
\(X_{\mathrm{dR}/S}\) records the full reflection, whereas the effective image
\(Q_{X/S}\) is the quotient actually presented by points of \(X\). Jets require
effective transport, so this chapter works with the \v{C}ech groupoid of
\(X\to Q_{X/S}\); no second infinitesimal relation is introduced.

The existing infinitesimal neighbourhoods are already constructed; the remaining
problem is to make their relation act on dependent objects.
Pointwise de Rham neighbourhoods describe which generalized points are
infinitesimally related. A differential-equation calculus must additionally
transport an arbitrary object over \(X\) along those relations, compose that
transport, and recognize when an equation is already stable under it. The
dependent-adjoint triple of the effective quotient supplies exactly these
operations:
\[
  \Sigma_{e_{X/S}}
  \dashv
  e_{X/S}^*
  \dashv
  \Pi_{e_{X/S}}.
\]

Composing pullback with the two adjacent dependent adjoints produces the
two endofunctors that will carry the calculus:
\[
  \mathsf{Disk}^{\mathrm{dR},\infty}_{X/S}
  :=
  e_{X/S}^*\Sigma_{e_{X/S}},
  \qquad
  J^\infty_{X/S}
  :=
  e_{X/S}^*\Pi_{e_{X/S}}.
\]
The left composite records infinitesimal extension along the relation and the
right composite records sections over those infinitesimal variations. They are
respectively the de Rham formal-disk monad and the infinite-jet comonad, and
the dependent-adjoint triple makes them adjoint:
\[
  \mathsf{Disk}^{\mathrm{dR},\infty}_{X/S}
  \dashv
  J^\infty_{X/S},
\]
Mate formation then identifies disk actions with jet coactions. The main
structural theorem adds the geometric content: effective \v{C}ech descent
identifies either Eilenberg--Moore presentation with a single object over the
quotient:
\[
  (\mathcal H_{\Sp})_{/Q_{X/S}}
  \simeq
  \operatorname{Alg}_{\mathsf{Disk}^{\mathrm{dR},\infty}}
  ((\mathcal H_{\Sp})_{/X})
  \simeq
  \operatorname{Coalg}_{J^\infty}
  ((\mathcal H_{\Sp})_{/X}).
\]
This three-way identification is the chapter's organizing result. It turns the
same effective relation into a quotient-side descent object, a disk action, or
a jet coaction; the later differential-equation calculus uses whichever
presentation exposes the next operation most directly.

The route from the effective relation to these presentations is summarized by
\[
\begin{tikzcd}[column sep=large,row sep=large]
&
\mathcal R^{\mathrm{dR}}_{X/S,\bullet}
  \ar[dl,"{\text{left correspondence}}"']
  \ar[dr,"{\text{right correspondence}}"]
&
\\
\mathsf{Disk}^{\mathrm{dR},\infty}_{X/S}
  \ar[d,"{\text{algebras}}"']
&&
J^\infty_{X/S}
  \ar[d,"{\text{coalgebras}}"]
\\
\operatorname{Alg}_{\mathsf{Disk}^{\mathrm{dR},\infty}}
  ((\mathcal H_{\Sp})_{/X})
&&
\operatorname{Coalg}_{J^\infty}
  ((\mathcal H_{\Sp})_{/X})
\\
&
(\mathcal H_{\Sp})_{/Q_{X/S}}
  \ar[ul,"\sim"]
  \ar[ur,"\sim"']
&
\end{tikzcd}
\]
The display is a route map, not an additional layer of structure. Its lower
equivalences are conclusions: the monad and comonad are constructed first,
mate formation compares their Eilenberg--Moore objects, and effective descent
then identifies both with the quotient slice.

Throughout the chapter, \textbf{crystalline transport} means the transport
already encoded by the effective de Rham relation generated by
\[
  X\xrightarrow{e_{X/S}}Q_{X/S}.
\]
It is therefore derived rather than independently chosen: a jet coaction
transposes to relation transport, its coalgebra laws give the unit and cocycle,
relation inversion gives the inverse, and effective descent identifies the
whole structure with pullback from \(Q_{X/S}\). The word \emph{Cartan} will be
used later for a different but related role: it describes an equation morphism
that preserves this crystalline transport.

The construction sits at the intersection of two established descriptions of
differential geometry. Jet-theoretic and synthetic approaches organize
differential operators and formally integrable equations through infinite jets
\cite{Marvan1986,KhavkineSchreiber2017,
AlfonsiYoung2023}; crystal and nonlinear
\(D\)-geometric approaches organize infinitesimal compatibility by descent
over de Rham objects
\cite{GaitsgoryRozenblyumCrystals,BeilinsonDrinfeld2004}. On the present spectral theorem domain, the square-zero-generated effective de
Rham relation simultaneously produces the jet calculus over \(X\) and the
quotient-side descent object over \(Q_{X/S}\).

The chapter follows that identification in three phases. First, relative de
Rham neighbourhoods are promoted from pointwise disks to functorial disk and
jet endofunctors, and substitution is fixed before any parameterized
construction is used. Second, mate formation and effective descent identify
disk actions, jet coactions, and crystalline transport; the resulting comonad
then supplies holonomic sections, differential operators, generalized
equations, solution objects, and Cartan prolongation. Third, base change,
formally \(\acute{e}\)tale descent, and translation models expose the geometric
content of the calculus. The final section keeps two boundaries separate:
reduced-classical comparison is a convergence problem for a further
localization, whereas finite order requires a new global filtration of the
intrinsic relation.

The conventions on constructed coherence from the preceding chapters remain in
force. In particular, once an adjunction, Beck--Chevalley comparison, descent
limit, or relation groupoid has supplied the relevant higher structure, later
prose uses that structure without reopening its coherence. Likewise, an
\emph{intrinsic formally integrable PDE} will mean an object already carrying
effective de Rham descent, whereas a \emph{generalized jet equation} will mean
an ambient map into an infinite jet object that need not yet carry that
descent. The intrinsic infinite jet functor \(J^\infty\) is constructed here;
a global tower \(J^k\) is not used until a compatible finite-order
globalization has actually been constructed.

\section{Formal disks and the infinite jet adjunction}
\label{sec:jets-formal-disk-adjunction}

The effective relation is already available; this section turns it into an
action on the whole slice over \(X\). We first fix the dependent slice
semantics and identify the pointwise formal disks. We then promote those disks
to the formal-disk monad and infinite-jet comonad, identify their adjunction,
and establish substitution. The reader should leave the section retaining the
disk--jet adjunction and its fibrewise interpretation, not the intermediate
bookkeeping used to construct them.

For orientation, the pointwise and functorial pictures fit into the following
pullback architecture. For \(a:T\to X\) and \(E\to X\), the middle square is
the formal disk over \(a\); the upper square pulls the coefficient object to
that disk. Dependent product down the left column is the fibrewise operation
that later computes \(a^*J^\infty E\).
\[
\begin{tikzcd}[column sep=huge,row sep=large]
\operatorname{ev}_a^*E
  \ar[r]
  \ar[d]
&
E
  \ar[d]
\\
T\times_{Q_{X/S}}X
  \ar[r,"\operatorname{ev}_a"]
  \ar[d,"\pi_a"']
&
X
  \ar[d,"e_{X/S}"]
\\
T
  \ar[r,"e_{X/S}a"']
&
Q_{X/S}.
\end{tikzcd}
\]
\subsection{Slice semantics and effective \v{C}ech descent}
\label{sec:jets-ambient-slices}

The effective quotient acts on coefficients through the slice semantics of the
ambient spectral Zariski \(\infty\)-topos. We fix those semantics once because
every later jet comparison is inherited from the same dependent adjoints,
Beck--Chevalley maps, and effective \v{C}ech descent.

Fix a morphism
\[
  p:X\longrightarrow S
\]
in the spectral Zariski \(\infty\)-topos.

For every map \(f:U\to V\) in the ambient infinity-topos we use the dependent
adjoint triple
\[
  \Sigma_f\dashv f^*\dashv\Pi_f.
\]
Equivalently, the three functors fit into
\[
\begin{tikzcd}[column sep=huge]
(\mathcal H_{\Sp})_{/V}
  \ar[r,shift left=1.4ex,"f^*"]
&
(\mathcal H_{\Sp})_{/U}
  \ar[l,shift left=1.4ex,"\Sigma_f"]
  \ar[l,shift right=1.4ex,"\Pi_f"']
\end{tikzcd}
\]
with the indicated adjunctions. Here \(\Sigma_f\) is composition with \(f\),
\(f^*\) is pullback, and \(\Pi_f\) is dependent product. These adjoints, together
with their Beck--Chevalley transformations, are the only slice-theoretic input
used below; see \cite{LurieHTT}.

\begin{theorem}[Slice descent and dependent adjoints]
\label{thm:parameterized-slice-cech-descent}
The contravariant slice construction
\[
  T\longmapsto(\mathcal H_{\Sp})_{/T}
\]
defines a functor
\[
  \mathcal H_{\Sp}^{\op}\longrightarrow\Pr^L
\]
whose transition functor along \(f:U\to V\) is pullback
\[
  f^*:(\mathcal H_{\Sp})_{/V}\to(\mathcal H_{\Sp})_{/U}.
\]
For every such \(f\), there is a dependent-adjoint triple
\[
  \Sigma_f\dashv f^*\dashv\Pi_f,
\]
and the associated Beck--Chevalley transformations are compatible with
Cartesian pasting.

Moreover, the slice construction carries colimits in \(\mathcal H_{\Sp}\) to
limits of infinity-categories. In particular, if
\[
  r:U\twoheadrightarrow T
\]
is an effective epimorphism with \v{C}ech nerve \(U_\bullet\), then restriction
induces a canonical equivalence
\[
  (\mathcal H_{\Sp})_{/T}
  \xrightarrow{\simeq}
  \operatorname*{Tot}_{[n]\in\Delta}
  (\mathcal H_{\Sp})_{/U_n}.
\]
All descent and transition maps used below are obtained from this
contravariant slice construction.
\end{theorem}

\begin{proof}
\leavevmode
\begin{enumerate}[label=\textup{(\arabic*)}]
    \item \textbf{Apply van Kampen slice descent.}
Colimits in an infinity-topos are van Kampen, so for every small colimit
presentation
\[
  T\simeq\operatorname*{colim}_{i\in I}T_i
\]
the canonical restriction functor is an equivalence
\[
  (\mathcal H_{\Sp})_{/T}
  \xrightarrow{\simeq}
  \operatorname*{lim}_{i\in I^{\op}}
  (\mathcal H_{\Sp})_{/T_i};
\]
see \cite[Theorem~6.1.3.9]{LurieHTT}.

    \item \textbf{Construct the dependent adjoints.}
Universality of colimits makes each
pullback functor \(f^*\) colimit-preserving, while local Cartesian closure gives
the dependent adjoints
\[
  \Sigma_f\dashv f^*\dashv\Pi_f.
\]

    \item \textbf{Record Beck--Chevalley coherence.}
Their Beck--Chevalley transformations and higher coherence are the canonical
ones supplied by Cartesian pasting and mate formation.

    \item \textbf{Present an effective epimorphism by its Čech nerve.}
If \(r:U\twoheadrightarrow T\) is effective, then
\[
  T\simeq
  \operatorname*{colim}_{[n]\in\Delta^{\op}}U_n.
\]

    \item \textbf{Identify effective Čech descent.}
Applying the van Kampen equivalence gives
\[
  (\mathcal H_{\Sp})_{/T}
  \simeq
  \operatorname*{lim}_{[n]\in\Delta}
  (\mathcal H_{\Sp})_{/U_n},
\]
which is the asserted effective \v{C}ech descent equivalence.
\end{enumerate}
\end{proof}

Let
\[
  \mathsf u:X\to\mathcal R^{\mathrm{dR}}_{X/S},
  \qquad
  \mathsf{inv}:
  \mathcal R^{\mathrm{dR}}_{X/S}
  \xrightarrow{\simeq}
  \mathcal R^{\mathrm{dR}}_{X/S}
\]
denote the unit and inversion maps of the de Rham relation groupoid. Its
first two nontrivial simplicial degrees are displayed by
\[
\begin{tikzcd}[column sep=huge]
\mathcal R^{\mathrm{dR}}_{X/S,2}
  \ar[r,shift left=2.2ex,"d_0"]
  \ar[r,"d_1" description]
  \ar[r,shift right=2.2ex,"d_2"']
&
\mathcal R^{\mathrm{dR}}_{X/S}
  \ar[r,shift left=1.2ex,"\mathsf s"]
  \ar[r,shift right=1.2ex,"\mathsf t"']
&
X .
\end{tikzcd}
\]

The slice machinery is now fixed. The next step is geometric: identify what the
effective relation looks like when centred at a generalized point.

\subsection{Formal disks as pullbacks of the de Rham relation}
\label{sec:formal-disks}

Two points of \(X\) are related exactly when they have the same effective de
Rham image. Pulling the source of that relation back along
\(a:T\to X\) therefore isolates the infinitesimal variations centred at
\(a\). The next theorem shows that this pullback is precisely the relative de
Rham neighbourhood of the graph of \(a\); no additional notion of formal disk
is introduced.

Let \(a:T\to X\) be an object of \((\mathcal H_{\Sp})_{/X}\), and let
\[
  \Gamma_a:T\longrightarrow T\times_SX
\]
be its graph. We use the relative de Rham neighbourhood along \(\Gamma_a\) as
the pointwise formal disk
\[
  \operatorname{Nbh}^{\mathrm{dR}}_{T/S}(T\times_SX)
\]
taken along \(\Gamma_a\). Write
\[
  \pi_a:
  \operatorname{Nbh}^{\mathrm{dR}}_{T/S}(T\times_SX)\to T,
  \qquad
  \operatorname{ev}_a:
  \operatorname{Nbh}^{\mathrm{dR}}_{T/S}(T\times_SX)\to X
\]
for the maps induced by the two projections of \(T\times_SX\).

\begin{theorem}[Effective presentation of the de Rham formal disk]
\label{thm:formal-disk-graph-neighbourhood}
\label{prop:disk-effective-quotient}
There are canonical equivalences
\[
  \operatorname{Nbh}^{\mathrm{dR}}_{T/S}(T\times_SX)
  \simeq
  T\times_{X_{\mathrm{dR}/S}}X
  \simeq
  T\times_{Q_{X/S}}X
  \simeq
  T\times_{X,\mathsf s}\mathcal R^{\mathrm{dR}}_{X/S}.
\]
Under the final equivalence \(\pi_a\) is the pullback of \(\mathsf s\), and
\(\operatorname{ev}_a\) is induced by \(\mathsf t\). In particular \(\pi_a\)
is an effective epimorphism.
The relation presentation is summarized by the Cartesian square
\[
\begin{tikzcd}[column sep=huge,row sep=large]
T\times_{X,\mathsf s}\mathcal R^{\mathrm{dR}}_{X/S}
  \ar[r]
  \ar[d,"\pi_a"']
&
\mathcal R^{\mathrm{dR}}_{X/S}
  \ar[d,"\mathsf s"]\\
T \ar[r,"a"']
&
X .
\end{tikzcd}
\]
Evaluation is obtained by following the upper horizontal map with \(\mathsf t\).
\end{theorem}

\begin{proof}
\leavevmode
\begin{enumerate}[label=\textup{(\arabic*)}]
    \item \textbf{Expand the relative de Rham neighbourhood.}
By definition,
\[
  \operatorname{Nbh}^{\mathrm{dR}}_{T/S}(T\times_SX)
  =
  (T\times_SX)
  \times_{(T\times_SX)_{\mathrm{dR}/S}}
  T_{\mathrm{dR}/S}.
\]

    \item \textbf{Apply left exactness.}
Left exactness of relative de Rham reflection gives
\[
  (T\times_SX)_{\mathrm{dR}/S}
  \simeq
  T_{\mathrm{dR}/S}\times_SX_{\mathrm{dR}/S}.
\]

    \item \textbf{Identify the reflected graph map.}
Under this equivalence the map from \(T_{\mathrm{dR}/S}\) is
\[
  (\id,a_{\mathrm{dR}/S}).
\]

    \item \textbf{Apply naturality and finite-limit pasting.}
Naturality of the de Rham unit gives
\[
  a_{\mathrm{dR}/S}\eta_{T/S}^{\mathrm{dR}}
  \simeq
  \eta_{X/S}^{\mathrm{dR}}a,
\]
so finite-limit pasting gives
\[
  \operatorname{Nbh}^{\mathrm{dR}}_{T/S}(T\times_SX)
  \simeq
  T\times_{X_{\mathrm{dR}/S}}X.
\]

    \item \textbf{Use the effective-image monomorphism.}
Since
\[
  m_{X/S}:Q_{X/S}\longrightarrow X_{\mathrm{dR}/S}
\]
is a monomorphism,
\[
  Q_{X/S}\times_{X_{\mathrm{dR}/S}}Q_{X/S}
  \simeq
  Q_{X/S}.
\]

    \item \textbf{Replace the reflected pullback by the effective quotient.}
Both maps into \(X_{\mathrm{dR}/S}\) factor through \(Q_{X/S}\), hence
\[
  T\times_{X_{\mathrm{dR}/S}}X
  \simeq
  T\times_{Q_{X/S}}X.
\]

    \item \textbf{Identify the relation presentation.}
Finally,
\[
  \mathcal R^{\mathrm{dR}}_{X/S}
  \simeq
  X\times_{Q_{X/S}}X
\]
gives the relation presentation.

    \item \textbf{Identify the structural maps and effectivity.}
The description of \(\pi_a\) and
\(\operatorname{ev}_a\) follows from the two relation projections, and
effectivity of \(\pi_a\) follows by base change from
\(e_{X/S}:X\twoheadrightarrow Q_{X/S}\).
\end{enumerate}
\end{proof}

Thus the pointwise formal disk has been identified intrinsically as the
effective pullback \(T\times_{Q_{X/S}}X\), with projection to the parameter and
evaluation in \(X\). What is still missing is functoriality in an arbitrary
object over \(X\).

\subsection{Jet endofunctors from the effective quotient}
\label{sec:jet-endofunctors}

The dependent-adjoint triple of \(e_{X/S}:X\to Q_{X/S}\) packages these
pointwise disks functorially. Pullback composed with the left adjoint gives the
formal-disk endofunctor; pullback composed with the right adjoint gives the jet
endofunctor. The remaining results in this subsection verify that these two
endofunctors are exactly the left and right correspondence actions of the same
\v{C}ech relation.

\subsubsection{Endofunctors from the effective quotient}

\begin{definition}[De Rham formal-disk and jet endofunctors]
Define endofunctors of \((\mathcal H_{\Sp})_{/X}\) by
\[
  \mathsf{Disk}^{\mathrm{dR},\infty}_{X/S}
  :=
  e_{X/S}^*\Sigma_{e_{X/S}},
  \qquad
  J^\infty_{X/S}
  :=
  e_{X/S}^*\Pi_{e_{X/S}}.
\]
For \(E\to X\), the object \(J^\infty_{X/S}E\to X\) is the relative infinite
jet object.
\end{definition}

For a fixed relative context \(X/S\), we suppress the subscript \(X/S\) from
\(\mathsf{Disk}^{\mathrm{dR},\infty}_{X/S}\) and \(J^\infty_{X/S}\) whenever
the base is already active.

\begin{theorem}[Effective-image formulas]
\label{thm:effective-image-jets}
There are canonical equivalences
\[
  \mathsf{Disk}^{\mathrm{dR},\infty}_{X/S}
  \simeq
  (\eta_{X/S}^{\mathrm{dR}})^*
  \Sigma_{\eta_{X/S}^{\mathrm{dR}}},
  \qquad
  J^\infty_{X/S}
  \simeq
  (\eta_{X/S}^{\mathrm{dR}})^*
  \Pi_{\eta_{X/S}^{\mathrm{dR}}}.
\]
These equivalences identify the monad and comonad structures constructed below.
Hence the de Rham formal-disk and jet calculus depends only on the effective de Rham quotient
\(X\to Q_{X/S}\).
\end{theorem}

\begin{proof}
\leavevmode
\begin{enumerate}[label=\textup{(\arabic*)}]
    \item \textbf{Record the effective-image factorization.}
The effective-image factorization is the commutative triangle
\[
\begin{tikzcd}[column sep=huge,row sep=large]
X \ar[rr,"\eta_{X/S}^{\mathrm{dR}}"]
  \ar[dr,"e_{X/S}"']
&& X_{\mathrm{dR}/S}\\
& Q_{X/S} \ar[ur,hook,"m_{X/S}"'] &
\end{tikzcd}
\]
with \(\eta_{X/S}^{\mathrm{dR}}=m_{X/S}e_{X/S}\).

    \item \textbf{Contract the dependent sum along the monomorphism.}
The monomorphism
\(m_{X/S}:Q_{X/S}\to X_{\mathrm{dR}/S}\) has invertible diagonal, and hence
\[
  m_{X/S}^*\Sigma_{m_{X/S}}
  \simeq
  \id_{(\mathcal H_{\Sp})_{/Q_{X/S}}}.
\]

    \item \textbf{State the dependent-product contraction.}
Likewise
\[
  m_{X/S}^*\Pi_{m_{X/S}}
  \simeq
  \id_{(\mathcal H_{\Sp})_{/Q_{X/S}}}.
\]

    \item \textbf{Compute the representing mapping spaces.}
Indeed, for \(Z,E\in(\mathcal H_{\Sp})_{/Q_{X/S}}\),
\[
\begin{aligned}
  \Map_{Q_{X/S}}
  \bigl(Z,m_{X/S}^*\Pi_{m_{X/S}}E\bigr)
  &\simeq
  \Map_{X_{\mathrm{dR}/S}}
  \bigl(\Sigma_{m_{X/S}}Z,\Pi_{m_{X/S}}E\bigr)\\
  &\simeq
  \Map_{Q_{X/S}}
  \bigl(m_{X/S}^*\Sigma_{m_{X/S}}Z,E\bigr)\\
  &\simeq
  \Map_{Q_{X/S}}(Z,E).
\end{aligned}
\]

    \item \textbf{Apply Yoneda.}
Yoneda gives the second equivalence.

    \item \textbf{Compose the dependent adjunctions.}
Composition of dependent sums and
products now gives
\[
\begin{aligned}
  (\eta_{X/S}^{\mathrm{dR}})^*
  \Sigma_{\eta_{X/S}^{\mathrm{dR}}}
  &\simeq
  e_{X/S}^*
  m_{X/S}^*
  \Sigma_{m_{X/S}}
  \Sigma_{e_{X/S}}\\
  &\simeq
  e_{X/S}^*\Sigma_{e_{X/S}},\\[0.4em]
  (\eta_{X/S}^{\mathrm{dR}})^*
  \Pi_{\eta_{X/S}^{\mathrm{dR}}}
  &\simeq
  e_{X/S}^*
  m_{X/S}^*
  \Pi_{m_{X/S}}
  \Pi_{e_{X/S}}\\
  &\simeq
  e_{X/S}^*\Pi_{e_{X/S}}.
\end{aligned}
\]

    \item \textbf{Transport the monad and comonad structure.}
Because the comparisons come from equivalences of the composing adjunctions,
they identify the induced monad and comonad structures.
\end{enumerate}
\end{proof}

The effective-image calculation removes the universal reflected target from the
active jet data: from this point through the jet construction, the map
\(e_{X/S}:X\to Q_{X/S}\) and its \v{C}ech relation are sufficient. The next
formula rewrites the two endofunctors directly as left and right actions of
that relation.

\begin{proposition}[Relation-correspondence formulas]
\label{prop:relation-jet-formulas}
There are canonical equivalences
\[
  \mathsf{Disk}^{\mathrm{dR},\infty}_{X/S}
  \simeq
  \Sigma_{\mathsf s}\mathsf t^*,
  \qquad
  J^\infty_{X/S}
  \simeq
  \Pi_{\mathsf s}\mathsf t^*.
\]
Transport along inversion also gives
\[
  \mathsf{Disk}^{\mathrm{dR},\infty}_{X/S}
  \simeq
  \Sigma_{\mathsf t}\mathsf s^*.
\]
\end{proposition}

\begin{proof}
\leavevmode
\begin{enumerate}[label=\textup{(\arabic*)}]
    \item \textbf{Identify the relation square as Cartesian.}
The square
\[
\begin{tikzcd}
\mathcal R^{\mathrm{dR}}_{X/S}
  \ar[r,"\mathsf t"]
  \ar[d,"\mathsf s"']
&
X \ar[d,"e_{X/S}"]
\\
X \ar[r,"e_{X/S}"']
&
Q_{X/S}
\end{tikzcd}
\]
is Cartesian.

    \item \textbf{Apply Beck--Chevalley.}
Left and right Beck--Chevalley give
\[
  e_{X/S}^*\Sigma_{e_{X/S}}
  \simeq
  \Sigma_{\mathsf s}\mathsf t^*,
  \qquad
  e_{X/S}^*\Pi_{e_{X/S}}
  \simeq
  \Pi_{\mathsf s}\mathsf t^*.
\]

    \item \textbf{Use relation inversion.}
Relation inversion exchanges \(\mathsf s\) and \(\mathsf t\), yielding the
final formula.
\end{enumerate}
\end{proof}

The degree-one correspondence is the conceptual interface:
\(\mathsf{Disk}^{\mathrm{dR},\infty}\) and \(J^\infty\) are the left and
right dependent actions of the same relation. The iterated formulas below have
one purpose: they verify that relation composition and units induce the
monad/comonad structure coherently. The individual vertex notation need not be
retained after that verification.

\subsubsection{The simplicial relation and iterated correspondences}

\begin{proposition}[Iterated relation correspondences]
\label{prop:iterated-relation-correspondences}
For every integer \(n\geq0\), let
\[
  \mathsf t_n:\mathcal R^{\mathrm{dR}}_{X/S,n}\to X
\]
denote the last-vertex projection. Thus
\[
  \mathsf t_n(x_0,\ldots,x_n)=x_n.
\]
For \(n=0\),
\[
  v_0=\mathsf t_0=\id_X,
\]
and in degree one
\[
  v_1=\mathsf s,
  \qquad
  \mathsf t_1=\mathsf t.
\]

There are canonical equivalences
\[
  (\mathsf{Disk}^{\mathrm{dR},\infty})^n
  \simeq
  \Sigma_{v_n}\mathsf t_n^*,
  \qquad
  (J^\infty)^n
  \simeq
  \Pi_{v_n}\mathsf t_n^*,
\]
compatible in \(n\) with the Cartesian pasting maps generated by insertion of
the relation unit and composition of adjacent relation arrows.

Let
\[
  \operatorname{rev}_n:\mathcal R^{\mathrm{dR}}_{X/S,n}\xrightarrow{\simeq}\mathcal R^{\mathrm{dR}}_{X/S,n},
  \qquad
  (x_0,\ldots,x_n)\longmapsto(x_n,\ldots,x_0)
\]
be reversal. Then
\[
  v_n\operatorname{rev}_n=\mathsf t_n,
  \qquad
  \mathsf t_n\operatorname{rev}_n=v_n.
\]
Transport along \(\operatorname{rev}_n\) identifies the right adjoint of
\(\Sigma_{v_n}\mathsf t_n^*\) with
\(\Pi_{v_n}\mathsf t_n^*\). Hence
\[
  (\mathsf{Disk}^{\mathrm{dR},\infty})^n\dashv(J^\infty)^n
\]
for every \(n\). Under these adjunctions the simplicial structure maps
inducing the monad structure on \(\mathsf{Disk}^{\mathrm{dR},\infty}\) are mates of the corresponding
simplicial structure maps inducing the comonad structure on \(J^\infty\).
In particular, the unit and multiplication of \(\mathsf{Disk}^{\mathrm{dR},\infty}\) are mates of the
counit and comultiplication of \(J^\infty\), and the simplicial identities of
\(\mathcal R^{\mathrm{dR}}_{X/S,\bullet}\) supply the remaining compatibility.
\end{proposition}

\begin{proof}
\leavevmode
\begin{enumerate}[label=\textup{(\arabic*)}]
    \item \textbf{Check the initial degrees.}
The assertion for \(n=0\) is tautological. For \(n=1\), the identities
\[
  v_1=\mathsf s,
  \qquad
  \mathsf t_1=\mathsf t
\]
reduce the statement to Proposition~\ref{prop:relation-jet-formulas}.

    \item \textbf{Form the induction square.}
Suppose the formulas have been constructed in degree \(n\). Form the
Cartesian square
\[
\begin{tikzcd}
\mathcal R^{\mathrm{dR}}_{X/S,n+1} \ar[r,"\pi_{1\ldots n+1}"] \ar[d,"\pi_{01}"']
  & \mathcal R^{\mathrm{dR}}_{X/S,n} \ar[d,"v_n"]\\
\mathcal R^{\mathrm{dR}}_{X/S,1} \ar[r,"\mathsf t"']
  & X,
\end{tikzcd}
\]
where
\[
  \pi_{01}(x_0,\ldots,x_{n+1})=(x_0,x_1),
  \qquad
  \pi_{1\ldots n+1}(x_0,\ldots,x_{n+1})=(x_1,\ldots,x_{n+1}).
\]

    \item \textbf{Identify the first and last vertices.}
By definition of the first and last vertices,
\[
  v_{n+1}=\mathsf s\,\pi_{01},
  \qquad
  \mathsf t_{n+1}=\mathsf t_n\pi_{1\ldots n+1}.
\]

    \item \textbf{Apply left Beck--Chevalley.}
Left Beck--Chevalley gives
\[
  \mathsf t^*\Sigma_{v_n}
  \simeq
  \Sigma_{\pi_{01}}\pi_{1\ldots n+1}^*.
\]
Therefore
\[
\begin{aligned}
  \mathsf{Disk}^{\mathrm{dR},\infty}(\mathsf{Disk}^{\mathrm{dR},\infty})^n
  &\simeq
  \Sigma_{\mathsf s}\mathsf t^*
  \Sigma_{v_n}\mathsf t_n^*\\
  &\simeq
  \Sigma_{\mathsf s}\Sigma_{\pi_{01}}\pi_{1\ldots n+1}^*\mathsf t_n^*\\
  &\simeq
  \Sigma_{v_{n+1}}\mathsf t_{n+1}^*.
\end{aligned}
\]

    \item \textbf{Apply right Beck--Chevalley.}
Right Beck--Chevalley similarly gives
\[
  \mathsf t^*\Pi_{v_n}
  \simeq
  \Pi_{\pi_{01}}\pi_{1\ldots n+1}^*,
\]
and hence
\[
\begin{aligned}
  J^\infty(J^\infty)^n
  &\simeq
  \Pi_{\mathsf s}\mathsf t^*
  \Pi_{v_n}\mathsf t_n^*\\
  &\simeq
  \Pi_{\mathsf s}\Pi_{\pi_{01}}\pi_{1\ldots n+1}^*\mathsf t_n^*\\
  &\simeq
  \Pi_{v_{n+1}}\mathsf t_{n+1}^*.
\end{aligned}
\]

    \item \textbf{Complete the inductive correspondence.}
This proves the iterated correspondence formulas.

    \item \textbf{Identify the formal right adjoint.}
The right adjoint of
\[
  \Sigma_{v_n}\mathsf t_n^*
\]
is
\[
  \Pi_{\mathsf t_n}v_n^*.
\]

    \item \textbf{Record the reversal identities.}
Reversal satisfies
\[
  v_n\operatorname{rev}_n=\mathsf t_n,
  \qquad
  \mathsf t_n\operatorname{rev}_n=v_n,
\]
and \(\operatorname{rev}_n\) is an equivalence.

    \item \textbf{Transport the right adjoint through reversal.}
Composition of dependent products and
transport along \(\operatorname{rev}_n\) therefore give
\[
  \Pi_{\mathsf t_n}v_n^*
  \simeq
  \Pi_{v_n}\mathsf t_n^*.
\]
Thus
\[
  (\mathsf{Disk}^{\mathrm{dR},\infty})^n\dashv(J^\infty)^n.
\]

    \item \textbf{Match the degree-one structure maps.}
In degree one, the relation unit \(\mathsf u:X\to\mathcal R^{\mathrm{dR}}_{X/S,1}\) induces
the monad unit by the dependent-sum unit and the comonad counit by the
dependent-product counit. These transformations are mates.

    \item \textbf{Match the degree-two structure maps.}
In degree two, the middle face
\[
  d_1:\mathcal R^{\mathrm{dR}}_{X/S,2}\to\mathcal R^{\mathrm{dR}}_{X/S,1},
  \qquad
  (x_0,x_1,x_2)\longmapsto(x_0,x_2),
\]
induces the monad multiplication by contraction with the counit of
\[
  \Sigma_{d_1}\dashv d_1^*,
\]
while it induces the comonad comultiplication by insertion of the unit of
\[
  d_1^*\dashv\Pi_{d_1}.
\]
These transformations are mates under the displayed iterated adjunctions.

    \item \textbf{Recover the higher coherences.}
In higher degrees, all structural composites are obtained by deleting or
repeating vertices of the same \v{C}ech simplex. Beck--Chevalley and mate formation are compatible with Cartesian pasting, so
the two systems of transformations agree. The simplicial identities then
supply the remaining monad--comonad compatibility.
\end{enumerate}
\end{proof}

The iterated correspondence has now done its job: the simplicial operators of
the \v{C}ech relation are identified with the structural maps of the disk
monad and jet comonad. The next theorem extracts the pointwise consequence
that will actually be used later.

\begin{theorem}[Fibrewise interpretation of jets]
\label{thm:fibrewise-jets}
For \(a:T\to X\) and \(E\to X\), there is a canonical equivalence in the
slice over \(T\)
\[
  a^*J^\infty_{X/S}E
  \simeq
  \Pi_{\pi_a}\operatorname{ev}_a^*E.
\]
Thus a generalized fibre of the jet object is the object of sections of \(E\)
over the intrinsic de Rham formal disk centred at that generalized point.
\end{theorem}

\begin{proof}
\leavevmode
\begin{enumerate}[label=\textup{(\arabic*)}]
    \item \textbf{Compute the pulled-back jet.}
Pull back the formula
\[
  J^\infty\simeq\Pi_{\mathsf s}\mathsf t^*
\]
along \(a\). Proposition~\ref{thm:formal-disk-graph-neighbourhood} identifies the
pullback of \(\mathsf s\) with \(\pi_a\), and right Beck--Chevalley gives the
result.
\end{enumerate}
\end{proof}

\subsubsection{Monad, comonad, and the disk--jet adjunction}

The correspondence calculus already determines the structure maps
geometrically. We now record the same structure in adjunction notation because
later mate and Eilenberg--Moore arguments use the units and counits
explicitly.

For the adjoint triple above, let
\[
  \eta^\Sigma:
  \id_{(\mathcal H_{\Sp})_{/X}}
  \longrightarrow
  e_{X/S}^*\Sigma_{e_{X/S}},
  \qquad
  \varepsilon^\Sigma:
  \Sigma_{e_{X/S}}e_{X/S}^*
  \longrightarrow
  \id_{(\mathcal H_{\Sp})_{/Q_{X/S}}}
\]
be the unit and counit of
\[
  \Sigma_{e_{X/S}}\dashv e_{X/S}^*,
\]
and let
\[
  \eta^\Pi:
  \id_{(\mathcal H_{\Sp})_{/Q_{X/S}}}
  \longrightarrow
  \Pi_{e_{X/S}}e_{X/S}^*,
  \qquad
  \varepsilon^\Pi:
  e_{X/S}^*\Pi_{e_{X/S}}
  \longrightarrow
  \id_{(\mathcal H_{\Sp})_{/X}}
\]
be the unit and counit of
\[
  e_{X/S}^*\dashv\Pi_{e_{X/S}}.
\]

\begin{proposition}[De Rham formal-disk monad]
The endofunctor \(\mathsf{Disk}^{\mathrm{dR},\infty}\) carries the canonical monad structure
\[
  (\mathsf{Disk}^{\mathrm{dR},\infty},\eta^\Sigma,\mu),
  \qquad
  \mu:=e_{X/S}^*\varepsilon^\Sigma_{\Sigma_{e_{X/S}}}:
  \mathsf{Disk}^{\mathrm{dR},\infty} \mathsf{Disk}^{\mathrm{dR},\infty}\to \mathsf{Disk}^{\mathrm{dR},\infty}.
\]
Under the relation formula, the unit is induced by \(\mathsf u:X\to\mathcal R^{\mathrm{dR}}_{X/S}\) and the
multiplication by
\[
  \mathrm{comp}:\mathcal R^{\mathrm{dR}}_{X/S}\times_{\mathsf t,X,\mathsf s}\mathcal R^{\mathrm{dR}}_{X/S}\to\mathcal R^{\mathrm{dR}}_{X/S}.
\]
\end{proposition}

\begin{proof}
\leavevmode
\begin{enumerate}[label=\textup{(\arabic*)}]
    \item \textbf{Identify the adjunction monad.}
The monad is the monad of the adjunction
\(\Sigma_{e_{X/S}}\dashv e_{X/S}^*\).

    \item \textbf{Translate the monad structure to the relation.}
Under the \v{C}ech correspondence, its unit is
the diagonal and its multiplication is relation composition.

    \item \textbf{Verify the monad laws.}
The monad laws are equivalently the unit and associativity identities of the
relation groupoid.
\end{enumerate}
\end{proof}

\begin{theorem}[Jet comonad]
\label{thm:jet-comonad}
The endofunctor \(J^\infty\) carries the canonical comonad structure
\[
  (J^\infty,\varepsilon^\Pi,\delta^J),
  \qquad
  \delta^J:=e_{X/S}^*\eta^\Pi_{\Pi_{e_{X/S}}}:
  J^\infty\to J^\infty J^\infty.
\]
Thus
\[
  \varepsilon^\Pi_E:J^\infty E\to E,
  \qquad
  \delta^J_E:J^\infty E\to J^\infty J^\infty E
\]
satisfy the counit and coassociativity laws supplied by the adjunction.
\end{theorem}

\begin{theorem}[Adjoint-triple mate calculus]
\label{thm:adjoint-triple-mates}
The adjoint triple
\[
  \Sigma_{e_{X/S}}\dashv e_{X/S}^*\dashv\Pi_{e_{X/S}}
\]
canonically induces an adjunction
\[
  \mathsf{Disk}^{\mathrm{dR},\infty}=e_{X/S}^*\Sigma_{e_{X/S}}
  \dashv
  e_{X/S}^*\Pi_{e_{X/S}}=J^\infty.
\]
For \(Z,E\in(\mathcal H_{\Sp})_{/X}\),
\[
\begin{aligned}
  \Map_X(\mathsf{Disk}^{\mathrm{dR},\infty} Z,E)
  &=
  \Map_X(e_{X/S}^*\Sigma_{e_{X/S}} Z,E)\\
  &\simeq
  \Map_{Q_{X/S}}(\Sigma_{e_{X/S}} Z,\Pi_{e_{X/S}} E)\\
  &\simeq
  \Map_X(Z,e_{X/S}^*\Pi_{e_{X/S}} E)\\
  &=
  \Map_X(Z,J^\infty E).
\end{aligned}
\]
Its unit is
\[
  \eta^{\mathrm{adj}}_Z:
  Z
  \xrightarrow{\eta^\Sigma_Z}
  e_{X/S}^*\Sigma_{e_{X/S}} Z
  \xrightarrow{e_{X/S}^*\eta^\Pi_{\Sigma_{e_{X/S}} Z}}
  e_{X/S}^*\Pi_{e_{X/S}}e_{X/S}^*\Sigma_{e_{X/S}} Z
  =
  J^\infty \mathsf{Disk}^{\mathrm{dR},\infty} Z,
\]
and its counit is
\[
  \varepsilon^{\mathrm{adj}}_E:
  \mathsf{Disk}^{\mathrm{dR},\infty} J^\infty E
  =
  e_{X/S}^*\Sigma_{e_{X/S}}e_{X/S}^*\Pi_{e_{X/S}} E
  \xrightarrow{e_{X/S}^*\varepsilon^\Sigma_{\Pi_{e_{X/S}} E}}
  e_{X/S}^*\Pi_{e_{X/S}} E
  \xrightarrow{\varepsilon^\Pi_E}
  E.
\]

Under mate correspondence for this composite adjunction,
the monad unit
\[
  \eta^\Sigma:\id\to \mathsf{Disk}^{\mathrm{dR},\infty}
\]
and the comonad counit
\[
  \varepsilon^\Pi:J^\infty\to\id
\]
are mates. Likewise the monad multiplication
\[
  \mu:\mathsf{Disk}^{\mathrm{dR},\infty} \mathsf{Disk}^{\mathrm{dR},\infty}\to \mathsf{Disk}^{\mathrm{dR},\infty}
\]
and the comonad comultiplication
\[
  \delta^J:J^\infty\to J^\infty J^\infty
\]
are mates with respect to the induced adjunctions
\[
  (\mathsf{Disk}^{\mathrm{dR},\infty})^2\dashv(J^\infty)^2
  \qquad\text{and}\qquad
  \mathsf{Disk}^{\mathrm{dR},\infty}\dashv J^\infty.
\]
These mate identifications are compatible with pasting.

\end{theorem}

\begin{proof}
\leavevmode
\begin{enumerate}[label=\textup{(\arabic*)}]
    \item \textbf{Compose the two adjunctions.}
The mapping-space equivalence is the composite of
\[
  \Sigma_{e_{X/S}}\dashv e_{X/S}^*
  \qquad\text{and}\qquad
  e_{X/S}^*\dashv\Pi_{e_{X/S}}.
\]
The displayed formulas for \(\eta^{\mathrm{adj}}\) and \(\varepsilon^{\mathrm{adj}}\) are the unit and counit of
the composite adjunction.

    \item \textbf{Take homotopy-coherent mates.}
In the ambient \((\infty,2)\)-category of infinity-categories, mate formation
identifies the relevant mapping spaces of natural transformations and is
compatible with pasting. Applying that calculus to the
unit and counit of
\(\Sigma_{e_{X/S}}\dashv e_{X/S}^*\dashv\Pi_{e_{X/S}}\) identifies
\(\eta^\Sigma\) with \(\varepsilon^\Pi\) and
\(e_{X/S}^*\varepsilon^\Sigma_{\Sigma_{e_{X/S}}}\) with
\(e_{X/S}^*\eta^\Pi_{\Pi_{e_{X/S}}}\) under the induced adjunctions. The triangle identities supply the required comparisons. This is the
adjunction calculus for infinity-categories; see
\cite{LurieHTT,LurieHigherAlgebra}.
\end{enumerate}
\end{proof}

\begin{corollary}[De Rham formal-disk--jet adjunction]
\label{thm:disk-jet-adjunction}
There is a canonical adjunction
\[
  \mathsf{Disk}^{\mathrm{dR},\infty}_{X/S}\dashv J^\infty_{X/S}.
\]
It is induced from the dependent-adjoint diagram
\[
\begin{tikzcd}[column sep=huge]
(\mathcal H_{\Sp})_{/Q_{X/S}}
  \ar[r,shift left=1.4ex,"e_{X/S}^*"]
&
(\mathcal H_{\Sp})_{/X}
  \ar[l,shift left=1.4ex,"\Sigma_{e_{X/S}}"]
  \ar[l,shift right=1.4ex,"\Pi_{e_{X/S}}"']
\end{tikzcd}
\]
by composing across the two adjunctions.
Under the relation correspondence it may also be written
\[
\begin{aligned}
\Map_X(\mathsf{Disk}^{\mathrm{dR},\infty} Z,E)
&\simeq \Map_{\mathcal R^{\mathrm{dR}}_{X/S}}(\mathsf t^*Z,\mathsf s^*E)\\
&\simeq \Map_{\mathcal R^{\mathrm{dR}}_{X/S}}(\mathsf s^*Z,\mathsf t^*E)\\
&\simeq \Map_X(Z,J^\infty E),
\end{aligned}
\]
where the middle equivalence is transport along inversion.
\end{corollary}

\begin{proposition}[Maps from de Rham formal disks]
\label{prop:maps-from-formal-disks}
For \(Z,E\in(\mathcal H_{\Sp})_{/X}\),
\[
  \Map_X(Z,J^\infty E)
  \simeq
  \Map_X(\mathsf{Disk}^{\mathrm{dR},\infty} Z,E).
\]
If \(Z=a:T\to X\) is a generalized point, then under the inversion presentation
\[
  \mathsf{Disk}^{\mathrm{dR},\infty}\simeq\Sigma_{\mathsf t}\mathsf s^*
\]
the object \(\mathsf{Disk}^{\mathrm{dR},\infty}(a)\) is canonically the de Rham formal disk
\(\operatorname{Nbh}^{\mathrm{dR}}_{T/S}(T\times_SX)\) regarded over \(X\) by evaluation. Hence
\[
  \Map_X(a,J^\infty E)
  \simeq
  \Map_X(\operatorname{Nbh}^{\mathrm{dR}}_{T/S}(T\times_SX),E).
\]
\end{proposition}

\begin{proof}
\leavevmode
\begin{enumerate}[label=\textup{(\arabic*)}]
    \item \textbf{Apply the disk--jet adjunction.}
The first equivalence is the adjunction \(\mathsf{Disk}^{\mathrm{dR},\infty}\dashv J^\infty\).

    \item \textbf{Use the inversion presentation.}
For the
second statement, the inversion form of
\Ref{prop:relation-jet-formulas} gives
\[
  \mathsf{Disk}^{\mathrm{dR},\infty}(a)
  \simeq
  \Sigma_{\mathsf t}\mathsf s^*(a).
\]

    \item \textbf{Identify the generalized de Rham formal disk.}
The pullback \(\mathsf s^*(a)\) is
\[
  T\times_{X,\mathsf s}\mathcal R^{\mathrm{dR}}_{X/S}
  \simeq
  \operatorname{Nbh}^{\mathrm{dR}}_{T/S}(T\times_SX)
\]
by \Ref{thm:formal-disk-graph-neighbourhood}, and postcomposition with \(\mathsf t\) is exactly
the evaluation map of that disk.
\end{enumerate}
\end{proof}

The reader may now compress the jet construction to two facts:
\[
  \mathsf{Disk}^{\mathrm{dR},\infty}_{X/S}
  \dashv
  J^\infty_{X/S},
  \qquad
  a^*J^\infty E
  \simeq
  \Pi_{\pi_a}\operatorname{ev}_a^*E.
\]
The remaining proposition records a closure property of the jet comonad needed
for later pullback constructions; after that, substitution completes the
section.

\begin{proposition}[Limit preservation]
\label{prop:jet-limits}
The functor \(J^\infty\) preserves all small limits. Consequently the
forgetful functor
\[
  U_J:\operatorname{Coalg}_{J^\infty}((\mathcal H_{\Sp})_{/X})\to(\mathcal H_{\Sp})_{/X}
\]
creates all small limits.
\end{proposition}

\begin{proof}
\leavevmode
\begin{enumerate}[label=\textup{(\arabic*)}]
    \item \textbf{Establish limit preservation.}
Pullback preserves all limits and \(\Pi_{e_{X/S}}\) is a right adjoint, hence
\(J^\infty=e_{X/S}^*\Pi_{e_{X/S}}\) preserves limits.

    \item \textbf{Create limits in the coalgebra category.}
The Eilenberg--Moore construction for a limit-preserving comonad
\cite{LurieHigherAlgebra} then gives the second assertion: form the limit of the underlying diagram and use
\(J^\infty(\lim E_i)\simeq\lim J^\infty E_i\) to assemble its contractibly unique coaction.
\end{enumerate}
\end{proof}

\subsection{Substitution of the jet calculus}
\label{sec:jet-base-change}

Because both jet endofunctors come from the same dependent-adjoint triple,
substitution is the Beck--Chevalley comparison for base change of \(e_{X/S}\). We establish that comparison now,
before equations and solution objects use the calculus parametrically.

\begin{theorem}[Base change]
\label{thm:jet-base-change}
Let \(g:S'\to S\), put
\[
  X':=X\times_SS',
\]
and write
\[
  g_X:X'\to X,
  \qquad
  g_Q:Q_{X'/S'}\to Q_{X/S}.
\]
There is a canonical Cartesian square
\[
\begin{tikzcd}
X' \ar[r,"g_X"] \ar[d,"e_{X'/S'}"']
  & X \ar[d,"e_{X/S}"]\\
Q_{X'/S'} \ar[r,"g_Q"']
  & Q_{X/S}.
\end{tikzcd}
\]
Moreover,
\[
  Q_{X'/S'}\simeq Q_{X/S}\times_SS'.
\]
Consequently there are canonical equivalences
\[
  g_X^*\mathsf{Disk}^{\mathrm{dR},\infty}_{X/S}
  \simeq
  \mathsf{Disk}^{\mathrm{dR},\infty}_{X'/S'}g_X^*,
  \qquad
  g_X^*J^\infty_{X/S}
  \simeq
  J^\infty_{X'/S'}g_X^*.
\]
They identify the monad and comonad structures and the adjunction
\(\mathsf{Disk}^{\mathrm{dR},\infty}\dashv J^\infty\). The canonical disk coactions and Eilenberg--Moore comparison introduced below
are transported by the same base-change equivalences.
\end{theorem}

\begin{proof}
\leavevmode
\begin{enumerate}[label=\textup{(\arabic*)}]
    \item \textbf{Base-change the effective quotient.}
Chapter~\ref{chap:spectral-de-Rham} gives the coherent base-change equivalence
\[
  Q_{X'/S'}\simeq Q_{X/S}\times_SS',
\]
and functoriality of the effective-image factorization identifies the displayed
square as Cartesian.

    \item \textbf{Apply left Beck--Chevalley.}
Left Beck--Chevalley gives
\[
\begin{aligned}
  g_X^*\mathsf{Disk}^{\mathrm{dR},\infty}_{X/S}
  &=
  g_X^*e_{X/S}^*\Sigma_{e_{X/S}}\\
  &\simeq
  e_{X'/S'}^*g_Q^*\Sigma_{e_{X/S}}\\
  &\simeq
  e_{X'/S'}^*\Sigma_{e_{X'/S'}}g_X^*
  =
  \mathsf{Disk}^{\mathrm{dR},\infty}_{X'/S'}g_X^*.
\end{aligned}
\]

    \item \textbf{Apply right Beck--Chevalley.}
Right Beck--Chevalley gives the analogous calculation for \(J^\infty\).

    \item \textbf{Transport the structural maps by mate coherence.}
These
comparisons are mate transformations for a coherently base-changed adjoint
triple, so they identify the units, counits, monad multiplication,
comonad comultiplication, and composite-adjunction unit.

    \item \textbf{Transport the later comparison structures.}
Naturality of the
later comparison coactions and mate equivalence then gives their base-change
compatibility as well.

    \item \textbf{Record higher Beck--Chevalley coherence.}
All higher coherence is the higher Beck--Chevalley
pasting coherence.
\end{enumerate}
\end{proof}

Arbitrary base change therefore transports the disk monad, the jet comonad, and
their adjunction together. This completes the section's interface: a
generalized point sees sections over its intrinsic de Rham disk, while a
change of base transports the entire construction by Beck--Chevalley. Later
base-change statements for equations, prolongations, and solutions will be
formal consequences of this result.

\section{Crystalline transport and intrinsic formally integrable PDEs}
\label{sec:jets-descent-pde}

The preceding section produced the disk--jet adjunction and its substitution
law. The remaining question is geometric: what does an action or coaction on an
object over \(X\) actually mean? This section answers that question once. Mate
formation identifies disk actions with jet coactions, and effective descent
identifies either one with an object over \(Q_{X/S}\). Crystalline transport is
then the degree-one relation-theoretic form of that same descent datum.

The section therefore proves one structural identification in progressively
more concrete presentations. Once the quotient, algebra, coalgebra, and
relation descriptions have been connected, later sections will use them
without re-establishing their equivalence.

\subsection{Homotopy-coherent mates and Eilenberg--Moore structures}

The disk--jet adjunction already transposes individual structure maps. The next
lemma is a bookkeeping result: it upgrades that transposition to the full
Eilenberg--Moore level, including structured morphisms. Its abstract notation
is used only for this passage and need not remain active afterwards.

\begin{lemma}[Adjoint monad--comonad duality]
\label{lem:homotopy-coherent-adjoint-monad-duality}
Let
\[
  T:\mathcal C\to\mathcal C
\]
be a homotopy-coherent monad admitting a right adjoint
\[
  T\dashv J.
\]
Transport the monad structure of \(T\) to \(J\) by taking right mates under the
iterated adjunctions
\[
  T^n\dashv J^n.
\]
Then \(J\) carries the resulting homotopy-coherent comonad structure, and mate
formation induces a canonical equivalence
\[
  \mathfrak M_{T\dashv J}:
  \operatorname{Alg}_T(\mathcal C)
  \xrightarrow{\simeq}
  \operatorname{Coalg}_J(\mathcal C)
\]
over \(\mathcal C\). Thus, if
\[
  U_T:\operatorname{Alg}_T(\mathcal C)\to\mathcal C,
  \qquad
  U_J:\operatorname{Coalg}_J(\mathcal C)\to\mathcal C
\]
are the forgetful functors, there is a canonical equivalence
\[
  U_J\mathfrak M_{T\dashv J}\simeq U_T.
\]

If
\[
  \eta^{T,J}:\id_{\mathcal C}\to JT,
  \qquad
  \varepsilon^{T,J}:TJ\to\id_{\mathcal C}
\]
are the unit and counit of \(T\dashv J\), then the equivalence sends an action
\[
  s:TE\to E
\]
to the coaction
\[
  s^\sharp
  :=
  J(s)\eta^{T,J}_E:
  E\to JE,
\]
and sends a coaction
\[
  \rho:E\to JE
\]
to the action
\[
  \rho^\flat
  :=
  \varepsilon^{T,J}_E\,T(\rho):
  TE\to E.
\]
These assignments are mutually inverse on structure spaces and preserve the
underlying objects and morphisms of \(\mathcal C\).
\end{lemma}

\begin{proof}
\leavevmode
\begin{enumerate}[label=\textup{(\arabic*)}]
    \item \textbf{Form the iterated adjunctions.}
In the ambient \((\infty,2)\)-category of infinity-categories, the adjunction
\(T\dashv J\) induces adjunctions
\[
  T^n\dashv J^n
\]
for every \(n\geq0\). Homotopy-coherent mate formation identifies the mapping
spaces of natural transformations among these iterates and is compatible with
whiskering and pasting; compare
\cite{LurieHTT,LurieHigherAlgebra}.

    \item \textbf{Construct the right-mate comonad.}
Taking right mates of the unit, multiplication, and higher coherence maps of
the monad \(T\) therefore produces the counit, comultiplication, and higher
coherence maps of a comonad \(J\).

    \item \textbf{Identify algebra and coalgebra structure spaces.}
The same mate equivalences identify, for
every object \(E\), the full space of \(T\)-algebra structures on \(E\) with the
full space of \(J\)-coalgebra structures on \(E\), and identify the corresponding
structured mapping spaces. Hence they assemble to an equivalence
\[
  \operatorname{Alg}_T(\mathcal C)
  \xrightarrow{\simeq}
  \operatorname{Coalg}_J(\mathcal C)
\]
over \(\mathcal C\).

    \item \textbf{Recover the objectwise formulas.}
The displayed formulas are the degree-one adjunct formulas for
\(T\dashv J\). Their composites are equivalent to the original structure maps
by the triangle identities, so they describe the objectwise effect of the
Eilenberg--Moore equivalence.
\end{enumerate}
\end{proof}

\begin{theorem}[Mate equivalence of \(\mathsf{Disk}^{\mathrm{dR},\infty}\)-algebras and \(J^\infty\)-coalgebras]
\label{thm:mate-em-equivalence}
Mate formation for
\[
  \mathsf{Disk}^{\mathrm{dR},\infty}\dashv J^\infty
\]
extends to a canonical equivalence
\[
  \mathfrak M:
  \operatorname{Alg}_{\mathsf{Disk}^{\mathrm{dR},\infty}}((\mathcal H_{\Sp})_{/X})
  \xrightarrow{\simeq}
  \operatorname{Coalg}_{J^\infty}((\mathcal H_{\Sp})_{/X}).
\]
If
\[
  U_{\mathsf{Disk}}:
  \operatorname{Alg}_{\mathsf{Disk}^{\mathrm{dR},\infty}}((\mathcal H_{\Sp})_{/X})
  \longrightarrow
  (\mathcal H_{\Sp})_{/X},
  \qquad
  U_J:
  \operatorname{Coalg}_{J^\infty}((\mathcal H_{\Sp})_{/X})
  \longrightarrow
  (\mathcal H_{\Sp})_{/X}
\]
are the two forgetful functors, then there is a canonical equivalence
\[
  U_J\mathfrak M
  \xrightarrow{\simeq}
  U_{\mathsf{Disk}}.
\]
Thus \(\mathfrak M\) is canonically an equivalence over
\((\mathcal H_{\Sp})_{/X}\).

Explicitly, if
\[
  s:\mathsf{Disk}^{\mathrm{dR},\infty} E\to E
\]
is a \(\mathsf{Disk}^{\mathrm{dR},\infty}\)-algebra action, its corresponding coaction is the adjunct
\[
  \rho_s
  :=
  s^\sharp
  =
  J^\infty(s)\eta^{\mathrm{adj}}_E:
  E\longrightarrow J^\infty E.
\]
Conversely, if
\[
  \rho:E\to J^\infty E
\]
is a \(J^\infty\)-coalgebra coaction, its corresponding action is the adjoint
\[
  s_\rho
  :=
  \rho^\flat
  =
  \varepsilon^{\mathrm{adj}}_E\,\mathsf{Disk}^{\mathrm{dR},\infty}(\rho):
  \mathsf{Disk}^{\mathrm{dR},\infty} E\longrightarrow E.
\]
These constructions are mutually inverse on the full
structure spaces.

Under this correspondence the algebra unit law
\[
  s\eta^\Sigma_E\simeq\id_E
\]
is the mate of the coalgebra counit law
\[
  \varepsilon^\Pi_E\rho_s\simeq\id_E,
\]
and the algebra associativity law
\[
  s\mu_E\simeq s\,\mathsf{Disk}^{\mathrm{dR},\infty}(s)
\]
is the mate of the coalgebra coassociativity law
\[
  \delta^J_E\rho_s
  \simeq
  J^\infty(\rho_s)\rho_s.
\]
Likewise, if a structured morphism has underlying map
\[
  f:E\to E',
\]
its mate has underlying map canonically identified with the same \(f\);
preservation of the algebra actions is equivalent to preservation of the
corresponding coalgebra coactions.
\end{theorem}

\begin{proof}
\leavevmode
\begin{enumerate}[label=\textup{(\arabic*)}]
    \item \textbf{Identify the right-mate comonad.}
By Proposition~\ref{prop:iterated-relation-correspondences}, the iterated
adjunctions
\[
  (\mathsf{Disk}^{\mathrm{dR},\infty})^n\dashv(J^\infty)^n
\]
identify the full monad structure of \(\mathsf{Disk}^{\mathrm{dR},\infty}\) with
the right-mate comonad structure of \(J^\infty\). In particular,
the right mates of
\[
  \eta^\Sigma:\id\to \mathsf{Disk}^{\mathrm{dR},\infty},
  \qquad
  \mu:(\mathsf{Disk}^{\mathrm{dR},\infty})^2\to \mathsf{Disk}^{\mathrm{dR},\infty}
\]
are precisely
\[
  \varepsilon^\Pi:J^\infty\to\id,
  \qquad
  \delta^J:J^\infty\to(J^\infty)^2,
\]
as identified by the simplicial relation calculus.

    \item \textbf{Apply adjoint--monad duality.}
Apply
Lemma~\ref{lem:homotopy-coherent-adjoint-monad-duality}
to the adjunction
\[
  \mathsf{Disk}^{\mathrm{dR},\infty}\dashv J^\infty.
\]
It produces a canonical equivalence
\[
  \mathfrak M:
  \operatorname{Alg}_{\mathsf{Disk}^{\mathrm{dR},\infty}}((\mathcal H_{\Sp})_{/X})
  \xrightarrow{\simeq}
  \operatorname{Coalg}_{J^\infty}((\mathcal H_{\Sp})_{/X})
\]
together with the canonical equivalence
\[
  U_J\mathfrak M
  \simeq
  U_{\mathsf{Disk}}.
\]
Hence the comparison is canonically defined over \((\mathcal H_{\Sp})_{/X}\).

    \item \textbf{Record the objectwise structure maps.}
The objectwise structure-map formulas supplied by the lemma are
\[
  s^\sharp
  =
  J^\infty(s)\eta^{\mathrm{adj}}_E,
  \qquad
  \rho^\flat
  =
  \varepsilon^{\mathrm{adj}}_E\mathsf{Disk}^{\mathrm{dR},\infty}(\rho).
\]

    \item \textbf{Verify the inverse structure maps.}
They are mutually inverse by the triangle identities for
\[
  \mathsf{Disk}^{\mathrm{dR},\infty}\dashv J^\infty,
\]
and the lemma identifies the full algebra and
coalgebra structure spaces, with these degree-one maps as their objectwise
components.

    \item \textbf{Apply mate correspondence to the structural laws.}
At the first two structural levels, the mate correspondence gives
\[
  s\eta^\Sigma_E\simeq\id_E
  \quad\longleftrightarrow\quad
  \varepsilon^\Pi_E\rho_s\simeq\id_E
\]
and
\[
  s\mu_E\simeq s\,\mathsf{Disk}^{\mathrm{dR},\infty}(s)
  \quad\longleftrightarrow\quad
  \delta^J_E\rho_s
  \simeq
  J^\infty(\rho_s)\rho_s.
\]

    \item \textbf{Compare structured morphisms.}
For an underlying morphism
\[
  f:E\to E',
\]
the same mate correspondence carries the algebra-morphism
compatibility
\[
  f\,s\simeq s'\,\mathsf{Disk}^{\mathrm{dR},\infty}(f)
\]
to the coalgebra-morphism compatibility
\[
  J^\infty(f)\rho_s\simeq\rho_{s'}f,
\]
and preserves the corresponding structured morphism data.

    \item \textbf{Conclude the Eilenberg--Moore equivalence.}
This proves the Eilenberg--Moore equivalence over
\((\mathcal H_{\Sp})_{/X}\).
\end{enumerate}
\end{proof}

\begin{corollary}[Comparison actions and coactions are mates]
\label{cor:comparison-actions-coactions-mates}
Let \(F\to Q_{X/S}\), put \(E=e_{X/S}^*F\), and define
\[
  s_F:=e_{X/S}^*\varepsilon^\Sigma_F:\mathsf{Disk}^{\mathrm{dR},\infty} E\to E,
  \qquad
  \rho_F:=e_{X/S}^*\eta^\Pi_F:E\to J^\infty E.
\]
Then
\[
  \rho_F\simeq s_F^\sharp.
\]
Consequently
\[
  \mathfrak M(E,s_F)\simeq(E,\rho_F)
\]
canonically over \(E\).
\end{corollary}

\begin{proof}
\leavevmode
\begin{enumerate}[label=\textup{(\arabic*)}]
    \item \textbf{Expand the adjunct.}
Using the formula for the unit \(\eta^{\mathrm{adj}}\) of
\(\mathsf{Disk}^{\mathrm{dR},\infty}\dashv J^\infty\),
\[
\begin{aligned}
  s_F^\sharp
  &=
  J^\infty(s_F)\eta^{\mathrm{adj}}_E\\
  &=
  e_{X/S}^*\Pi_{e_{X/S}}(e_{X/S}^*\varepsilon^\Sigma_F)\,
  e_{X/S}^*\eta^\Pi_{\Sigma_{e_{X/S}} E}\,
  \eta^\Sigma_E.
\end{aligned}
\]

    \item \textbf{Use naturality of the unit.}
Naturality of
\(\eta^\Pi:\id_{(\mathcal H_{\Sp})_{/Q_{X/S}}}\to\Pi_{e_{X/S}}e_{X/S}^*\)
with respect to
\(\varepsilon^\Sigma_F:\Sigma_{e_{X/S}}e_{X/S}^*F\to F\) gives
\[
  \Pi_{e_{X/S}}e_{X/S}^*(\varepsilon^\Sigma_F)\,\eta^\Pi_{\Sigma_{e_{X/S}}e_{X/S}^*F}
  \simeq
  \eta^\Pi_F\,\varepsilon^\Sigma_F.
\]

    \item \textbf{Apply the triangle identity.}
Applying \(e_{X/S}^*\) and composing with \(\eta^\Sigma_E\) therefore yields
\[
\begin{aligned}
  s_F^\sharp
  &\simeq
  e_{X/S}^*\eta^\Pi_F\,
  e_{X/S}^*\varepsilon^\Sigma_F\,
  \eta^\Sigma_E\\
  &\simeq
  e_{X/S}^*\eta^\Pi_F
  =
  \rho_F,
\end{aligned}
\]
where the second equivalence is the triangle identity for
\(\Sigma_{e_{X/S}}\dashv e_{X/S}^*\).
\end{enumerate}
\end{proof}

\subsubsection{The canonical jet coaction on a disk object}

The Eilenberg--Moore equivalence supplies a distinguished jet coaction on every
disk object. The following construction names it ``free'' in anticipation of
the universal property proved later; until that theorem is reached, only the
explicitly constructed comparison coaction is used.

\begin{construction}[Free jet coalgebra on the disk object]
\label{con:free-disk-coalgebra}
For \(Z\in(\mathcal H_{\Sp})_{/X}\), define
\[
  \rho_Z^{\mathbb F}
  :=
  e_{X/S}^*\eta^\Pi_{\Sigma_{e_{X/S}} Z}:
  \mathsf{Disk}^{\mathrm{dR},\infty} Z\longrightarrow J^\infty \mathsf{Disk}^{\mathrm{dR},\infty} Z,
\]
and put
\[
  \mathbb F(Z)
  :=
  \bigl(\mathsf{Disk}^{\mathrm{dR},\infty} Z,\rho_Z^{\mathbb F}\bigr).
\]
This is the canonical comparison \(J^\infty\)-coalgebra on the object
\(e_{X/S}^*\Sigma_{e_{X/S}} Z\) supplied by the adjunction
\(e_{X/S}^*\dashv\Pi_{e_{X/S}}\). Equivalently, by
Theorem~\ref{thm:mate-em-equivalence}, it is the coalgebra corresponding to the
free \(\mathsf{Disk}^{\mathrm{dR},\infty}\)-algebra action
\[
  \mu_Z:\mathsf{Disk}^{\mathrm{dR},\infty} \mathsf{Disk}^{\mathrm{dR},\infty} Z\to \mathsf{Disk}^{\mathrm{dR},\infty} Z.
\]
Its coalgebra laws are supplied by the adjunction and mate calculus. Its free
universal property is proved later in
Theorem~\ref{thm:free-forgetful-cofree}.
\end{construction}

\begin{proposition}[Unit factorization for the de Rham formal-disk--jet adjunction]
\label{prop:disk-jet-unit-factorization}
The unit \(\eta^{\mathrm{adj}}_Z:Z\to J^\infty \mathsf{Disk}^{\mathrm{dR},\infty} Z\) satisfies
\[
  \eta^{\mathrm{adj}}_Z\simeq\rho_Z^{\mathbb F}\eta^\Sigma_Z.
\]
\end{proposition}

\begin{proof}
\leavevmode
\begin{enumerate}[label=\textup{(\arabic*)}]
    \item \textbf{Invoke the composite-adjunction unit formula.}
The assertion is the unit formula in Theorem~\ref{thm:adjoint-triple-mates}.
\end{enumerate}
\end{proof}

\begin{lemma}[Mate calculus]
\label{lem:jet-mates}
Let \(s:\mathsf{Disk}^{\mathrm{dR},\infty} Z\to E\), and let
\[
  s^\sharp:Z\to J^\infty E
\]
be its adjunct. Then
\[
  \varepsilon^\Pi_Es^\sharp\simeq s\eta^\Sigma_Z.
\]
For every \(f:E\to F\), the adjunct of \(fs\) is
\[
  J^\infty(f)s^\sharp.
\]
Both identities are natural in the corresponding mapping spaces.
\end{lemma}

\begin{proof}
\leavevmode
\begin{enumerate}[label=\textup{(\arabic*)}]
    \item \textbf{Write the adjunct through the adjunction unit.}
Under the adjunction
\[
  \mathsf{Disk}^{\mathrm{dR},\infty}\dashv J^\infty,
\]
the adjunct of
\[
  s:\mathsf{Disk}^{\mathrm{dR},\infty} Z\to E
\]
is canonically
\[
  s^\sharp
  \simeq
  J^\infty(s)\eta^{\mathrm{adj}}_Z,
\]
where
\[
  \eta^{\mathrm{adj}}_Z:Z\to J^\infty \mathsf{Disk}^{\mathrm{dR},\infty} Z
\]
is the unit of the formal-disk--jet adjunction.

    \item \textbf{Apply naturality of the comonad counit.}
Hence naturality of the comonad
counit
\[
  \varepsilon^\Pi:J^\infty\to\id_{(\mathcal H_{\Sp})_{/X}}
\]
gives
\[
\begin{aligned}
  \varepsilon^\Pi_Es^\sharp
  &\simeq
  \varepsilon^\Pi_EJ^\infty(s)\eta^{\mathrm{adj}}_Z\\
  &\simeq
  s\,\varepsilon^\Pi_{\mathsf{Disk}^{\mathrm{dR},\infty} Z}\eta^{\mathrm{adj}}_Z.
\end{aligned}
\]

    \item \textbf{Insert the disk--jet unit factorization.}
By
Proposition~\ref{prop:disk-jet-unit-factorization},
\[
  \eta^{\mathrm{adj}}_Z\simeq\rho_Z^{\mathbb F}\eta^\Sigma_Z.
\]

    \item \textbf{Use the comparison-coalgebra counit.}
The map
\[
  \rho_Z^{\mathbb F}
  =
  e_{X/S}^*\eta^\Pi_{\Sigma_{e_{X/S}} Z}:
  \mathsf{Disk}^{\mathrm{dR},\infty} Z\to J^\infty \mathsf{Disk}^{\mathrm{dR},\infty} Z
\]
is the canonical comparison coalgebra structure supplied by
\(e_{X/S}^*\dashv\Pi_{e_{X/S}}\), so its counit identity is
\[
  \varepsilon^\Pi_{\mathsf{Disk}^{\mathrm{dR},\infty} Z}\rho_Z^{\mathbb F}
  \simeq
  \id_{\mathsf{Disk}^{\mathrm{dR},\infty} Z}.
\]

    \item \textbf{Conclude the first mate identity.}
Therefore
\[
\begin{aligned}
  \varepsilon^\Pi_Es^\sharp
  &\simeq
  s\,\varepsilon^\Pi_{\mathsf{Disk}^{\mathrm{dR},\infty} Z}\rho_Z^{\mathbb F}\eta^\Sigma_Z\\
  &\simeq
  s\eta^\Sigma_Z,
\end{aligned}
\]
which proves the first identity.

    \item \textbf{Fix the target morphism.}
Now let
\[
  f:E\to F.
\]

    \item \textbf{Apply target naturality.}
Naturality of the mapping-space equivalence
\[
  \Map_X(\mathsf{Disk}^{\mathrm{dR},\infty} Z,-)
  \simeq
  \Map_X(Z,J^\infty-)
\]
in the target variable identifies the adjunct of
\[
  fs:\mathsf{Disk}^{\mathrm{dR},\infty} Z\to F
\]
with
\[
  J^\infty(f)s^\sharp:
  Z\to J^\infty F.
\]

    \item \textbf{Retain the higher homotopies.}
Since these identities are obtained inside the natural equivalence of mapping
spaces defining the adjunction, they are automatically compatible with all
higher homotopies in those mapping spaces.
\end{enumerate}
\end{proof}

\subsection{Crystalline transport and effective de Rham descent}
\label{sec:crystalline-descent}

The previous subsection established the formal equivalence between disk actions
and jet coactions. We now supply its geometric meaning. Because
\(\mathcal R^{\mathrm{dR}}_{X/S,\bullet}\) is already the \v{C}ech nerve of the
effective quotient, effective descent identifies the Eilenberg--Moore
categories with the quotient slice. Transposing a coaction then exposes the
degree-one transport, while the coalgebra laws become exactly the unit and
cocycle laws of the relation groupoid.

\begin{proposition}[Effective \v{C}ech descent for the quotient]
\label{prop:effective-cech-descent-Q}
Pullback along the effective epimorphism
\[
  e_{X/S}:X\twoheadrightarrow Q_{X/S}
\]
induces a canonical equivalence
\[
  (\mathcal H_{\Sp})_{/Q_{X/S}}
  \xrightarrow{\simeq}
  \operatorname*{Tot}_{[n]\in\Delta}
  (\mathcal H_{\Sp})_{/\mathcal R^{\mathrm{dR}}_{X/S,n}}.
\]
Under this equivalence an object over \(Q_{X/S}\) determines its degree-zero
object over \(X\), its degree-one crystalline transport, the \v{C}ech cocycle, and all
higher simplicial coherence simultaneously.
\end{proposition}

\begin{proof}
\leavevmode
\begin{enumerate}[label=\textup{(\arabic*)}]
    \item \textbf{Identify the effective Čech nerve.}
Theorem~\ref{thm:relation-effective-presentation} of
Chapter~\ref{chap:spectral-de-Rham} identifies
\(\mathcal R^{\mathrm{dR}}_{X/S,\bullet}\) with the \v{C}ech nerve of \(e_{X/S}\).

    \item \textbf{Apply slice descent.}
Apply
Theorem~\ref{thm:parameterized-slice-cech-descent}.
\end{enumerate}
\end{proof}

\begin{proposition}[Conservativity of effective pullback]
\label{prop:effective-pullback-conservative}
The pullback functor
\[
  e_{X/S}^*:(\mathcal H_{\Sp})_{/Q_{X/S}}\longrightarrow(\mathcal H_{\Sp})_{/X}
\]
is conservative.
\end{proposition}

\begin{proof}
\leavevmode
\begin{enumerate}[label=\textup{(\arabic*)}]
    \item \textbf{Invoke effective descent for epimorphisms.}
Effective epimorphisms are effective descent morphisms in an infinity-topos
\cite{LurieHTT}.

    \item \textbf{Pull the morphism back to the Čech nerve.}
If \(f:F\to G\) becomes an equivalence
after pullback to \(X\), then every degree of its pullback to the \v{C}ech
nerve \(\mathcal R^{\mathrm{dR}}_{X/S,\bullet}\) is an equivalence.

    \item \textbf{Reconstruct the original morphism.}
Under
Proposition~\ref{prop:effective-cech-descent-Q}, the original morphism is the
totalization of those degreewise morphisms, hence is an equivalence.
\end{enumerate}
\end{proof}

\begin{theorem}[Crystalline descent theorem]
\label{thm:effective-jet-descent}
The comparison functors
\[
  K_\Sigma:
  (\mathcal H_{\Sp})_{/Q_{X/S}}
  \longrightarrow
  \operatorname{Alg}_{\mathsf{Disk}^{\mathrm{dR},\infty}}
  ((\mathcal H_{\Sp})_{/X}),
\]
\[
  F\longmapsto
  \bigl(
    e_{X/S}^*F,
    e_{X/S}^*\varepsilon^\Sigma_F
  \bigr),
\]
and
\[
  K_\Pi:
  (\mathcal H_{\Sp})_{/Q_{X/S}}
  \longrightarrow
  \operatorname{Coalg}_{J^\infty}
  ((\mathcal H_{\Sp})_{/X}),
\]
\[
  F\longmapsto
  \bigl(
    e_{X/S}^*F,
    e_{X/S}^*\eta^\Pi_F
  \bigr),
\]
are equivalences of infinity-categories.

Moreover, if
\[
  \mathfrak M:
  \operatorname{Alg}_{\mathsf{Disk}^{\mathrm{dR},\infty}}
  ((\mathcal H_{\Sp})_{/X})
  \xrightarrow{\simeq}
  \operatorname{Coalg}_{J^\infty}
  ((\mathcal H_{\Sp})_{/X})
\]
is the mate equivalence of
Theorem~\ref{thm:mate-em-equivalence}, then
\[
  \mathfrak M K_\Sigma
  \simeq
  K_\Pi
\]
canonically over \((\mathcal H_{\Sp})_{/X}\). Hence there are canonical
equivalences
\[
  (\mathcal H_{\Sp})_{/Q_{X/S}}
  \simeq
  \operatorname{Alg}_{\mathsf{Disk}^{\mathrm{dR},\infty}}
  ((\mathcal H_{\Sp})_{/X})
  \simeq
  \operatorname{Coalg}_{J^\infty}
  ((\mathcal H_{\Sp})_{/X}).
\]
\end{theorem}

\begin{proof}
\leavevmode
\begin{enumerate}[label=\textup{(\arabic*)}]
    \item \textbf{Establish presentability.}
The categories
\[
  (\mathcal H_{\Sp})_{/X},
  \qquad
  (\mathcal H_{\Sp})_{/Q_{X/S}}
\]
are presentable.

    \item \textbf{Establish conservativity.}
Pullback
\[
  e_{X/S}^*:
  (\mathcal H_{\Sp})_{/Q_{X/S}}
  \longrightarrow
  (\mathcal H_{\Sp})_{/X}
\]
is conservative by
Proposition~\ref{prop:effective-pullback-conservative}.

    \item \textbf{Verify realization preservation.}
It is also a left
adjoint, since
\[
  e_{X/S}^*\dashv\Pi_{e_{X/S}},
\]
and therefore preserves all geometric realizations required by the monadic
Barr--Beck theorem.

    \item \textbf{Identify the induced monad.}
For the adjunction
\[
  \Sigma_{e_{X/S}}\dashv e_{X/S}^*,
\]
the induced monad is
\[
  e_{X/S}^*\Sigma_{e_{X/S}}
  =
  \mathsf{Disk}^{\mathrm{dR},\infty}.
\]

    \item \textbf{Apply monadic Barr--Beck.}
Barr--Beck therefore gives
\[
  K_\Sigma:
  (\mathcal H_{\Sp})_{/Q_{X/S}}
  \xrightarrow{\simeq}
  \operatorname{Alg}_{\mathsf{Disk}^{\mathrm{dR},\infty}}
  ((\mathcal H_{\Sp})_{/X});
\]
see \cite[Theorem~4.7.3.5]{LurieHigherAlgebra}.

    \item \textbf{Fix the pulled-back object.}
For \(F\to Q_{X/S}\), put
\[
  E=e_{X/S}^*F.
\]

    \item \textbf{Identify the comparison structures.}
The disk-algebra comparison structure is
\[
  s_F
  =
  e_{X/S}^*\varepsilon^\Sigma_F:
  \mathsf{Disk}^{\mathrm{dR},\infty}E\to E,
\]
whereas the jet-coalgebra comparison structure is
\[
  \rho_F
  =
  e_{X/S}^*\eta^\Pi_F:
  E\to J^\infty E.
\]

    \item \textbf{Apply the mate comparison.}
Corollary~\ref{cor:comparison-actions-coactions-mates} identifies these maps as
mates:
\[
  \rho_F\simeq s_F^\sharp.
\]

    \item \textbf{Identify the comparison functors.}
Consequently
\[
  \mathfrak M K_\Sigma\simeq K_\Pi
\]
naturally over \((\mathcal H_{\Sp})_{/X}\).

    \item \textbf{Conclude the coalgebraic presentation.}
Since \(K_\Sigma\) and \(\mathfrak M\) are equivalences, so is \(K_\Pi\).
Thus the quotient slice, disk-algebra category, and jet-coalgebra category are
three presentations of the same effective de Rham descent category.
\end{enumerate}
\end{proof}

At this point the reader may replace the preceding constructions by one
interface:
\[
  \boxed{
  \text{object over }Q_{X/S}
  \;\Longleftrightarrow\;
  \text{formal-disk action}
  \;\Longleftrightarrow\;
  \text{jet coaction}
  \;\Longleftrightarrow\;
  \text{effective de Rham descent}.
  }
\]
The results that follow unpack the degree-one geometric content of this
equivalence as crystalline transport along \(\mathcal R^{\mathrm{dR}}_{X/S}\),
so that unit, cocycle, inverse, and higher compatibility can be read directly
from the same \v{C}ech groupoid.
\begin{theorem}[Transport, algebra actions, and coalgebra coactions]
\label{thm:transport-equivalence}
For \(E\to X\), define the spaces
\[
\begin{aligned}
  \operatorname{Alg}_{\mathsf{Disk}^{\mathrm{dR},\infty}}(E)
  &:=
  \operatorname{hofib}_{E}
  \left(
    \operatorname{Alg}_{\mathsf{Disk}^{\mathrm{dR},\infty}}((\mathcal H_{\Sp})_{/X})^\simeq
    \longrightarrow
    (\mathcal H_{\Sp})_{/X}^\simeq
  \right),\\
  \operatorname{Coalg}_{J^\infty}(E)
  &:=
  \operatorname{hofib}_{E}
  \left(
    \operatorname{Coalg}_{J^\infty}((\mathcal H_{\Sp})_{/X})^\simeq
    \longrightarrow
    (\mathcal H_{\Sp})_{/X}^\simeq
  \right),
\end{aligned}
\]
and let
\[
  \operatorname{Desc}_{e_{X/S}}(E)
\]
be the homotopy fibre over \(E\) of degree-zero evaluation from the maximal
subgroupoid of
\[
  \operatorname*{Tot}_{[n]\in\Delta}
  (\mathcal H_{\Sp})_{/\mathcal R^{\mathrm{dR}}_{X/S,n}}.
\]
Then there are canonical equivalences
\[
  \operatorname{Alg}_{\mathsf{Disk}^{\mathrm{dR},\infty}}(E)
  \simeq
  \operatorname{Coalg}_{J^\infty}(E)
  \simeq
  \operatorname{Desc}_{e_{X/S}}(E).
\]
Thus a \(\mathsf{Disk}^{\mathrm{dR},\infty}\)-algebra structure, a jet coalgebra structure, and coherent
de Rham-relation descent on the fixed underlying object are three
presentations of the same structure space.
\end{theorem}

\begin{proof}
\leavevmode
\begin{enumerate}[label=\textup{(\arabic*)}]
    \item \textbf{Take fibres of the mate equivalence.}
The first equivalence is obtained by taking homotopy fibres of the mate
equivalence over \((\mathcal H_{\Sp})_{/X}\) in
Theorem~\ref{thm:mate-em-equivalence}.

    \item \textbf{Take fibres of the descent equivalence.}
The second is obtained by combining
Theorem~\ref{thm:effective-jet-descent} with
Proposition~\ref{prop:effective-cech-descent-Q} and again taking the homotopy
fibre of the degree-zero underlying-object functor.

    \item \textbf{Identify the fixed underlying object.}
The homotopy fibres record the equivalence with the fixed object \(E\) at
the level of maximal subgroupoids.
\end{enumerate}
\end{proof}

\begin{construction}[Canonical crystalline transport]
\label{con:canonical-descent-action-coaction}
Let \(F\to Q_{X/S}\), put \(E:=e_{X/S}^*F\), and use the comparison action
\(s_F\) and coaction \(\rho_F\) of
Corollary~\ref{cor:comparison-actions-coactions-mates}. On
\[
  \mathcal R^{\mathrm{dR}}_{X/S}=X\times_{Q_{X/S}}X
\]
the two composites \(e_{X/S}\mathsf s\) and \(e_{X/S}\mathsf t\) agree. Hence the two
pullbacks of \(F\) are canonically identified, giving
\[
  \theta_F:
  \mathsf s^*E
  =
  (e_{X/S}\mathsf s)^*F
  \xrightarrow{\simeq}
  (e_{X/S}\mathsf t)^*F
  =
  \mathsf t^*E.
\]
This is the crystalline transport of \(F\), displayed by
\[
\begin{tikzcd}[column sep=huge,row sep=large]
\mathsf s^*E \ar[r,"\theta_F","\sim"']
  \ar[d]
&
\mathsf t^*E \ar[d]\\
\mathcal R^{\mathrm{dR}}_{X/S}
  \ar[r,equal]
&
\mathcal R^{\mathrm{dR}}_{X/S}.
\end{tikzcd}
\]
Under right Beck--Chevalley, the transpose of the comparison coaction \(\rho_F\) is
precisely \(\theta_F\); thus the descent presentation and the direct
coalgebraic construction of
Proposition~\ref{prop:jet-coaction-relation-transport} agree.
\end{construction}

\begin{proposition}[Crystalline transport extracted from a jet coaction]
\label{prop:jet-coaction-relation-transport}
Let
\[
  \rho:E\longrightarrow J^\infty E
\]
be a \(J^\infty\)-coalgebra structure. Under
\[
  J^\infty\simeq\Pi_{\mathsf s}\mathsf t^*,
\]
its right-adjoint transpose is a canonical morphism
\[
  \theta_\rho:\mathsf s^*E\longrightarrow \mathsf t^*E.
\]
The coalgebra laws canonically make this morphism an equivalence. Writing
\[
  \mathcal R^{\mathrm{dR}}_{X/S,2}=X\times_{Q_{X/S}}X\times_{Q_{X/S}}X
\]
and \(p_{ij}:\mathcal R^{\mathrm{dR}}_{X/S,2}\to\mathcal R^{\mathrm{dR}}_{X/S}\) for the edge projections, one has
\[
  \mathsf u^*\theta_\rho\simeq\id_E,
\]
\[
  p_{12}^*\theta_\rho\circ p_{01}^*\theta_\rho
  \simeq
  p_{02}^*\theta_\rho,
\]
and
\[
  \theta_\rho^{-1}\simeq\mathsf{inv}^*\theta_\rho.
\]
Thus the unit, cocycle, inverse, and remaining relation compatibility are
derived from the coalgebra laws and the \v{C}ech-groupoid structure.
\end{proposition}

\begin{proof}
\leavevmode
\begin{enumerate}[label=\textup{(\arabic*)}]
    \item \textbf{Transpose the jet coaction.}
Under the right Beck--Chevalley equivalence
\[
  J^\infty E
  \simeq
  \Pi_{\mathsf s}\mathsf t^*E,
\]
the map \(\rho:E\to J^\infty E\) is adjoint to a uniquely determined map
\[
  \theta_\rho:\mathsf s^*E\to \mathsf t^*E.
\]

    \item \textbf{Recover the unit law.}
The comonad counit
\[
  \varepsilon^\Pi_E:J^\infty E\to E
\]
is induced, under the relation correspondence, by restriction along the
relation unit \(\mathsf u\). Transposing the coalgebra counit law
\[
  \varepsilon^\Pi_E\rho\simeq\id_E
\]
therefore gives
\[
  \mathsf u^*\theta_\rho\simeq\id_E.
\]

    \item \textbf{Identify relation composition with monad multiplication.}
By the formal-disk monad proposition, the monad multiplication
\[
  \mu:(\mathsf{Disk}^{\mathrm{dR},\infty})^2\to \mathsf{Disk}^{\mathrm{dR},\infty}
\]
is induced by composition in the de Rham relation groupoid.

    \item \textbf{Identify the comonad mate.}
By
Theorem~\ref{thm:adjoint-triple-mates}, its mate is the comultiplication
\[
  \delta^J:J^\infty\to J^\infty J^\infty.
\]

    \item \textbf{Transpose coassociativity to the relation.}
Consequently, after applying right Beck--Chevalley twice, the two sides of the
coalgebra coassociativity law
\[
  \delta^J_E\rho
  \simeq
  J^\infty(\rho)\rho
\]
transpose respectively to
\[
  p_{02}^*\theta_\rho
  \qquad\text{and}\qquad
  p_{12}^*\theta_\rho\circ p_{01}^*\theta_\rho.
\]

    \item \textbf{Conclude the cocycle identity.}
Hence
\[
  p_{12}^*\theta_\rho\circ p_{01}^*\theta_\rho
  \simeq
  p_{02}^*\theta_\rho.
\]

    \item \textbf{Define the first return simplex.}
To derive invertibility, define
\[
  a:\mathcal R^{\mathrm{dR}}_{X/S}\to \mathcal R^{\mathrm{dR}}_{X/S,2},
  \qquad
  a(x_0,x_1):=(x_0,x_1,x_0).
\]
Then
\[
  p_{01}a=\id_{\mathcal R^{\mathrm{dR}}_{X/S}},
  \qquad
  p_{12}a=\mathsf{inv},
  \qquad
  p_{02}a=\mathsf u\,\mathsf s.
\]

    \item \textbf{Derive a left inverse.}
Pulling back the cocycle identity along \(a\) gives
\[
  \mathsf{inv}^*\theta_\rho\circ\theta_\rho
  \simeq
  \mathsf s^*\mathsf u^*\theta_\rho
  \simeq
  \id_{\mathsf s^*E}.
\]

    \item \textbf{Define the second return simplex.}
Likewise, for
\[
  b:\mathcal R^{\mathrm{dR}}_{X/S}\to \mathcal R^{\mathrm{dR}}_{X/S,2},
  \qquad
  b(x_0,x_1):=(x_1,x_0,x_1),
\]
one has
\[
  p_{01}b=\mathsf{inv},
  \qquad
  p_{12}b=\id_{\mathcal R^{\mathrm{dR}}_{X/S}},
  \qquad
  p_{02}b=\mathsf u\,\mathsf t,
\]

    \item \textbf{Derive a right inverse.}
and hence
\[
  \theta_\rho\circ\mathsf{inv}^*\theta_\rho
  \simeq
  \mathsf t^*\mathsf u^*\theta_\rho
  \simeq
  \id_{\mathsf t^*E}.
\]

    \item \textbf{Identify inverse transport.}
Thus \(\theta_\rho\) is an equivalence with
\[
  \theta_\rho^{-1}\simeq\mathsf{inv}^*\theta_\rho.
\]

    \item \textbf{Recover higher simplicial coherence.}
For the higher coherences, use
Proposition~\ref{prop:iterated-relation-correspondences}:
\[
  (J^\infty)^n
  \simeq
  \Pi_{v_n}\mathsf t_n^*
\]
identifies every iterated coalgebra structure map with transport over the
\(n\)-simplex \(\mathcal R^{\mathrm{dR}}_{X/S,n}\). The coalgebra laws transpose to the face and degeneracy identities of the
\v{C}ech nerve. The same simplicial object therefore supplies the remaining
transport compatibility.
\end{enumerate}
\end{proof}

\begin{proposition}[Crystalline transport extracted from a \(\mathsf{Disk}^{\mathrm{dR},\infty}\)-algebra]
\label{prop:disk-algebra-relation-transport}
Let
\[
  s:\mathsf{Disk}^{\mathrm{dR},\infty} E\longrightarrow E
\]
be a \(\mathsf{Disk}^{\mathrm{dR},\infty}\)-algebra structure. Under
\[
  \mathsf{Disk}^{\mathrm{dR},\infty}\simeq\Sigma_{\mathsf s}\mathsf t^*,
\]
its left-adjoint transpose is a canonical equivalence
\[
  \tau_s:\mathsf t^*E\xrightarrow{\simeq}\mathsf s^*E.
\]
Let
\[
  \rho:=\mathfrak M(E,s):E\longrightarrow J^\infty E
\]
be the corresponding coalgebra structure of
Theorem~\ref{thm:mate-em-equivalence}. Then
\[
  \tau_s
  \simeq
  \theta_\rho^{-1}
  \simeq
  \mathsf{inv}^*\theta_\rho.
\]
The algebra unit and associativity coherences are the reverse-oriented mates of
the relation-unit and cocycle coherences.
\end{proposition}

\begin{proof}
\leavevmode
\begin{enumerate}[label=\textup{(\arabic*)}]
    \item \textbf{Identify the correspondence adjunction through inversion.}
The adjunction
\[
  \mathsf{Disk}^{\mathrm{dR},\infty}\simeq\Sigma_{\mathsf s}\mathsf t^*
  \dashv
  J^\infty\simeq\Pi_{\mathsf s}\mathsf t^*
\]
is obtained from the correspondence adjunction after transporting the
right-adjoint presentation through relation inversion.

    \item \textbf{Transpose the disk action.}
Hence the mate of
\[
  s:\Sigma_{\mathsf s}\mathsf t^*E\to E
\]
is obtained by transposing \(s\) to
\[
  \tau_s:\mathsf t^*E\to \mathsf s^*E
\]
and then transporting along
\(\mathsf{inv}:\mathcal R^{\mathrm{dR}}_{X/S}\xrightarrow{\simeq}\mathcal R^{\mathrm{dR}}_{X/S}\).

    \item \textbf{Transpose the coalgebra coaction.}
The corresponding coalgebra coaction
\(\rho=s^\sharp\) transposes to
\[
  \theta_\rho:\mathsf s^*E\to \mathsf t^*E.
\]

    \item \textbf{Identify inversion with inverse transport.}
The two transposes are therefore exchanged by inversion. By
Proposition~\ref{prop:jet-coaction-relation-transport},
\[
  \theta_\rho^{-1}\simeq\mathsf{inv}^*\theta_\rho,
\]
so
\[
  \tau_s\simeq\theta_\rho^{-1}\simeq\mathsf{inv}^*\theta_\rho.
\]

    \item \textbf{Transport the higher coherences.}
Compatibility of mates with pasting identifies the algebra unit and
associativity coherences with the oppositely oriented relation unit and
cocycle coherences, including all higher homotopies.
\end{enumerate}
\end{proof}

\begin{proposition}[Degree-one transport and cocycle coherence]
\label{prop:degree-one-transport-cocycle}
For the canonical structures of
Construction~\ref{con:canonical-descent-action-coaction},
\[
  \mathsf u^*\theta_F\simeq\id_E,
\]
\[
  p_{12}^*\theta_F\circ p_{01}^*\theta_F
  \simeq
  p_{02}^*\theta_F,
\]
and
\[
  \theta_F^{-1}\simeq\mathsf{inv}^*\theta_F.
\]
The remaining simplicial compatibility is inherited from the single object
\(F\) on \(Q_{X/S}\).
\end{proposition}

\begin{proof}
\leavevmode
\begin{enumerate}[label=\textup{(\arabic*)}]
    \item \textbf{Apply the relation-transport theorem.}
Apply Proposition~\ref{prop:jet-coaction-relation-transport} to
\(\rho_F\); equivalently, use the degree-one and degree-two part of
Proposition~\ref{prop:effective-cech-descent-Q}.
\end{enumerate}
\end{proof}

\begin{proposition}[Explicit mate relation]
\label{prop:mate-relation-explicit}
The canonical formal-disk action and jet coaction satisfy
\[
  \varepsilon^\Pi_E\rho_F\simeq\id_E,
  \qquad
  \rho_Fs_F\simeq J^\infty(s_F)\rho_E^{\mathbb F}.
\]
Thus the formal-disk action and jet coaction are two adjoint presentations of the
same effective descent structure.
\end{proposition}

\begin{proof}
\leavevmode
\begin{enumerate}[label=\textup{(\arabic*)}]
    \item \textbf{Apply the counit identity.}
The first equation is the triangle identity for
\(e_{X/S}^*\dashv\Pi_{e_{X/S}}\).

    \item \textbf{Use the coalgebra-morphism equation.}
For the second, the action
\[
  s_F:(\mathsf{Disk}^{\mathrm{dR},\infty} E,\mu_E)\longrightarrow(E,s_F)
\]
is the canonical morphism from the free \(\mathsf{Disk}^{\mathrm{dR},\infty}\)-algebra to the comparison
algebra. Under the equivalence over \((\mathcal H_{\Sp})_{/X}\) of
Theorem~\ref{thm:effective-jet-descent}, the same underlying map is a morphism
from the free jet coalgebra
\[
  (\mathsf{Disk}^{\mathrm{dR},\infty} E,\rho_E^{\mathbb F})
\]
to \((E,\rho_F)\). The coalgebra-morphism equation is precisely
\[
  \rho_Fs_F\simeq J^\infty(s_F)\rho_E^{\mathbb F}.
\]
\end{enumerate}
\end{proof}

The quotient, disk, jet, and relation descriptions have now been identified as
four interfaces to one descent problem. We can therefore introduce the
intrinsic category used later without adding another layer of structure.

\subsection{Intrinsic formally integrable PDEs}
\label{sec:PDE-topos}

We use the quotient slice for objects whose infinitesimal compatibility is
already built in. The terminology must be kept distinct from the ambient
equation calculus introduced in the next section: an intrinsic formally
integrable PDE is already a descended object, whereas a generalized jet
equation will merely be a map into an infinite jet object and may still fail
to preserve crystalline transport.
\begin{definition}[Intrinsic formally integrable PDE]
Define
\[
  \operatorname{PDE}^{\mathrm{fi}}_{X/S}
  :=
  (\mathcal H_{\Sp})_{/Q_{X/S}}.
\]
An object of this slice is an \textbf{intrinsic formally integrable spectral
PDE}.

Here \emph{formally integrable} means that the corresponding object over \(X\)
carries coherent effective descent along the finite-infinitesimal de Rham
relation
\[
  \mathcal R^{\mathrm{dR}}_{X/S,\bullet}
  \simeq
  \check C_\bullet(X\to Q_{X/S}).
\]
Equivalently, it carries the associated formal-disk action, jet coaction, or
crystalline transport. The terminology asserts this infinitesimal descent
property; it does not assert analytic existence, regularity, or convergence of
actual solutions.
\end{definition}
By Theorem~\ref{thm:effective-jet-descent}, the intrinsic category has the two
equivalent presentations over \(X\):
\[
  \operatorname{PDE}^{\mathrm{fi}}_{X/S}
  \simeq
  \operatorname{Alg}_{\mathsf{Disk}^{\mathrm{dR},\infty}}
  ((\mathcal H_{\Sp})_{/X})
  \simeq
  \operatorname{Coalg}_{J^\infty}
  ((\mathcal H_{\Sp})_{/X}).
\]
Thus ``intrinsic PDE'', ``disk algebra'', and ``jet coalgebra'' describe the
same already-integrable object through different interfaces. They should not
be conflated with a generalized equation \(E\to J^\infty Y\), which is an
ambient presentation whose compatibility will be tested later by Cartan
prolongation.
\begin{remark}[Relation to crystals and nonlinear \(D\)-geometry]
The slice over \(Q_{X/S}\) is the nonlinear descent category associated to the
de Rham relation groupoid. This is formally parallel to crystals over a
derived de Rham prestack and to nonlinear \(D\)-geometry
\cite{GaitsgoryRozenblyumCrystals,BeilinsonDrinfeld2004}. When the reduced-classical comparison of
Chapter~\ref{chap:spectral-de-Rham} is invertible, this description agrees with
the corresponding reduced-classical de Rham-prestack presentation.
\end{remark}

\begin{proposition}[Underlying object detects intrinsic morphisms]
\label{prop:underlying-object-detects-intrinsic-morphisms}
Let \(F,G\to Q_{X/S}\) be intrinsic formally integrable PDEs and put
\[
  E:=e_{X/S}^*F,
  \qquad
  E':=e_{X/S}^*G.
\]
Let
\[
  \theta_F:\mathsf s^*E\xrightarrow{\simeq}\mathsf t^*E,
  \qquad
  \theta_G:\mathsf s^*E'\xrightarrow{\simeq}\mathsf t^*E'
\]
be their canonical crystalline transports.

For a map
\[
  f:E\to E'
\]
in \((\mathcal H_{\Sp})_{/X}\), the following are equivalent:
\begin{enumerate}
\item \(f\) descends to a map
\[
  F\to G
\]
over \(Q_{X/S}\);
\item \(f\) extends to a morphism between the corresponding effective
\v{C}ech-descent objects;
\item \(f\) preserves crystalline transport together with its induced higher
\v{C}ech coherence; in degree one this includes a specified homotopy filling
the square
\[
\begin{tikzcd}
\mathsf s^*E \ar[r,"\theta_F"] \ar[d,"\mathsf s^*f"']
  & \mathsf t^*E \ar[d,"\mathsf t^*f"]\\
\mathsf s^*E' \ar[r,"\theta_G"']
  & \mathsf t^*E',
\end{tikzcd}
\]
compatible with the cocycle on
\(\mathcal R^{\mathrm{dR}}_{X/S,2}\) and the remaining simplicial data;
\item \(f\) is a \(\mathsf{Disk}^{\mathrm{dR},\infty}\)-algebra morphism;
\item \(f\) is a \(J^\infty\)-coalgebra morphism.
\end{enumerate}
Thus preserving crystalline transport means preserving the effective-descent
datum encoded by the degree-one square and its simplicial compatibilities.
\end{proposition}

\begin{proof}
\leavevmode
\begin{enumerate}[label=\textup{(\arabic*)}]
    \item \textbf{Use the Čech totalization.}
Proposition~\ref{prop:effective-cech-descent-Q} gives a canonical equivalence
\[
  (\mathcal H_{\Sp})_{/Q_{X/S}}
  \xrightarrow{\simeq}
  \operatorname*{Tot}_{[n]\in\Delta}
  (\mathcal H_{\Sp})_{/\mathcal R^{\mathrm{dR}}_{X/S,n}}.
\]
Under this equivalence, the objects \(F\) and \(G\) determine their underlying
degree-zero objects \(E\) and \(E'\), the degree-one transports
\(\theta_F\) and \(\theta_G\), their degree-two cocycle homotopies, and all
higher simplicial coherence simultaneously.

    \item \textbf{Characterize descended morphisms.}
Consequently a map
\[
  f:E\to E'
\]
descends to \(Q_{X/S}\) exactly when it is the degree-zero component of a morphism
between these two objects of the totalization. In degree one this requires the
displayed transport square to commute up to the specified homotopy, and in
higher degrees it requires precisely the compatibility of that homotopy with
the cocycle and the remaining simplicial compatibility. This proves the equivalence
of \textup{(1)}, \textup{(2)}, and \textup{(3)} on the full mapping spaces.

    \item \textbf{Pass to the monadic and comonadic presentations.}
Theorem~\ref{thm:effective-jet-descent} gives equivalences of
infinity-categories over \((\mathcal H_{\Sp})_{/X}\),
\[
  (\mathcal H_{\Sp})_{/Q_{X/S}}
  \simeq
  \operatorname{Alg}_{\mathsf{Disk}^{\mathrm{dR},\infty}}((\mathcal H_{\Sp})_{/X})
  \simeq
  \operatorname{Coalg}_{J^\infty}((\mathcal H_{\Sp})_{/X}).
\]
Under these equivalences a morphism over \(Q_{X/S}\) is carried respectively to a
\(\mathsf{Disk}^{\mathrm{dR},\infty}\)-algebra morphism and a \(J^\infty\)-coalgebra morphism. Hence
\textup{(1)}, \textup{(4)}, and \textup{(5)} are equivalent as well.

    \item \textbf{Retain the full mapping-space coherence.}
These equivalences hold on the complete mapping spaces and therefore retain
the higher compatibilities carried by the totalization.
\end{enumerate}
\end{proof}

The intrinsic descent category is therefore fixed. Quotient-side language will
be used when descent and limits are the issue; coalgebra language will be used
when operators, equations, and prolongation are the issue.

\section{Differential operators and generalized equations}
\label{sec:jets-operators-equations}

We now use the jet comonad as a calculus rather than as a descent classifier.
Its two principal uses have opposite variance. Maps out of a jet object give
differential operators, while maps into a jet object give generalized
equations. Holonomic sections test ordinary solutions of those equations;
formal-disk families test infinitesimal compatibility. This variance split is
the organizing principle of the section.
\[
\begin{tikzcd}[column sep=huge,row sep=large]
&
J^\infty_{X/S}
  \ar[dl,"{\operatorname{coKl}(-)}"']
  \ar[dr,"{(\id\downarrow -)}"]
&
\\
\operatorname{DiffOp}_{X/S}
  \ar[d,"{\text{infinite prolongation}}"']
&&
\operatorname{Eqn}_{X/S}
  \ar[dl,"{\text{actual solutions}}"']
  \ar[d,"{\text{formal-disk solutions}}"]
\\
\operatorname{Coalg}_{J^\infty}
&
\operatorname{Sol}
&
\operatorname{FSol}.
\end{tikzcd}
\]
The diagram records roles rather than new functors: the constructions
themselves are defined below.
\subsection{Free and cofree jet coalgebras and holonomicity}
\label{sec:free-cofree-jets}
The calculus uses two universal coalgebras. The disk object
\(\mathbb F(Z)\) carries the comparison coaction constructed above, and
\(\mathbb K(E)\) is cofree on an arbitrary object over \(X\). Their adjunctions
identify the terminal coalgebra and hence the holonomic prolongation of an
ordinary section.
The equation calculus uses two universal coalgebras: the disk coalgebra \(\mathbb F(Z)\) introduced above and the cofree coalgebra on an arbitrary object over \(X\). Their adjunctions also identify the canonical coalgebra on \(X\), from which holonomic prolongation follows.

\begin{construction}[Cofree jet coalgebras]
For \(E\in(\mathcal H_{\Sp})_{/X}\), define
\[
  \mathbb K(E):=(J^\infty E,\delta^J_E).
\]
The disk coalgebra \(\mathbb F(Z)\) was constructed in
Construction~\ref{con:free-disk-coalgebra}; \(\mathbb K(E)\) is the cofree jet coalgebra.
\end{construction}

\begin{theorem}[Free--forgetful--cofree adjunctions]
\label{thm:free-forgetful-cofree}
Let
\[
  U_J:\operatorname{Coalg}_{J^\infty}((\mathcal H_{\Sp})_{/X})\to(\mathcal H_{\Sp})_{/X}
\]
be the forgetful functor. There are canonical adjunctions
\[
  \mathbb F\dashv U_J\dashv\mathbb K.
\]
They may be displayed as
\[
\begin{tikzcd}[column sep=huge]
(\mathcal H_{\Sp})_{/X}
  \ar[r,shift left=1.4ex,"\mathbb F"]
  \ar[r,shift right=1.4ex,"\mathbb K"']
&
\operatorname{Coalg}_{J^\infty}((\mathcal H_{\Sp})_{/X})
  \ar[l,"U_J" description]
\end{tikzcd}
\]
with \(\mathbb F\dashv U_J\) and \(U_J\dashv\mathbb K\).
In particular,
\[
  \Map_{\operatorname{Coalg}_{J^\infty}((\mathcal H_{\Sp})_{/X})}(\mathbb F(Z),C)
  \simeq
  \Map_X(Z,U_JC),
\]
and
\[
  \Map_{\operatorname{Coalg}_{J^\infty}((\mathcal H_{\Sp})_{/X})}(C,\mathbb K(E))
  \simeq
  \Map_X(U_JC,E).
\]
Under effective descent, \(\mathbb F(Z)\) corresponds to
\(\Sigma_{e_{X/S}} Z\to Q_{X/S}\), while \(\mathbb K(E)\) corresponds to
\(\Pi_{e_{X/S}} E\to Q_{X/S}\).
\end{theorem}

\begin{proof}
\leavevmode
\begin{enumerate}[label=\textup{(\arabic*)}]
    \item \textbf{Recall the Eilenberg--Moore adjunctions.}
The Eilenberg--Moore construction for monads and comonads
\cite{LurieHigherAlgebra} gives the free and cofree adjunctions used here.

    \item \textbf{Construct the free disk algebra.}
The monad \(\mathsf{Disk}^{\mathrm{dR},\infty}\) has the canonical free-algebra functor
\[
  Z\longmapsto(\mathsf{Disk}^{\mathrm{dR},\infty} Z,\mu_Z)
\]
left adjoint to the forgetful functor
\[
  \operatorname{Alg}_{\mathsf{Disk}^{\mathrm{dR},\infty}}((\mathcal H_{\Sp})_{/X})\to(\mathcal H_{\Sp})_{/X}.
\]

    \item \textbf{Compute its jet-coalgebra mate.}
By Corollary~\ref{cor:comparison-actions-coactions-mates} applied to the object
\(\Sigma_{e_{X/S}} Z\in(\mathcal H_{\Sp})_{/Q_{X/S}}\), the mate of
\[
  \mu_Z
  =
  e_{X/S}^*\varepsilon^\Sigma_{\Sigma_{e_{X/S}} Z}
\]
is
\[
  \rho_Z^{\mathbb F}
  =
  e_{X/S}^*\eta^\Pi_{\Sigma_{e_{X/S}} Z}.
\]

    \item \textbf{Identify the free jet coalgebra.}
Therefore Theorem~\ref{thm:mate-em-equivalence} carries the free
\(\mathsf{Disk}^{\mathrm{dR},\infty}\)-algebra on \(Z\) to
\[
  \mathbb F(Z)=(\mathsf{Disk}^{\mathrm{dR},\infty} Z,\rho_Z^{\mathbb F}).
\]

    \item \textbf{Transport the free--forgetful adjunction.}
Since the mate equivalence is over \((\mathcal H_{\Sp})_{/X}\), it transports the
free--forgetful adjunction to
\[
  \mathbb F\dashv U_J.
\]

    \item \textbf{Construct the cofree coalgebra.}
For the comonad \(J^\infty\), the cofree coalgebra on \(E\) is
\[
  \mathbb K(E)=(J^\infty E,\delta^J_E),
\]
with counit \(\varepsilon^\Pi_E\), and it is right adjoint to \(U_J\).

    \item \textbf{Assemble the adjoint triple.}
This proves
\[
  \mathbb F\dashv U_J\dashv\mathbb K
\]
and the displayed mapping-space equivalences.

    \item \textbf{Identify the descended adjoints.}
Finally, under
Theorem~\ref{thm:effective-jet-descent}, the left and right adjoints of
\(e_{X/S}^*\) are \(\Sigma_{e_{X/S}}\) and \(\Pi_{e_{X/S}}\). The uniqueness
of adjoints identifies the descended objects corresponding to
\(\mathbb F(Z)\) and \(\mathbb K(E)\) with
\(\Sigma_{e_{X/S}} Z\) and \(\Pi_{e_{X/S}} E\), respectively.
\end{enumerate}
\end{proof}

The theorem now justifies the terminology used earlier:
\(\mathbb F(Z)\) is the free jet coalgebra on \(Z\), whereas
\(\mathbb K(E)\) is cofree on \(E\). From this point onward the adjectives
``free'' and ``cofree'' refer to these proved universal properties.

\begin{proposition}[Terminal jet coalgebra]
The terminal object
\[
  \id_{Q_{X/S}}:Q_{X/S}\to Q_{X/S}
\]
of \((\mathcal H_{\Sp})_{/Q_{X/S}}\) pulls back to the terminal object
\[
  \id_X:X\to X
\]
of \((\mathcal H_{\Sp})_{/X}\). Its comonadic comparison structure is a canonical
coaction
\[
  \rho_X:X\longrightarrow J^\infty X,
\]
and \((X,\rho_X)\) is terminal in
\(\operatorname{Coalg}_{J^\infty}((\mathcal H_{\Sp})_{/X})\). The counit
\[
  \varepsilon^\Pi_X:J^\infty X\to X
\]
is an equivalence with inverse \(\rho_X\).
\end{proposition}

\begin{proof}
\leavevmode
\begin{enumerate}[label=\textup{(\arabic*)}]
    \item \textbf{Pull back the terminal descent object.}
Under Theorem~\ref{thm:effective-jet-descent},
\[
  (\mathcal H_{\Sp})_{/Q_{X/S}}
  \simeq
  \operatorname{Coalg}_{J^\infty}((\mathcal H_{\Sp})_{/X}),
\]
and the underlying-object functor corresponds to \(e_{X/S}^*\). The terminal
object of \((\mathcal H_{\Sp})_{/Q_{X/S}}\) therefore determines a terminal coalgebra whose
underlying object is
\[
  e_{X/S}^*(\id_{Q_{X/S}})=\id_X.
\]

    \item \textbf{Apply the coalgebra counit equation.}
The coalgebra counit equation gives
\[
  \varepsilon^\Pi_X\rho_X\simeq\id_X.
\]

    \item \textbf{Use terminality of the jet object.}
By Proposition~\ref{prop:jet-limits}, \(J^\infty X\) is terminal in
\((\mathcal H_{\Sp})_{/X}\).

    \item \textbf{Identify the reverse composite.}
Hence the reverse composite is the contractibly unique
endomorphism of a terminal object and is equivalent to the identity.

    \item \textbf{Conclude invertibility of the counit.}
Thus
\(\rho_X\simeq\left(\varepsilon^\Pi_X\right)^{-1}\).
\end{enumerate}
\end{proof}

\begin{definition}[Infinite prolongation of a section]
For a section \(\phi:X\to E\), define
\[
  j^\infty\phi
  :=
  J^\infty(\phi)\rho_X:
  X\to J^\infty E.
\]
\end{definition}

Naturality of the counit gives
\[
  \varepsilon^\Pi_Ej^\infty\phi\simeq\phi.
\]

\begin{definition}[Holonomic section]
The terminology follows the infinite-jet formulation of holonomicity in
geometric PDE theory \cite{KhavkineSchreiber2017}. A section \(\sigma:X\to J^\infty E\) is \textbf{holonomic} if it is a
jet-coalgebra morphism
\[
  (X,\rho_X)\longrightarrow\mathbb K(E).
\]
\end{definition}

\begin{theorem}[Holonomic recognition]
\label{thm:holonomic-recognition}
There is a canonical equivalence
\[
  \Map_{\operatorname{Coalg}_{J^\infty}((\mathcal H_{\Sp})_{/X})}
  \bigl((X,\rho_X),\mathbb K(E)\bigr)
  \simeq
  \Map_X(X,E).
\]
Under this equivalence \(\phi\) corresponds to \(j^\infty\phi\). The
holonomic section fits into the commutative triangle
\[
\begin{tikzcd}[column sep=huge,row sep=large]
X \ar[rr,"\sigma"] \ar[dr,"\varepsilon^\Pi_E\sigma"']
&& J^\infty E \ar[dl,"\varepsilon^\Pi_E"]\\
& E &
\end{tikzcd}
\]
and is holonomic precisely when
\(\sigma\simeq j^\infty(\varepsilon^\Pi_E\sigma)\).
\end{theorem}

\begin{proof}
\leavevmode
\begin{enumerate}[label=\textup{(\arabic*)}]
    \item \textbf{Apply the cofree adjunction.}
Apply the cofree adjunction \(U_J\dashv\mathbb K\) to the terminal coalgebra.
\end{enumerate}
\end{proof}

\subsection{Differential operators as co-Kleisli morphisms}
\label{sec:differential-operators}

The cofree adjunction turns the jet comonad into the co-Kleisli calculus of
operators. A differential operator is therefore represented by a map
\(J^\infty E\to F\); comultiplication supplies composition, and infinite
prolongation is simply the corresponding morphism between cofree
coalgebras.

\begin{definition}[Crystalline differential operator]
Following the co-Kleisli formulation of infinite-jet operators used in
categorical treatments of PDEs
\cite{Marvan1986,KhavkineSchreiber2017}, for
\(E,F\in(\mathcal H_{\Sp})_{/X}\) a crystalline differential operator
\[
  D:E\rightsquigarrow F
\]
is a morphism
\[
  \widehat D:J^\infty E\to F
\]
in \((\mathcal H_{\Sp})_{/X}\). Its action on a section \(\phi:X\to E\) is
\[
  D(\phi):=\widehat D\,j^\infty\phi.
\]
The action is represented by the composite
\[
\begin{tikzcd}[column sep=huge,row sep=large]
X \ar[r,"j^\infty\phi"] \ar[dr,"D(\phi)"']
&
J^\infty E \ar[d,"\widehat D"]\\
& F .
\end{tikzcd}
\]
\end{definition}

\begin{construction}[Differential-operator infinity-category]
Define \(\operatorname{DiffOp}_{X/S}\) to have the same objects as \((\mathcal H_{\Sp})_{/X}\) and mapping
spaces
\[
  \Map_{\operatorname{DiffOp}_{X/S}}(E,F)
  :=
  \Map_{\operatorname{Coalg}_{J^\infty}((\mathcal H_{\Sp})_{/X})}
  \bigl(\mathbb K(E),\mathbb K(F)\bigr).
\]
Equivalently, \(\operatorname{DiffOp}_{X/S}\) is the full infinity-subcategory of jet
coalgebras spanned by the cofree objects, with each cofree object \(\mathbb K(E)\) relabelled by \(E\).
\end{construction}

\begin{proposition}[Co-Kleisli identification]
\label{prop:cokleisli}
There are canonical equivalences
\[
  \Map_{\operatorname{DiffOp}_{X/S}}(E,F)
  \simeq
  \Map_X(J^\infty E,F).
\]
Under these equivalences the identity is represented by \(\varepsilon^\Pi_E\), and
the composite of
\[
  \widehat D_1:J^\infty E_1\to E_2,
  \qquad
  \widehat D_2:J^\infty E_2\to E_3
\]
is represented by
\[
  \widehat D_2\,J^\infty(\widehat D_1)\delta^J_{E_1}.
\]
Hence
\[
  \operatorname{DiffOp}_{X/S}\simeq\operatorname{coKl}(J^\infty_{X/S}).
\]
For two composable operators, the co-Kleisli composite is represented by
\[
\begin{tikzcd}[column sep=huge,row sep=large]
J^\infty E_1
  \ar[r,"\delta^J_{E_1}"]
  \ar[drr,"\widehat D_2\circ_{\mathrm{coKl}}\widehat D_1"']
&
J^\infty J^\infty E_1
  \ar[r,"J^\infty(\widehat D_1)"]
&
J^\infty E_2
  \ar[d,"\widehat D_2"]\\
&& E_3 .
\end{tikzcd}
\]
\end{proposition}

\begin{proof}
\leavevmode
\begin{enumerate}[label=\textup{(\arabic*)}]
    \item \textbf{Apply the cofree adjunction.}
Apply the cofree adjunction to maps between cofree coalgebras.

    \item \textbf{Identify co-Kleisli composition.}
Composition is
composition in the coalgebra infinity-category, whose classifying map is the co-Kleisli composition formula.

    \item \textbf{Inherit the higher coherence.}
All higher associativity and unit coherence is
therefore inherited from that infinity-category.
\end{enumerate}
\end{proof}

\begin{proposition}[Order-zero operators]
Every morphism \(f:E\to F\) in \((\mathcal H_{\Sp})_{/X}\) determines an order-zero
operator
\[
  \widehat f^{(0)}:=f\varepsilon^\Pi_E:J^\infty E\to F.
\]
These assignments form an identity-on-objects functor
\[
  \operatorname{ord}_0:(\mathcal H_{\Sp})_{/X}\to\operatorname{DiffOp}_{X/S}.
\]
Full faithfulness of \(\operatorname{ord}_0\) requires additional hypotheses beyond the unrestricted slice.
\end{proposition}

\begin{proof}
\leavevmode
\begin{enumerate}[label=\textup{(\arabic*)}]
    \item \textbf{Construct the order-zero functor.}
The comonad counit defines the canonical identity-on-objects functor to the
co-Kleisli infinity-category.

    \item \textbf{Verify compatibility with composition.}
Naturality of the counit and the comonad laws
identify the co-Kleisli composite of two order-zero maps with the order-zero map
of their ordinary composite.
\end{enumerate}
\end{proof}

\begin{definition}[Infinite prolongation of an operator]
Let \(D:E\rightsquigarrow F\) be a differential operator represented by
\[
  \widehat D:J^\infty E\to F.
\]
Define its infinite prolongation by
\[
  \operatorname{prol}_{\mathrm{op}}^\infty(\widehat D)
  :=
  J^\infty(\widehat D)\delta^J_E:
  J^\infty E\to J^\infty F.
\]
\end{definition}

\begin{theorem}[Cofree realization of differential operators]
Let \(D:E\rightsquigarrow F\) be represented by
\[
  \widehat D:J^\infty E\to F.
\]
Then \(\operatorname{prol}_{\mathrm{op}}^\infty(\widehat D)\) is the contractibly unique
cofree jet-coalgebra morphism
\[
  \mathbb K(E)\to\mathbb K(F)
\]
classified by \(\widehat D\). It satisfies
\[
  \varepsilon^\Pi_F\operatorname{prol}_{\mathrm{op}}^\infty(\widehat D)\simeq\widehat D,
\]
and infinite prolongation identifies \(\operatorname{DiffOp}_{X/S}\) fully faithfully with the
full subcategory of cofree jet coalgebras.
\end{theorem}

\begin{proof}
\leavevmode
\begin{enumerate}[label=\textup{(\arabic*)}]
    \item \textbf{Apply the cofree adjunction.}
By the cofree adjunction \(U_J\dashv\mathbb K\), the morphism classified by
\(\widehat D\) has underlying map
\(J^\infty(\widehat D)\delta^J_E\); the displayed identity is the comonad
counit law.
\end{enumerate}
\end{proof}

\begin{proposition}[Prolongation respects co-Kleisli composition]
\label{prop:prolongation-composition}
For composable differential operators
\[
  D_1:E_1\rightsquigarrow E_2,
  \qquad
  D_2:E_2\rightsquigarrow E_3,
\]
represented by
\[
  \widehat D_1:J^\infty E_1\to E_2,
  \qquad
  \widehat D_2:J^\infty E_2\to E_3,
\]
one has a canonical equivalence of jet-coalgebra morphisms
\[
  \operatorname{prol}_{\mathrm{op}}^\infty
  (\widehat D_2\circ_{\mathrm{coKl}}\widehat D_1)
  \simeq
  \operatorname{prol}_{\mathrm{op}}^\infty(\widehat D_2)\circ \operatorname{prol}_{\mathrm{op}}^\infty(\widehat D_1).
\]
For every section \(\phi:X\to E_1\),
\[
  (D_2\circ_{\mathrm{coKl}}D_1)(\phi)
  \simeq
  \widehat D_2J^\infty(\widehat D_1)
  \delta^J_{E_1}j^\infty\phi.
\]
\end{proposition}

\begin{proof}
\leavevmode
\begin{enumerate}[label=\textup{(\arabic*)}]
    \item \textbf{Identify the classified composite.}
Composition of the two cofree coalgebra morphisms is the cofree morphism whose
classifying map is the co-Kleisli composite.

    \item \textbf{Evaluate on a section.}
The section formula is obtained
by precomposition with \(j^\infty\phi\).
\end{enumerate}
\end{proof}
Differential operators are thus maps \emph{out of} jet objects, with infinite
prolongation realizing their co-Kleisli composition coalgebraically. The
second use of \(J^\infty\) reverses the variance: an equation is presented by a
map \emph{into} a jet object.

\subsection{Generalized jet equations}
\label{sec:generalized-PDEs}

Homotopy equalizers and derived zero loci of differential operators both
produce morphisms into \(J^\infty Y\). Rather than force such derived
equations to be embedded subobjects, we retain the full morphism
\(E\to J^\infty Y\) and organize these presentations in a comma
\(\infty\)-category.
\begin{construction}[Equation cut out by differential operators]
Let
\[
  D_0,D_1:Y\rightsquigarrow F
\]
be differential operators represented by
\[
  \widehat D_0,\widehat D_1:J^\infty Y\rightrightarrows F.
\]
Define the derived equation object
\[
  E(D_0,D_1)
  :=
  \operatorname{Eq}(\widehat D_0,\widehat D_1)
\]
with its canonical map to \(J^\infty Y\). This is the homotopy equalizer in
\((\mathcal H_{\Sp})_{/X}\).
\end{construction}

\begin{construction}[Derived zero locus]
Let \(W\to X\) be equipped with a zero section
\[
  0:X\to W,
\]
and let \(f:J^\infty Y\to W\). Define
\[
  Z(f):=J^\infty Y\times_WX.
\]
It is characterized by the Cartesian square
\[
\begin{tikzcd}[column sep=huge,row sep=large]
Z(f) \ar[r] \ar[d]
&
J^\infty Y \ar[d,"f"]\\
X \ar[r,"0"']
&
W .
\end{tikzcd}
\]
The displayed homotopy fibre is the derived zero locus.
\end{construction}

\begin{definition}[Generalized and embedded jet equation]
The formulation by morphisms into an infinite jet object is the derived
analogue of generalized equations in geometric PDE theory
\cite{Marvan1986,KhavkineSchreiber2017}. A \textbf{generalized jet equation} on \(Y\) relative to \(X/S\) is a
morphism
\[
  i:E\longrightarrow J^\infty Y
\]
in \((\mathcal H_{\Sp})_{/X}\). It is \textbf{embedded} if \(i\) is a monomorphism.
\end{definition}

\begin{construction}[Infinity-category of generalized jet equations]
\label{con:equation-category}
Define
\[
  \operatorname{Eqn}_{X/S}
  :=
  \bigl(\id_{(\mathcal H_{\Sp})_{/X}}\downarrow J^\infty_{X/S}\bigr).
\]
Thus an object is a triple
\[
  (Y,E,i),
  \qquad
  i:E\longrightarrow J^\infty_{X/S}Y,
\]
and a morphism
\[
  (Y,E,i)\longrightarrow(Y',E',i')
\]
consists of maps
\[
  f:Y\to Y',
  \qquad
  g:E\to E',
\]
together with a commutative square
\[
\begin{tikzcd}
E \ar[r,"i"] \ar[d,"g"']
  & J^\infty Y \ar[d,"J^\infty(f)"]\\
E' \ar[r,"i'"']
  & J^\infty Y'.
\end{tikzcd}
\]
Equivalently,
\[
  \operatorname{Eqn}_{X/S}
  \simeq
  \operatorname{Fun}(\Delta^1,(\mathcal H_{\Sp})_{/X})
  \times_{(\mathcal H_{\Sp})_{/X}\times(\mathcal H_{\Sp})_{/X}}
  ((\mathcal H_{\Sp})_{/X}\times(\mathcal H_{\Sp})_{/X}),
\]
where the first map is
\((\operatorname{ev}_0,\operatorname{ev}_1)\) and the second sends
\((E,Y)\) to \((E,J^\infty Y)\). Hence the equation infinity-category is
constructed by a finite limit in \(\operatorname{Cat}_\infty\).

For fixed \(Y\), write
\[
  \operatorname{Eqn}_{X/S}(Y)
  :=
  \operatorname{fib}_Y
  \left(
    \operatorname{Eqn}_{X/S}
    \longrightarrow
    (\mathcal H_{\Sp})_{/X}
  \right).
\]
Equivalently,
\[
  \operatorname{Eqn}_{X/S}(Y)
  \simeq
  ((\mathcal H_{\Sp})_{/X})_{/J^\infty Y}.
\]
\end{construction}

\begin{remark}[Generalized and embedded equations]
A monomorphism in an infinity-topos is \((-1)\)-truncated
\cite{LurieHTT}. Derived zero loci and homotopy equalizers naturally carry
higher homotopy data. Generalized jet equations retain this data by allowing
arbitrary morphisms into a jet object, while embedded equations form the
monomorphism-valued special case.
\end{remark}

\begin{proposition}[Pullback and simultaneous systems]
Let \(f:Y\to Z\) and let \(i:F\to J^\infty Z\) be a generalized jet equation. Its
pullback to \(Y\) is
\[
  J^\infty Y\times_{J^\infty Z}F
  \longrightarrow
  J^\infty Y.
\]
If \(i\) is embedded, so is its pullback. More generally, any small family
\[
  \{E_\alpha\to J^\infty Y\}_{\alpha\in I}
\]
has a simultaneous system given by its limit in
\(((\mathcal H_{\Sp})_{/X})_{/J^\infty Y}\). For embedded members this is their
intersection as subobjects.
\end{proposition}

\begin{proof}
\leavevmode
\begin{enumerate}[label=\textup{(\arabic*)}]
    \item \textbf{Construct the pulled-back equation.}
The first assertion is functoriality of \(J^\infty\) followed by pullback.

    \item \textbf{Preserve embeddedness.}
Monomorphisms are pullback-stable.

    \item \textbf{Construct simultaneous systems.}
Slices of an infinity-topos have all small
limits, and limits of monomorphisms are monomorphisms.
\end{enumerate}
\end{proof}

\begin{proposition}[Embedded predicates are internally classified]
If
\[
  i:E\hookrightarrow J^\infty Y
\]
is embedded, then it is classified by a morphism
\[
  \chi_i:J^\infty Y\longrightarrow\Omega_X^{\mathrm{sub}}
\]
to the subobject classifier of \((\mathcal H_{\Sp})_{/X}\). Thus embedded equations are
exactly proposition-valued differential predicates on the infinite jet object.
Generalized jet equations retain the full morphism data into \(J^\infty Y\).
\end{proposition}

\begin{proof}
\leavevmode
\begin{enumerate}[label=\textup{(\arabic*)}]
    \item \textbf{Apply the subobject-classifier theorem.}
The slice \((\mathcal H_{\Sp})_{/X}\) is an infinity-topos, and its
monomorphisms into a fixed object are classified by its subobject classifier
\cite{LurieHTT}.
\end{enumerate}
\end{proof}

The equation category separates the ambient variable \(Y\) from the equation
object \(E\) while retaining derived intersections and higher homotopy. It
still says nothing about whether \(E\) itself is transported coherently along
the de Rham relation. That is the role of Cartan structure.

\subsection{Cartan equations as coalgebra morphisms}
\label{sec:cartan-presentations}

The structured subclass consists of generalized equations whose source already
carries a jet-coalgebra presentation and whose map into the ambient jet object
preserves crystalline transport.

\begin{definition}[Jet-coalgebra presentation]
A \textbf{jet-coalgebra presentation} over \(X/S\) is a jet coalgebra
\[
  (E,\rho),
  \qquad
  \rho:E\longrightarrow J^\infty E.
\]
By Theorem~\ref{thm:effective-jet-descent}, every such structure is
equivalently the pullback presentation of an intrinsic formally integrable PDE
\[
  F\longrightarrow Q_{X/S}.
\]
\end{definition}

The term ``Cartan'' is reserved for this compatibility of an equation morphism
with the already-defined crystalline transport. Thus \emph{crystalline}
describes descent along the effective de Rham relation, while \emph{Cartan}
describes preservation of that descent by an equation morphism. The cofree
adjunction will recognize such morphisms as ordinary maps from the transported
source to the ambient object.

\begin{definition}[Cartan equation morphism]
Let \((E,\rho)\) be a jet coalgebra and let \(Y\to X\). A
map
\[
  y:E\to Y
\]
constructs the Cartan equation morphism
\[
  i_y:=J^\infty(y)\rho:E\to J^\infty Y.
\]
The counit identity gives the commutative diagram
\[
\begin{tikzcd}[column sep=huge,row sep=large]
E \ar[r,"\rho"] \ar[d,"y"']
&
J^\infty E \ar[r,"J^\infty(y)"]
&
J^\infty Y \ar[d,"\varepsilon^\Pi_Y"]\\
Y \ar[rr,equal]
&& Y .
\end{tikzcd}
\]
It is \textbf{embedded} if \(i_y\) is a monomorphism.
\end{definition}

\begin{theorem}[Cofree recognition of Cartan equation morphisms]
\label{thm:cartan-recognition}
For a jet coalgebra \((E,\rho)\) and \(Y\in(\mathcal H_{\Sp})_{/X}\), the cofree
adjunction gives
\[
  \Map_{\operatorname{Coalg}_{J^\infty}((\mathcal H_{\Sp})_{/X})}
  \bigl((E,\rho),\mathbb K(Y)\bigr)
  \simeq
  \Map_X(E,Y).
\]
The coalgebra map classified by \(y:E\to Y\) has underlying map
\[
  J^\infty(y)\rho:E\to J^\infty Y.
\]
Conversely every coalgebra map
\[
  i:(E,\rho)\to\mathbb K(Y)
\]
is canonically of this form with classifying map \(y=\varepsilon^\Pi_Yi\).
\end{theorem}

\begin{proof}
\leavevmode
\begin{enumerate}[label=\textup{(\arabic*)}]
    \item \textbf{Use the cofree mapping-space equivalence.}
The mapping-space equivalence is the cofree adjunction
\(U_J\dashv\mathbb K\), as in the Eilenberg--Moore formalism
\cite{LurieHigherAlgebra}.

    \item \textbf{Recover the classifying map.}
If \(i\) is a
coalgebra map, then
\[
\begin{aligned}
  J^\infty(\varepsilon^\Pi_Yi)\rho
  &\simeq
  J^\infty(\varepsilon^\Pi_Y)J^\infty(i)\rho\\
  &\simeq
  J^\infty(\varepsilon^\Pi_Y)\delta^J_Yi\\
  &\simeq i
\end{aligned}
\]
by the coalgebra-morphism identity and the comonad counit law.
\end{enumerate}
\end{proof}

Cartan equations are therefore the generalized equations whose sources already
carry crystalline transport and whose ambient equation maps preserve it. The
next question is different: how should such equations be tested against
sections?

\subsection{Actual and parameterized solutions}
\label{sec:actual-solutions}

An actual solution compares an ordinary section of \(Y\) with its holonomic
infinite jet and asks for a lift through the equation. The homotopy pullback
below records the section, the lift, and the path identifying the two. Dependent
products package the same condition uniformly in parameters over the base.
\begin{definition}[Solution space]
Let \(i:E\to J^\infty Y\) be a generalized jet equation. Define its solution
space by the homotopy pullback
\[
  \operatorname{Sol}(i)
  :=
  \Map_X(X,Y)
  \times_{\Map_X(X,J^\infty Y)}
  \Map_X(X,E),
\]
where the first map is \(\phi\mapsto j^\infty\phi\) and the second is induced
by \(i\). Equivalently,
\[
\begin{tikzcd}[column sep=huge,row sep=large]
\operatorname{Sol}(i)
  \ar[r]
  \ar[d]
&
\Map_X(X,E) \ar[d,"i_*"]\\
\Map_X(X,Y)
  \ar[r,"j^\infty"']
&
\Map_X(X,J^\infty Y)
\end{tikzcd}
\]
is Cartesian. Thus a point of \(\operatorname{Sol}(i)\) already consists of a section
\(\phi:X\to Y\), a map \(a:X\to E\), and the path
\(ia\simeq j^\infty\phi\).
\end{definition}

\begin{construction}[Cartan structure on a parameterized copy of the base]
\label{con:parameterized-base-coaction}
Let \(t:T\to S\), put
\[
  X_T:=X\times_ST,
\]
and regard \(X_T\) as an object of \((\mathcal H_{\Sp})_{/X}\) through the projection to
\(X\). The base-change equivalence
\[
  Q_{X_T/T}\simeq Q_{X/S}\times_ST
\]
makes \(Q_{X_T/T}\) an object over \(Q_{X/S}\), and
\[
  X_T\simeq e_{X/S}^*Q_{X_T/T}.
\]
The unit of \(e_{X/S}^*\dashv\Pi_{e_{X/S}}\) constructs
\[
  \rho_T:
  X_T\longrightarrow J^\infty_{X/S}X_T.
\]
For a \(T\)-parameterized section \(\phi:X_T\to Y\), define
\[
  j_T^\infty\phi
  :=
  J^\infty(\phi)\rho_T:
  X_T\to J^\infty Y.
\]
\end{construction}

\begin{proposition}[Compatibility with base-changed prolongation]
\label{prop:parameterized-prolongation-basechange}
Let \(g_X:X_T\to X\) be the projection, put \(Y_T:=g_X^*Y\), and let
\(\widetilde\phi:X_T\to Y_T\) correspond to \(\phi:X_T\to Y\). Under the
canonical equivalences
\[
\begin{aligned}
  \Map_X(X_T,J^\infty_{X/S}Y)
  &\simeq
  \Map_{X_T}
  \bigl(X_T,g_X^*J^\infty_{X/S}Y\bigr)\\
  &\simeq
  \Map_{X_T}
  \bigl(X_T,J^\infty_{X_T/T}Y_T\bigr),
\end{aligned}
\]
the map \(j_T^\infty\phi\) corresponds to
\[
  j^\infty\widetilde\phi:
  X_T\to J^\infty_{X_T/T}Y_T.
\]
\end{proposition}

\begin{proof}
\leavevmode
\begin{enumerate}[label=\textup{(\arabic*)}]
    \item \textbf{Base-change the effective quotient.}
Write \(e_{X_T/T}:X_T\to Q_{X_T/T}\) for the base change of \(e_{X/S}\). The
base-change theorem of Chapter~\ref{chap:spectral-de-Rham} gives
\[
  Q_{X_T/T}\simeq Q_{X/S}\times_ST
\]
and therefore an equivalence of relation groupoids
\[
  \mathcal R^{\mathrm{dR}}_{X_T/T,\bullet}
  \simeq
  \mathcal R^{\mathrm{dR}}_{X/S,\bullet}\times_ST.
\]

    \item \textbf{Transpose the parameterized coaction.}
Under
\[
  J^\infty_{X/S}\simeq\Pi_{\mathsf s}\mathsf t^*,
\]
the coaction
\[
  \rho_T:X_T\to J^\infty_{X/S}X_T
\]
is adjoint to the canonical crystalline transport
\[
  \theta_T:\mathsf s^*X_T\xrightarrow{\simeq}\mathsf t^*X_T.
\]

    \item \textbf{Identify transport on the parameter.}
Since \(X_T=X\times_ST\), both sides identify with
\(\mathcal R^{\mathrm{dR}}_{X/S}\times_ST\), and \(\theta_T\) is the canonical identity on the
parameter \(T\).

    \item \textbf{Compute the transpose of parameterized prolongation.}
Consequently the relation transpose of
\[
  j_T^\infty\phi=J^\infty(\phi)\rho_T
\]
is
\[
  \mathsf s^*X_T
  \xrightarrow{\theta_T}
  \mathsf t^*X_T
  \xrightarrow{\mathsf t^*\phi}
  \mathsf t^*Y.
\]

    \item \textbf{Compare the relation transposes.}
Transporting this map through
\(\mathcal R^{\mathrm{dR}}_{X_T/T}\simeq \mathcal R^{\mathrm{dR}}_{X/S}\times_ST\) identifies it with
\[
  v_1^*X_T
  \longrightarrow
  \mathsf t^*Y_T
\]
obtained by restricting \(\widetilde\phi\) to the evaluation vertex. By the
relation formula for the base-changed jet comonad, this is precisely the
transpose of
\[
  j^\infty\widetilde\phi:
  X_T\to J^\infty_{X_T/T}Y_T.
\]

    \item \textbf{Apply right Beck--Chevalley.}
Right Beck--Chevalley identifies
\(g_X^*J^\infty_{X/S}Y\) with
\(J^\infty_{X_T/T}Y_T\), so the two prolongations correspond as asserted.
\end{enumerate}
\end{proof}

\begin{construction}[Parameterized prolongation map over the base]
Since \(p=p_Qe_{X/S}\), composition of dependent products gives
\[
  \Pi_p\simeq\Pi_{p_Q}\Pi_{e_{X/S}}.
\]
Define
\[
  j^\infty_{/S}:
  \Pi_pY\longrightarrow\Pi_pJ^\infty Y
\]
to be \(\Pi_{p_Q}\) applied to the unit
\[
  \Pi_{e_{X/S}} Y
  \longrightarrow
  \Pi_{e_{X/S}}e_{X/S}^*\Pi_{e_{X/S}} Y.
\]
\end{construction}

\begin{definition}[Parameterized solution object]
For \(i:E\to J^\infty Y\), define
\[
  \mathbf{Sol}_{X/S}(i)
  :=
  \Pi_pE
  \times_{\Pi_pJ^\infty Y}
  \Pi_pY
\]
in the slice over \(S\), where the lower map is \(j^\infty_{/S}\).
Thus
\[
\begin{tikzcd}[column sep=huge,row sep=large]
\mathbf{Sol}_{X/S}(i)
  \ar[r]
  \ar[d]
&
\Pi_pE \ar[d,"\Pi_p(i)"]\\
\Pi_pY
  \ar[r,"j^\infty_{/S}"']
&
\Pi_pJ^\infty Y
\end{tikzcd}
\]
is Cartesian in \((\mathcal H_{\Sp})_{/S}\).
\end{definition}

\begin{theorem}[Parameterized solution object]
\label{thm:parameterized-solutions}
For every \(T\to S\), there is a canonical equivalence
\[
  \Map_S\bigl(T,\mathbf{Sol}_{X/S}(i)\bigr)
  \simeq
  \operatorname{Sol}(i_T),
\]
where \(i_T\) is the base-changed equation over \(X_T/T\). Thus
\(\mathbf{Sol}_{X/S}(i)\) represents parameterized solutions.
\end{theorem}

\begin{proof}
\leavevmode
\begin{enumerate}[label=\textup{(\arabic*)}]
    \item \textbf{Apply dependent product.}
Dependent product gives
\[
  \Map_S(T,\Pi_pZ)
  \simeq
  \Map_X(X_T,Z).
\]

    \item \textbf{Identify parameterized prolongation.}
Proposition~\ref{prop:parameterized-prolongation-basechange} identifies the
\(T\)-point of \(j^\infty_{/S}\) with ordinary base-changed jet prolongation.

    \item \textbf{Identify the solution pullback.}
Taking \(T\)-points of the defining pullback of
\(\mathbf{Sol}_{X/S}(i)\) therefore gives the homotopy pullback defining
\(\operatorname{Sol}(i_T)\).
\end{enumerate}
\end{proof}

Actual solutions have now been represented both pointwise and in families, but
they test only holonomic sections over \(X\). Formal integrability requires a
stronger infinitesimal test: the equation must also admit compatible families
over every de Rham formal disk.

\subsection{Formal-disk families and formal solutions}
\label{sec:formal-solutions}

The disk monad supplies those test objects. A formal solution over a parameter
\(Z\) consists of a lift from the disk object into the equation together with
an ambient disk family whose canonical Cartan lift agrees with that equation
map.
\begin{definition}[Parameterized de Rham formal-disk family]
For \(Z,Y\in(\mathcal H_{\Sp})_{/X}\), a \(Z\)-parameterized formal-disk family in
\(Y\) is a morphism
\[
  \phi:\mathsf{Disk}^{\mathrm{dR},\infty} Z\to Y.
\]
Its ambient jet map is the constructed map
\[
  \sigma_\phi
  :=
  J^\infty(\phi)\rho_Z^{\mathbb F}:
  \mathsf{Disk}^{\mathrm{dR},\infty} Z\to J^\infty Y.
\]
\end{definition}

\begin{theorem}[Constructed Cartan lift]
For every
\[
  \phi:\mathsf{Disk}^{\mathrm{dR},\infty} Z\to Y,
\]
the map \(\sigma_\phi\) underlies the contractibly unique cofree jet-coalgebra
morphism
\[
  \widetilde\sigma_\phi:
  \mathbb F(Z)\to\mathbb K(Y)
\]
classified by \(\phi\). Hence
\[
  \Map_X(\mathsf{Disk}^{\mathrm{dR},\infty} Z,Y)
  \simeq
  \Map_{\operatorname{Coalg}_{J^\infty}((\mathcal H_{\Sp})_{/X})}
  \bigl(\mathbb F(Z),\mathbb K(Y)\bigr).
\]
The Cartan structure is therefore determined by \(\phi\).
\end{theorem}

\begin{proof}
\leavevmode
\begin{enumerate}[label=\textup{(\arabic*)}]
    \item \textbf{Apply the cofree adjunction.}
Apply the cofree adjunction to the free jet coalgebra.

    \item \textbf{Identify the underlying Cartan lift.}
The cofree morphism
classified by \(\phi\) has underlying map
\[
  J^\infty(\phi)\rho_Z^{\mathbb F}.
\]
\end{enumerate}
\end{proof}

\begin{definition}[Parameterized formal-solution space]
For a generalized jet equation
\[
  i:E\to J^\infty Y,
\]
define
\[
  \operatorname{FSol}_i(Z)
  :=
  \Map_X(\mathsf{Disk}^{\mathrm{dR},\infty} Z,E)
  \times_{\Map_X(\mathsf{Disk}^{\mathrm{dR},\infty} Z,J^\infty Y)}
  \Map_X(\mathsf{Disk}^{\mathrm{dR},\infty} Z,Y),
\]
where the second map sends \(\phi\) to
\[
  J^\infty(\phi)\rho_Z^{\mathbb F}.
\]
Equivalently, the defining square
\[
\begin{tikzcd}[column sep=huge,row sep=large]
\operatorname{FSol}_i(Z)
  \ar[r]
  \ar[d]
&
\Map_X(\mathsf{Disk}^{\mathrm{dR},\infty} Z,E)
  \ar[d,"i_*"]\\
\Map_X(\mathsf{Disk}^{\mathrm{dR},\infty} Z,Y)
  \ar[r,"\phi\mapsto J^\infty(\phi)\rho_Z^{\mathbb F}"']
&
\Map_X(\mathsf{Disk}^{\mathrm{dR},\infty} Z,J^\infty Y)
\end{tikzcd}
\]
is Cartesian. Thus a point of \(\operatorname{FSol}_i(Z)\) contains maps
\[
  \ell:\mathsf{Disk}^{\mathrm{dR},\infty} Z\to E,
  \qquad
  \phi:\mathsf{Disk}^{\mathrm{dR},\infty} Z\to Y,
\]
together with the compatibility path
\[
  i\ell
  \simeq
  J^\infty(\phi)\rho_Z^{\mathbb F}.
\]
\end{definition}

\begin{remark}[Literal de Rham formal-disk interpretation]
If \(Z=a:T\to X\) is a generalized point, then
\(\mathsf{Disk}^{\mathrm{dR},\infty}(a)\) is the relative de Rham neighbourhood
\(\operatorname{Nbh}^{\mathrm{dR}}_{T/S}(T\times_SX)\). Hence \(\operatorname{FSol}_i(a)\) is literally
the homotopy-pullback space of maps from that de Rham formal disk into the equation and
the ambient object.
\end{remark}

\begin{proposition}[Adjoint formula for a de Rham formal-disk family]
\label{prop:adjoint-cartan-lift}
Let
\[
  \sigma_\phi
  =
  J^\infty(\phi)\rho_Z^{\mathbb F}:
  \mathsf{Disk}^{\mathrm{dR},\infty} Z\to J^\infty Y,
\]
and let
\[
  \sigma_\phi^\sharp:
  Z\to J^\infty J^\infty Y
\]
be its adjunct under
\[
  \mathsf{Disk}^{\mathrm{dR},\infty}\dashv J^\infty.
\]
Then
\[
  \sigma_\phi^\sharp
  \simeq
  \delta^J_Y\sigma_\phi\eta^\Sigma_Z.
\]
\end{proposition}

\begin{proof}
\leavevmode
\begin{enumerate}[label=\textup{(\arabic*)}]
    \item \textbf{Use the coalgebra-morphism identity.}
The cofree map
\[
  \widetilde\sigma_\phi:
  \mathbb F(Z)\to\mathbb K(Y)
\]
is a coalgebra morphism, so
\[
  J^\infty(\sigma_\phi)\rho_Z^{\mathbb F}
  \simeq
  \delta^J_Y\sigma_\phi.
\]

    \item \textbf{Insert the adjunction-unit factorization.}
Using
\[
  \eta^{\mathrm{adj}}_Z
  \simeq
  \rho_Z^{\mathbb F}\eta^\Sigma_Z
\]
from
Proposition~\ref{prop:disk-jet-unit-factorization},
\[
\begin{aligned}
  \sigma_\phi^\sharp
  &=
  J^\infty(\sigma_\phi)\eta^{\mathrm{adj}}_Z\\
  &\simeq
  J^\infty(\sigma_\phi)\rho_Z^{\mathbb F}\eta^\Sigma_Z\\
  &\simeq
  \delta^J_Y\sigma_\phi\eta^\Sigma_Z.
\end{aligned}
\]
\end{enumerate}
\end{proof}

\begin{proposition}[Adjoint form of a formal solution]
\label{prop:adjoint-formal-solution}
Let
\[
  (\ell,\phi,h)
\]
be a point of \(\operatorname{FSol}_i(Z)\), so
\[
  h:
  i\ell\simeq\sigma_\phi,
\]
and let
\[
  \ell^\sharp:
  Z\to J^\infty E
\]
be the adjunct of \(\ell\). The path \(h\) induces a path
\[
  J^\infty(i)\ell^\sharp
  \simeq
  \delta^J_Yi\ell\eta^\Sigma_Z,
\]
and mate calculus gives
\[
  \varepsilon^\Pi_E\ell^\sharp
  \simeq
  \ell\eta^\Sigma_Z.
\]
\end{proposition}

\begin{proof}
\leavevmode
\begin{enumerate}[label=\textup{(\arabic*)}]
    \item \textbf{Identify the adjunct of the equation composite.}
The adjunct of \(i\ell\) is
\[
  J^\infty(i)\ell^\sharp.
\]

    \item \textbf{Transport the formal-solution path.}
Transport the path
\[
  i\ell\simeq\sigma_\phi
\]
through the adjunction and apply
Proposition~\ref{prop:adjoint-cartan-lift}.

    \item \textbf{Apply mate calculus.}
The final identity is
Lemma~\ref{lem:jet-mates}.
\end{enumerate}
\end{proof}

\section{Cartan prolongation and formal integrability}
\label{sec:jets-cartan-formal-integrability}

The preceding section separated four roles: operators are co-Kleisli maps,
generalized equations are maps into jets, actual solutions test holonomic
sections, and formal solutions test de Rham disk families. What remains is to
force an arbitrary generalized equation to carry the crystalline compatibility
detected by the last test.

Cartan prolongation supplies that universal repair. It is a pullback in jet
coalgebras that retains the original ambient equation while imposing
compatibility with the comonad. Its underlying object represents formal
solutions, and the evaluation back to \(E\) measures whether any new
compatibility was actually added. The section is organized around that single
recognition map.
\[
\begin{tikzcd}[column sep=huge,row sep=large]
&
\operatorname{Prol}^{\mathrm{Cart}}(i)
  \ar[dl,"p_1"']
  \ar[dr,"p_2"]
&
\\
\mathbb K(E)
  \ar[dr,"\mathbb K(i)"']
&&
\mathbb K(Y)
  \ar[dl,"\widetilde\delta^J_Y"]
\\
&
\mathbb K(J^\infty Y)
&
\end{tikzcd}
\]
The pullback itself constructs the Cartan structure. The subsequent
prolongation-completeness criterion asks exactly when its evaluation to \(E\)
is already an equivalence; no separate notion of compatibility is introduced.
\subsection{Cartan prolongation and its universal property}
\label{sec:cartan-prolongation}

Fix a generalized jet equation
\[
  i:E\longrightarrow J^\infty Y.
\]

\begin{construction}[Infinite Cartan prolongation]
The cofree functor gives
\[
  \mathbb K(i):
  \mathbb K(E)\to\mathbb K(J^\infty Y).
\]
There is also the cofree coalgebra map
\[
  \widetilde\delta^J_Y:
  \mathbb K(Y)\to\mathbb K(J^\infty Y)
\]
classified by \(\id_{J^\infty Y}\); its underlying map is \(\delta^J_Y\). Define
\[
  \operatorname{Prol}^{\mathrm{Cart}}(i)
  :=
  \mathbb K(E)
  \times_{\mathbb K(J^\infty Y)}
  \mathbb K(Y)
\]
in \(\operatorname{Coalg}_{J^\infty}((\mathcal H_{\Sp})_{/X})\). Since the forgetful functor creates
limits, its underlying object is
\[
  \operatorname{Prol}(i)
  :=
  J^\infty E
  \times_{J^\infty J^\infty Y}
  J^\infty Y.
\]
Write
\[
  p_1:\operatorname{Prol}(i)\to J^\infty E,
  \qquad
  p_2:\operatorname{Prol}(i)\to J^\infty Y
\]
for the two projections. The construction is encoded by the Cartesian diagram
\[
\begin{tikzcd}[column sep=huge,row sep=large]
\operatorname{Prol}(i)
  \ar[r,"p_1"]
  \ar[d,"p_2"']
&
J^\infty E \ar[d,"J^\infty(i)"]\\
J^\infty Y
  \ar[r,"\delta^J_Y"']
&
J^\infty J^\infty Y .
\end{tikzcd}
\]
\end{construction}

\begin{theorem}[Canonical coalgebra on the prolongation]
\label{thm:canonical-prolongation-coalgebra}
The object \(\operatorname{Prol}(i)\) carries the canonical jet-coalgebra structure underlying
\(\operatorname{Prol}^{\mathrm{Cart}}(i)\). Write its coaction as
\[
  \rho_{\operatorname{Prol}(i)}:\operatorname{Prol}(i)\longrightarrow J^\infty \operatorname{Prol}(i).
\]
Both projections \(p_1\) and \(p_2\) are coalgebra maps.
\end{theorem}

\begin{proof}
\leavevmode
\begin{enumerate}[label=\textup{(\arabic*)}]
    \item \textbf{Create the pullback coalgebra.}
Apply the limit-creation statement of Proposition~\ref{prop:jet-limits} to the
defining pullback of coalgebras.
\end{enumerate}
\end{proof}

Fix \(Y\). Define
\[
  \operatorname{CartEq}_{X/S}(Y)
  :=
  \operatorname{Coalg}_{J^\infty}((\mathcal H_{\Sp})_{/X})_{/\mathbb K(Y)}.
\]
Let
\[
  V_Y:
  \operatorname{CartEq}_{X/S}(Y)
  \longrightarrow
  \operatorname{Eqn}_{X/S}(Y)
\]
be the underlying equation-morphism functor.

\begin{theorem}[Cartan prolongation adjunction]
\label{thm:cartan-prolongation-adjunction}
The construction above defines a functor
\[
  \operatorname{Prol}^{\mathrm{Cart}}_Y:\operatorname{Eqn}_{X/S}(Y)\to\operatorname{CartEq}_{X/S}(Y)
\]
and there is a canonical adjunction
\[
  V_Y\dashv \operatorname{Prol}^{\mathrm{Cart}}_Y.
\]
The adjunction is displayed by
\[
\begin{tikzcd}[column sep=huge]
\operatorname{CartEq}_{X/S}(Y)
  \ar[r,shift left=1.4ex,"V_Y"]
&
\operatorname{Eqn}_{X/S}(Y)
  \ar[l,shift left=1.4ex,"\operatorname{Prol}^{\mathrm{Cart}}_Y"]
\end{tikzcd}
\]
with \(V_Y\) left adjoint to Cartan prolongation.
Equivalently, for \(j:C\to\mathbb K(Y)\) in \(\operatorname{CartEq}_{X/S}(Y)\) and
\(i:E\to J^\infty Y\) in \(\operatorname{Eqn}_{X/S}(Y)\),
\[
  \Map_{\operatorname{CartEq}_{X/S}(Y)}(j,\operatorname{Prol}^{\mathrm{Cart}}_Y(i))
  \simeq
  \Map_{\operatorname{Eqn}_{X/S}(Y)}(V_Yj,i).
\]
\end{theorem}

\begin{proof}
\leavevmode
\begin{enumerate}[label=\textup{(\arabic*)}]
    \item \textbf{Unpack maps into the Cartan pullback.}
A map \(j\to \operatorname{Prol}^{\mathrm{Cart}}_Y(i)\) is a coalgebra map
\[
  a:C\to\mathbb K(E)
\]
whose composite with \(\mathbb K(i)\) agrees with
\(\widetilde\delta^J_Yj\).

    \item \textbf{Apply the cofree adjunction.}
The cofree adjunction identifies \(a\) with a map
\[
  \bar a:U_JC\to E.
\]

    \item \textbf{Compare the classifying maps.}
The classifying map of \(\mathbb K(i)a\) is \(i\bar a\), while the classifying
map of \(\widetilde\delta^J_Yj\) is
\[
  \varepsilon^\Pi_{J^\infty Y}\delta^J_YU_Jj
  \simeq U_Jj.
\]

    \item \textbf{Identify the mapping-space adjunction.}
Hence the compatibility is exactly the path
\[
  i\bar a\simeq U_Jj
\]
that defines a map from \(V_Yj\) to \(i\) in the slice \(\operatorname{Eqn}_{X/S}(Y)\). The
construction is natural on the full mapping spaces, giving the adjunction.
\end{enumerate}
\end{proof}

\begin{theorem}[Universal formal-solution object]
\label{thm:universal-formal-solution}
For every \(Z\in(\mathcal H_{\Sp})_{/X}\), there is a natural equivalence
\[
  \Map_X(Z,\operatorname{Prol}(i))
  \simeq
  \operatorname{FSol}_i(Z).
\]
Thus the identity of \(\operatorname{Prol}(i)\) determines the universal formal-solution
family.
\end{theorem}

\begin{proof}
\leavevmode
\begin{enumerate}[label=\textup{(\arabic*)}]
    \item \textbf{Map the free coalgebra into the pullback.}
Using \(\mathbb F\dashv U_J\) and the pullback defining \(\operatorname{Prol}^{\mathrm{Cart}}(i)\),
\[
\begin{aligned}
  \Map_X(Z,\operatorname{Prol}(i))
  &\simeq
  \Map_{\operatorname{Coalg}_{J^\infty}((\mathcal H_{\Sp})_{/X})}(\mathbb F(Z),\operatorname{Prol}^{\mathrm{Cart}}(i))\\
  &\simeq
  \Map_{\operatorname{Coalg}_{J^\infty}((\mathcal H_{\Sp})_{/X})}(\mathbb F(Z),\mathbb K(E))
  \times_{\Map_{\operatorname{Coalg}_{J^\infty}((\mathcal H_{\Sp})_{/X})}(\mathbb F(Z),\mathbb K(J^\infty Y))}
  \Map_{\operatorname{Coalg}_{J^\infty}((\mathcal H_{\Sp})_{/X})}(\mathbb F(Z),\mathbb K(Y))\\
  &\simeq
  \Map_X(\mathsf{Disk}^{\mathrm{dR},\infty} Z,E)
  \times_{\Map_X(\mathsf{Disk}^{\mathrm{dR},\infty} Z,J^\infty Y)}
  \Map_X(\mathsf{Disk}^{\mathrm{dR},\infty} Z,Y).
\end{aligned}
\]

    \item \textbf{Identify the formal-solution space.}
The right side is \(\operatorname{FSol}_i(Z)\).
\end{enumerate}
\end{proof}

Under the formal-disk--jet adjunction, the maps \(p_1\) and \(p_2\) are adjunct to
maps
\[
  \ell_i^{\mathrm{univ}}:\mathsf{Disk}^{\mathrm{dR},\infty} \operatorname{Prol}(i)\to E,
  \qquad
  \phi_i^{\mathrm{univ}}:\mathsf{Disk}^{\mathrm{dR},\infty} \operatorname{Prol}(i)\to Y.
\]
The defining pullback path gives
\[
  i\ell_i^{\mathrm{univ}}
  \simeq
  J^\infty(\phi_i^{\mathrm{univ}})\rho_{\operatorname{Prol}(i)}^{\mathbb F}.
\]
This is the formal-solution family classified by \(\id_{\operatorname{Prol}(i)}\).

\begin{definition}[Prolongation evaluation map]
Define
\[
  c_i:\operatorname{Prol}(i)\to E
\]
by
\[
  c_i:=\varepsilon^\Pi_Ep_1.
\]
Equivalently,
\[
\begin{tikzcd}[column sep=huge,row sep=large]
\operatorname{Prol}(i)
  \ar[rr,"p_1"]
  \ar[dr,"c_i"']
&&
J^\infty E \ar[dl,"\varepsilon^\Pi_E"]\\
& E &
\end{tikzcd}
\]
commutes.
\end{definition}

\begin{proposition}[Prolongation evaluation as the adjunction counit]
\label{prop:formal-solution-value-map-adjunction-counit}
The prolongation evaluation map satisfies
\[
  c_i\simeq \ell_i^{\mathrm{univ}}\eta^\Sigma_{\operatorname{Prol}(i)},
  \qquad
  ic_i\simeq p_2.
\]
As a morphism in \(\operatorname{Eqn}_{X/S}(Y)\),
\[
  c_i:V_Y\operatorname{Prol}^{\mathrm{Cart}}_Y(i)\longrightarrow i
\]
is the counit of \(V_Y\dashv \operatorname{Prol}^{\mathrm{Cart}}_Y\).
\end{proposition}

\begin{proof}
\leavevmode
\begin{enumerate}[label=\textup{(\arabic*)}]
    \item \textbf{Identify evaluation through adjunction.}
The first identity is Lemma~\ref{lem:jet-mates} applied to the adjunct pair
\(\ell_i^{\mathrm{univ}}\dashv p_1\).

    \item \textbf{Compute the ambient equation map.}
Naturality of the comonad counit and the pullback
relation \(J^\infty(i)p_1\simeq\delta^J_Yp_2\) give
\[
\begin{aligned}
  ic_i
  &\simeq
  i\varepsilon^\Pi_Ep_1\\
  &\simeq
  \varepsilon^\Pi_{J^\infty Y}J^\infty(i)p_1\\
  &\simeq
  \varepsilon^\Pi_{J^\infty Y}\delta^J_Yp_2\\
  &\simeq p_2.
\end{aligned}
\]

    \item \textbf{Recognize the adjunction counit.}
The mapping-space construction in the proof of
Theorem~\ref{thm:cartan-prolongation-adjunction} identifies this morphism with
the adjunction counit.
\end{enumerate}
\end{proof}

\begin{proposition}[Evaluation of formal solutions]
\label{prop:formal-solution-value-evaluation}
Under the natural equivalence
\[
  \Map_X(Z,\operatorname{Prol}(i))\simeq\operatorname{FSol}_i(Z),
\]
the map induced by \(c_i\),
\[
  \Map_X(Z,\operatorname{Prol}(i))\to\Map_X(Z,E),
\]
sends a formal solution represented by
\(\ell:\mathsf{Disk}^{\mathrm{dR},\infty} Z\to E\) to its value
\[
  \ell\eta^\Sigma_Z:Z\to E.
\]
\end{proposition}

\begin{proof}
\leavevmode
\begin{enumerate}[label=\textup{(\arabic*)}]
    \item \textbf{Identify the adjunct of the formal-solution lift.}
If \(f:Z\to \operatorname{Prol}(i)\) represents the formal solution, then \(p_1f\) is the
adjunct of \(\ell\).

    \item \textbf{Compute prolongation evaluation.}
Hence
\[
  c_if
  =
  \varepsilon^\Pi_Ep_1f
  \simeq
  \ell\eta^\Sigma_Z
\]
by Lemma~\ref{lem:jet-mates}.
\end{enumerate}
\end{proof}

\begin{theorem}[Actual solutions as maps into Cartan prolongation]
\label{thm:actual-solutions-prolongation}
For every generalized jet equation
\[
  i:E\longrightarrow J^\infty Y,
\]
there is a canonical natural equivalence
\[
  \operatorname{Sol}(i)
  \xrightarrow{\simeq}
  \Map_{\operatorname{Coalg}_{J^\infty}((\mathcal H_{\Sp})_{/X})}
  \bigl((X,\rho_X),\operatorname{Prol}^{\mathrm{Cart}}(i)\bigr).
\]
Thus the Cartan prolongation represents actual solutions internally in the
jet-coalgebra category. The representing correspondence is organized by
\[
\begin{tikzcd}[column sep=huge,row sep=large]
(X,\rho_X)
  \ar[r]
  \ar[dr]
&
\operatorname{Prol}^{\mathrm{Cart}}(i)
  \ar[d]
\\
&
\mathbb K(Y),
\end{tikzcd}
\]
where the upper map is the coalgebraic form of a solution and the vertical map
is the prolonged ambient equation.
\end{theorem}

\begin{proof}
\leavevmode
\begin{enumerate}[label=\textup{(\arabic*)}]
    \item \textbf{Map the terminal coalgebra into the pullback.}
Mapping from \((X,\rho_X)\) into the defining pullback of \(\operatorname{Prol}^{\mathrm{Cart}}(i)\) gives
\[
\begin{aligned}
  &\Map_{\operatorname{Coalg}_{J^\infty}((\mathcal H_{\Sp})_{/X})}
  \bigl((X,\rho_X),\operatorname{Prol}^{\mathrm{Cart}}(i)\bigr)\\
  &\quad\simeq
  \Map_{\operatorname{Coalg}_{J^\infty}((\mathcal H_{\Sp})_{/X})}
  \bigl((X,\rho_X),\mathbb K(E)\bigr)
  \times_{\Map_{\operatorname{Coalg}_{J^\infty}((\mathcal H_{\Sp})_{/X})}
  ((X,\rho_X),\mathbb K(J^\infty Y))}
  \Map_{\operatorname{Coalg}_{J^\infty}((\mathcal H_{\Sp})_{/X})}
  \bigl((X,\rho_X),\mathbb K(Y)\bigr).
\end{aligned}
\]

    \item \textbf{Apply the cofree adjunction.}
The cofree adjunction identifies the three mapping spaces with
\[
  \Map_X(X,E),
  \qquad
  \Map_X(X,J^\infty Y),
  \qquad
  \Map_X(X,Y).
\]

    \item \textbf{Identify the two comparison maps.}
The first comparison map is induced by \(i\). The second sends
\(\phi:X\to Y\) to
\[
  J^\infty(\phi)\rho_X=j^\infty\phi.
\]

    \item \textbf{Recognize the solution pullback.}
Hence the displayed pullback is exactly the homotopy pullback defining
\(\operatorname{Sol}(i)\).
\end{enumerate}
\end{proof}

\begin{corollary}[Actual solutions induce formal solutions]
\label{prop:actual-to-formal-solutions}
Every point of \(\operatorname{Sol}(i)\) canonically determines a coalgebra map
\[
  (X,\rho_X)\longrightarrow \operatorname{Prol}^{\mathrm{Cart}}(i)
\]
and hence an underlying map
\[
  X\longrightarrow \operatorname{Prol}(i),
\]
equivalently an \(X\)-parameterized formal solution. The construction is
natural in morphisms of generalized jet equations.
\end{corollary}

\begin{proof}
\leavevmode
\begin{enumerate}[label=\textup{(\arabic*)}]
    \item \textbf{Forget the represented coalgebra structure.}
Apply \Ref{thm:actual-solutions-prolongation} and then forget the canonical
coalgebra structure on \(\operatorname{Prol}^{\mathrm{Cart}}(i)\).
\end{enumerate}
\end{proof}

Cartan prolongation has now produced both a universal Cartan equation and a
single comparison back to the original equation. Formal integrability is the
recognition problem for that comparison.

\subsection{Prolongation-completeness and intrinsic formal integrability}
\label{sec:prolongation-completeness}

We call an equation prolongation-complete when the evaluation
\(c_i:\operatorname{Prol}(i)\to E\) is invertible. The results in this
subsection identify that one condition in three useful forms: unique formal
solvability from a value, a Cartesian comonad square for an already-Cartan
equation, and a canonically determined jet-coalgebra structure on an arbitrary
prolongation-complete equation.
\begin{definition}[Prolongation-complete generalized jet equation]
A generalized jet equation
\[
  i:E\to J^\infty Y
\]
is \textbf{prolongation-complete} if
\[
  c_i:\operatorname{Prol}(i)\xrightarrow{\simeq}E.
\]
\end{definition}

\begin{theorem}[Formal-solution recognition of prolongation-completeness]
\label{thm:prolongation-complete-formal-solutions}
The equation \(i\) is prolongation-complete if and only if, for every
\(Z\in(\mathcal H_{\Sp})_{/X}\), evaluation is an equivalence
\[
  \operatorname{FSol}_i(Z)
  \xrightarrow{\simeq}
  \Map_X(Z,E).
\]
Equivalently, every generalized jet equation point carries a coherently unique
formal solution determined by its value. This criterion concerns formal infinitesimal solutions. Analytic existence
of actual solutions belongs to a separate analytic problem.
\end{theorem}

\begin{proof}
\leavevmode
\begin{enumerate}[label=\textup{(\arabic*)}]
    \item \textbf{Identify evaluation with mapping out of the counit.}
By Theorem~\ref{thm:universal-formal-solution} and
Proposition~\ref{prop:formal-solution-value-evaluation}, the displayed map is
\(\Map_X(Z,c_i)\).

    \item \textbf{Apply Yoneda.}
It is an equivalence for every \(Z\) exactly when
\(c_i\) is an equivalence, by Yoneda.
\end{enumerate}
\end{proof}

\begin{proposition}[Split prolongation evaluation for a Cartan equation morphism]
\label{prop:split-formal-solution-value-map}
Let
\[
  i:(E,\rho)\to\mathbb K(Y)
\]
be a jet-coalgebra morphism. There is a canonical coalgebra map
\[
  h_i:(E,\rho)\to \operatorname{Prol}^{\mathrm{Cart}}(i)
\]
whose composites with the two projections are \(\rho\) and \(i\). It satisfies
\[
  c_ih_i\simeq\id_E.
\]
Thus the prolongation evaluation map of every Cartan equation morphism is canonically a split
epimorphism.
\end{proposition}

\begin{proof}
\leavevmode
\begin{enumerate}[label=\textup{(\arabic*)}]
    \item \textbf{Construct the split comparison.}
The coalgebra-morphism identity
\[
  J^\infty(i)\rho\simeq\delta^J_Yi
\]
constructs \(h_i\) by the pullback universal property of \(\operatorname{Prol}^{\mathrm{Cart}}(i)\).

    \item \textbf{Compute the prolongation evaluation.}
Then
\[
  c_ih_i
  \simeq
  \varepsilon^\Pi_E\rho
  \simeq
  \id_E.
\]
\end{enumerate}
\end{proof}

\begin{theorem}[Cartesian criterion]
\label{thm:prolongation-cartesian}
For a Cartan equation morphism
\[
  i:(E,\rho)\to\mathbb K(Y),
\]
the following are equivalent:
\begin{enumerate}
\item the underlying equation \(i:E\to J^\infty Y\) is prolongation-complete;
\item the canonical comparison \(h_i:E\to \operatorname{Prol}(i)\) is an equivalence;
\item the Cartan square
\[
\begin{tikzcd}
E \ar[r,"\rho"] \ar[d,"i"']
  & J^\infty E \ar[d,"J^\infty(i)"]\\
J^\infty Y \ar[r,"\delta^J_Y"']
  & J^\infty J^\infty Y
\end{tikzcd}
\]
is Cartesian.
\end{enumerate}
\end{theorem}

\begin{proof}
\leavevmode
\begin{enumerate}[label=\textup{(\arabic*)}]
    \item \textbf{Identify the prolongation pullback.}
The object \(\operatorname{Prol}(i)\) is the pullback of \(J^\infty(i)\) along
\(\delta^J_Y\), and \(h_i\) is the comparison from \(E\) to that pullback.

    \item \textbf{Compare the two equivalences.}
Since
\(c_ih_i\simeq\id_E\), the map \(c_i\) is an equivalence exactly when
\(h_i\) is.

    \item \textbf{Recognize the Cartesian square.}
The latter is an equivalence exactly when the displayed square is
Cartesian.
\end{enumerate}
\end{proof}

\begin{theorem}[Tautological prolongation-completeness]
\label{thm:tautological-prolongation-completeness}
Let \(\rho:E\to J^\infty E\) be any jet-coalgebra coaction. Regard \(\rho\)
itself as a generalized jet equation on \(Y=E\). Then its prolongation evaluation map is an
equivalence. Hence every intrinsic formally integrable PDE has a canonical
tautological prolongation-complete ambient equation.
\end{theorem}

\begin{proof}
\leavevmode
\begin{enumerate}[label=\textup{(\arabic*)}]
    \item \textbf{Form the tautological prolongation.}
For \(i=\rho\),
\[
  \operatorname{Prol}(\rho)
  =
  J^\infty E
  \times_{J^\infty J^\infty E}
  J^\infty E,
\]
where the two maps are \(J^\infty(\rho)\) and \(\delta^J_E\). Write
\(p_1,p_2:\operatorname{Prol}(\rho)\rightrightarrows J^\infty E\) for the projections, so
\[
  J^\infty(\rho)p_1\simeq\delta^J_Ep_2.
\]

    \item \textbf{Identify the two projections.}
Applying \(J^\infty(\varepsilon^\Pi_E)\) and using
\(\varepsilon^\Pi_E\rho\simeq\id_E\) and
\(J^\infty(\varepsilon^\Pi_E)\delta^J_E\simeq\id\) gives
\[
  p_1\simeq p_2.
\]

    \item \textbf{Apply the second counit identity.}
Applying \(\varepsilon^\Pi_{J^\infty E}\) instead gives
\[
  \rho\varepsilon^\Pi_Ep_1\simeq p_2\simeq p_1.
\]

    \item \textbf{Construct the comparison maps.}
The coalgebra law constructs
\[
  h=(\rho,\rho):E\to \operatorname{Prol}(\rho),
\]
and the prolongation evaluation map is \(c_\rho=\varepsilon^\Pi_Ep_1\).

    \item \textbf{Verify the inverse composites.}
Hence
\[
  c_\rho h\simeq\id_E,
\]
while the two identities above give
\[
  hc_\rho\simeq\id_{\operatorname{Prol}(\rho)}.
\]

    \item \textbf{Conclude the equivalence.}
Thus \(h\) and \(c_\rho\) are inverse equivalences.
\end{enumerate}
\end{proof}

\begin{theorem}[Embedded Cartan equations are prolongation-complete]
If
\[
  i:(E,\rho)\to\mathbb K(Y)
\]
is a Cartan equation morphism whose underlying map
\(E\to J^\infty Y\) is a monomorphism, then \(c_i:\operatorname{Prol}(i)\to E\) is an
equivalence.
\end{theorem}

\begin{proof}
\leavevmode
\begin{enumerate}[label=\textup{(\arabic*)}]
    \item \textbf{Preserve the embedded equation.}
The functor \(J^\infty\) preserves monomorphisms because it preserves limits.
Hence \(J^\infty(i)\) is a monomorphism, and its pullback
\[
  p_2:\operatorname{Prol}(i)\to J^\infty Y
\]
is a monomorphism.

    \item \textbf{Cancel along the monomorphism.}
The split section \(h_i\) satisfies
\[
  p_2h_ic_i
  \simeq
  ic_i
  \simeq
  p_2.
\]
The monomorphism property of \(p_2\) gives \(h_ic_i\simeq\id\), while
\(c_ih_i\simeq\id\) was already constructed.
\end{enumerate}
\end{proof}

The preceding criteria detect when the prolongation evaluation is invertible.
To turn that existence statement into a canonical intrinsic structure, one
further point is needed: compatible Cartan lifts must have contractible moduli.
The next lemma isolates exactly this categorical uniqueness step. Its abstract
notation is used once, in the theorem that follows, and is not part of the
later PDE interface.

\begin{lemma}[Contractible lift over an invertible counit]
\label{lem:contractible-adjunction-lift}
Let
\[
  V:\mathcal A\rightleftarrows\mathcal B:P
\]
be an adjunction \(V\dashv P\), with \(V\) conservative, and let
\[
  \varepsilon^V:VP\to\id_{\mathcal B}
\]
be its counit. If
\[
  \varepsilon^V_b:VPb\xrightarrow{\simeq}b
\]
is an equivalence, then the homotopy fibre
\[
  \operatorname{hofib}_{b}
  \left(
    V^\simeq:\mathcal A^\simeq\to\mathcal B^\simeq
  \right)
\]
is contractible.
\end{lemma}

\begin{proof}
\leavevmode
\begin{enumerate}[label=\textup{(\arabic*)}]
    \item \textbf{Choose the distinguished lift.}
The pair
\[
  (Pb,\varepsilon^V_b)
\]
defines a point of the fibre.

    \item \textbf{Choose an arbitrary lift.}
Let
\[
  (a,\varphi),
  \qquad
  \varphi:Va\xrightarrow{\simeq}b,
\]
be any other point.

    \item \textbf{Apply the adjunction.}
Adjunction identifies the space of maps
\[
  a\longrightarrow Pb
\]
whose adjunct is \(\varphi\) with a contractible homotopy fibre of
\[
  \Map_{\mathcal A}(a,Pb)
  \xrightarrow{\simeq}
  \Map_{\mathcal B}(Va,b).
\]

    \item \textbf{Construct the comparison lift.}
Hence there is a contractibly unique map
\[
  \widetilde\varphi:a\to Pb
\]
such that
\[
  \varepsilon^V_b\,V(\widetilde\varphi)\simeq\varphi.
\]

    \item \textbf{Use conservativity.}
Because both \(\varepsilon^V_b\) and \(\varphi\) are equivalences,
\(V(\widetilde\varphi)\) is an equivalence; conservativity of \(V\) therefore
makes \(\widetilde\varphi\) an equivalence.

    \item \textbf{Compute the automorphism space.}
The same adjunction calculation with
\((a,\varphi)=(Pb,\varepsilon^V_b)\) makes the automorphism space of the
distinguished point contractible.

    \item \textbf{Conclude contractibility.}
Thus every point is connected to it through
a contractible space of equivalences, and the entire homotopy fibre is
contractible.
\end{enumerate}
\end{proof}

\begin{theorem}[Intrinsic structure of a prolongation-complete equation]
\label{thm:prolongation-complete-intrinsic-structure}
If
\[
  i:E\to J^\infty Y
\]
is prolongation-complete, then \(E\) carries a canonical
\(J^\infty\)-coalgebra structure
\[
  \rho_i:E\to J^\infty E
\]
for which
\[
  i:(E,\rho_i)\to\mathbb K(Y)
\]
is a jet-coalgebra morphism.

Writing
\[
  p_1:\operatorname{Prol}(i)\to J^\infty E
\]
for the first projection of the Cartan prolongation, one has canonically
\[
  \rho_i
  \simeq
  p_1c_i^{-1},
\]
where \(c_i^{-1}\) denotes the contractibly unique inverse of
\[
  c_i:\operatorname{Prol}(i)\xrightarrow{\simeq}E.
\]

The space of compatible Cartan structures on the fixed ambient equation \(i\)
is the homotopy fibre
\[
  \mathcal S_i
  :=
  \operatorname{hofib}_{i}
  \left(
    V_Y^\simeq:
    \operatorname{CartEq}_{X/S}(Y)^\simeq
    \longrightarrow
    \operatorname{Eqn}_{X/S}(Y)^\simeq
  \right).
\]
A point consists of a Cartan equation together with a specified equivalence
of its underlying ambient equation with \(i\). The space \(\mathcal S_i\) is contractible.
Consequently \(E\) descends canonically, through effective infinitesimal
descent, to an object of
\[
  \operatorname{PDE}^{\mathrm{fi}}_{X/S}.
\]
\end{theorem}

\begin{proof}
\leavevmode
\begin{enumerate}[label=\textup{(\arabic*)}]
    \item \textbf{Invoke the canonical prolongation coalgebra.}
The Cartan prolongation \(\operatorname{Prol}^{\mathrm{Cart}}(i)\) carries its canonical coalgebra structure by
Theorem~\ref{thm:canonical-prolongation-coalgebra}.

    \item \textbf{Use prolongation-completeness.}
Its underlying object is
\(\operatorname{Prol}(i)\), and prolongation-completeness says that
\[
  c_i:\operatorname{Prol}(i)\xrightarrow{\simeq}E
\]
is an equivalence.

    \item \textbf{Transport the coaction across the evaluation equivalence.}
Let
\[
  \rho_{\operatorname{Prol}(i)}:\operatorname{Prol}(i)\to J^\infty \operatorname{Prol}(i)
\]
be the prolongation coaction. Transporting this coaction across
\(c_i\) constructs
\[
\begin{aligned}
  \rho_i:
  E
  &\xrightarrow{c_i^{-1}}
  \operatorname{Prol}(i)
  \xrightarrow{\rho_{\operatorname{Prol}(i)}}
  J^\infty \operatorname{Prol}(i)
  \xrightarrow{J^\infty(c_i)}
  J^\infty E.
\end{aligned}
\]

    \item \textbf{Record the contractible inverse choice.}
The space of inverses of an equivalence is contractible.

    \item \textbf{Introduce the first projection.}
To compute the transported coaction, use the first projection
\[
  p_1:\operatorname{Prol}(i)\to J^\infty E.
\]

    \item \textbf{Use the coalgebra-morphism identity.}
Since
\[
  p_1:\operatorname{Prol}^{\mathrm{Cart}}(i)\to\mathbb K(E)
\]
is a coalgebra morphism,
\[
  J^\infty(p_1)\rho_{\operatorname{Prol}(i)}
  \simeq
  \delta^J_Ep_1.
\]

    \item \textbf{Apply the comonad counit.}
Using
\[
  c_i=\varepsilon^\Pi_Ep_1
\]
and the comonad counit identity gives
\[
\begin{aligned}
  J^\infty(c_i)\rho_{\operatorname{Prol}(i)}
  &=
  J^\infty(\varepsilon^\Pi_E)
  J^\infty(p_1)\rho_{\operatorname{Prol}(i)}\\
  &\simeq
  J^\infty(\varepsilon^\Pi_E)\delta^J_Ep_1\\
  &\simeq
  p_1.
\end{aligned}
\]

    \item \textbf{Identify the transported coaction.}
After transport across \(c_i\), this is precisely
\[
  \rho_i\simeq p_1c_i^{-1}.
\]

    \item \textbf{Identify the second projection.}
The second projection
\[
  p_2:\operatorname{Prol}^{\mathrm{Cart}}(i)\to\mathbb K(Y)
\]
is also a coalgebra morphism, and
Proposition~\ref{prop:formal-solution-value-map-adjunction-counit} gives
\[
  p_2\simeq ic_i.
\]

    \item \textbf{Transport the ambient equation.}
Transporting \(p_2\) across
\[
  \operatorname{Prol}^{\mathrm{Cart}}(i)\xrightarrow{\simeq}(E,\rho_i)
\]
therefore makes
\[
  i:(E,\rho_i)\to\mathbb K(Y)
\]
a coalgebra morphism.

    \item \textbf{Identify the Cartan prolongation adjunction.}
It remains to prove the asserted contractible uniqueness. By
Theorem~\ref{thm:cartan-prolongation-adjunction},
\[
  V_Y\dashv \operatorname{Prol}^{\mathrm{Cart}}_Y,
\]
and
Proposition~\ref{prop:formal-solution-value-map-adjunction-counit} identifies the counit at
\(i\) with
\[
  c_i:V_Y\operatorname{Prol}^{\mathrm{Cart}}_Y(i)\to i.
\]

    \item \textbf{Establish conservativity.}
The functor \(V_Y\) is conservative: a morphism in \(\operatorname{CartEq}_{X/S}(Y)\)
whose underlying equation morphism is an equivalence has underlying
map over \(X\) an equivalence, and the forgetful functor from
\(J^\infty\)-coalgebras reflects equivalences.

    \item \textbf{Apply the contractible-fibre lemma.}
Since \(c_i\) is an equivalence, apply
Lemma~\ref{lem:contractible-adjunction-lift} to
\[
  V_Y\dashv \operatorname{Prol}^{\mathrm{Cart}}_Y.
\]
It gives
\[
  \mathcal S_i\simeq *.
\]

    \item \textbf{Interpret contractible uniqueness.}
Thus the compatible coalgebra structure is canonical in the precise sense that
its moduli space is contractible.

    \item \textbf{Identify effective jet descent.}
Finally,
Theorem~\ref{thm:effective-jet-descent} identifies
\[
  \operatorname{Coalg}_{J^\infty}((\mathcal H_{\Sp})_{/X})
  \simeq
  (\mathcal H_{\Sp})_{/Q_{X/S}}.
\]

    \item \textbf{Descend the canonical coalgebra.}
Hence the canonical coalgebra \((E,\rho_i)\) determines the stated intrinsic
formally integrable PDE.
\end{enumerate}
\end{proof}

\begin{corollary}[Cartan factorization of a prolongation-complete equation]
Let \(i:E\to J^\infty Y\) be prolongation-complete and put
\[
  y_i:=\varepsilon^\Pi_Yi:E\to Y.
\]
Then
\[
  i\simeq J^\infty(y_i)\rho_i.
\]
\end{corollary}

\begin{proof}
\leavevmode
\begin{enumerate}[label=\textup{(\arabic*)}]
    \item \textbf{Use the intrinsic Cartan structure.}
Theorem~\ref{thm:prolongation-complete-intrinsic-structure} makes \(i\) a coalgebra
morphism into \(\mathbb K(Y)\).

    \item \textbf{Apply Cartan recognition.}
Apply
Theorem~\ref{thm:cartan-recognition}.
\end{enumerate}
\end{proof}

\begin{theorem}[Quotient-side criterion for prolongation-completeness]
\label{thm:quotient-prolongation-criterion}
Let \(F\to Q_{X/S}\) be an intrinsic formally integrable PDE, put
\[
  E=e_{X/S}^*F,
\]
and let
\[
  y:E\to Y
\]
be a map with associated Cartan equation
\[
  i_y:E\to J^\infty Y.
\]
Put
\[
  G:=\Pi_{e_{X/S}}Y
\]
and let
\[
  \bar y:F\to G
\]
be the adjunct of \(y\).

Under effective descent, the Cartan prolongation of \(i_y\) corresponds to
\[
  \Pi_{e_{X/S}}e_{X/S}^*F
  \times_{\Pi_{e_{X/S}}e_{X/S}^*G}
  G.
\]
The canonical Cartan comparison corresponds to the map
\[
  F\longrightarrow
  \Pi_{e_{X/S}}e_{X/S}^*F
  \times_{\Pi_{e_{X/S}}e_{X/S}^*G}
  G
\]
whose two components are
\[
  \eta^\Pi_F:F\to\Pi_{e_{X/S}}e_{X/S}^*F
  \qquad\text{and}\qquad
  \bar y:F\to G.
\]
The compatibility required to define this map is the naturality equivalence
\[
  \Pi_{e_{X/S}}e_{X/S}^*(\bar y)\,\eta^\Pi_F
  \simeq
  \eta^\Pi_G\bar y
\]
for the monad unit
\[
  \eta^\Pi:\id_{(\mathcal H_{\Sp})_{/Q_{X/S}}}
  \longrightarrow
  \Pi_{e_{X/S}}e_{X/S}^*.
\]

Consequently \(i_y\) is prolongation-complete if and only if
\[
\begin{tikzcd}
F \ar[r,"\eta^\Pi_F"] \ar[d,"\bar y"']
  & \Pi_{e_{X/S}}e_{X/S}^*F
      \ar[d,"\Pi_{e_{X/S}}e_{X/S}^*(\bar y)"]\\
G \ar[r,"\eta^\Pi_G"']
  & \Pi_{e_{X/S}}e_{X/S}^*G
\end{tikzcd}
\]
is Cartesian.
\end{theorem}

\begin{proof}
\leavevmode
\begin{enumerate}[label=\textup{(\arabic*)}]
    \item \textbf{Compute the quotient-side adjunct.}
The adjunct of
\[
  y:e_{X/S}^*F\to Y
\]
under
\[
  e_{X/S}^*\dashv\Pi_{e_{X/S}}
\]
is
\[
  \bar y:
  F
  \xrightarrow{\eta^\Pi_F}
  \Pi_{e_{X/S}}e_{X/S}^*F
  \xrightarrow{\Pi_{e_{X/S}}(y)}
  \Pi_{e_{X/S}}Y
  =
  G.
\]

    \item \textbf{Identify the canonical coaction.}
The canonical coaction on
\[
  E=e_{X/S}^*F
\]
is
\[
  \rho_F=e_{X/S}^*\eta^\Pi_F.
\]

    \item \textbf{Identify the Cartan equation as a pullback.}
Hence
\[
\begin{aligned}
  i_y
  &=J^\infty(y)\rho_F\\
  &=e_{X/S}^*\Pi_{e_{X/S}}(y)\,e_{X/S}^*\eta^\Pi_F\\
  &=e_{X/S}^*\bigl(\Pi_{e_{X/S}}(y)\eta^\Pi_F\bigr)\\
  &=e_{X/S}^*\bar y.
\end{aligned}
\]

    \item \textbf{Translate the cofree objects.}
Under
Theorem~\ref{thm:effective-jet-descent}, the cofree coalgebra on
\[
  E=e_{X/S}^*F
\]
corresponds to
\[
  \Pi_{e_{X/S}}e_{X/S}^*F.
\]
Likewise,
\[
  \mathbb K(Y)
\]
corresponds to
\[
  G=\Pi_{e_{X/S}}Y,
\]
and
\[
  \mathbb K(J^\infty Y)
\]
corresponds to
\[
  \Pi_{e_{X/S}}e_{X/S}^*\Pi_{e_{X/S}}Y
  =
  \Pi_{e_{X/S}}e_{X/S}^*G.
\]

    \item \textbf{Translate the first cofree morphism.}
Since
\[
  i_y=e_{X/S}^*\bar y,
\]
the cofree coalgebra morphism
\[
  \mathbb K(i_y):
  \mathbb K(E)\to\mathbb K(J^\infty Y)
\]
corresponds under effective descent to
\[
  \Pi_{e_{X/S}}e_{X/S}^*(\bar y):
  \Pi_{e_{X/S}}e_{X/S}^*F
  \longrightarrow
  \Pi_{e_{X/S}}e_{X/S}^*G.
\]

    \item \textbf{Translate the second cofree morphism.}
The cofree coalgebra map
\[
  \widetilde\delta^J_Y:
  \mathbb K(Y)\to\mathbb K(J^\infty Y)
\]
classified by
\[
  \id_{J^\infty Y}
\]
corresponds to the monad unit
\[
  \eta^\Pi_G:
  G\longrightarrow\Pi_{e_{X/S}}e_{X/S}^*G.
\]

    \item \textbf{Preserve the Cartan pullback.}
The effective-descent equivalence preserves limits. Therefore the Cartan
prolongation
\[
  \operatorname{Prol}^{\mathrm{Cart}}(i_y)
  =
  \mathbb K(E)
  \times_{\mathbb K(J^\infty Y)}
  \mathbb K(Y)
\]
corresponds to
\[
  \Pi_{e_{X/S}}e_{X/S}^*F
  \times_{\Pi_{e_{X/S}}e_{X/S}^*G}
  G.
\]

    \item \textbf{Apply naturality of the monad unit.}
Naturality of the monad unit gives the canonical equivalence
\[
  \Pi_{e_{X/S}}e_{X/S}^*(\bar y)\eta^\Pi_F
  \simeq
  \eta^\Pi_G\bar y.
\]

    \item \textbf{Construct the quotient-side comparison.}
Thus the pair
\[
  (\eta^\Pi_F,\bar y)
\]
defines a canonical comparison map
\[
  F\longrightarrow
  \Pi_{e_{X/S}}e_{X/S}^*F
  \times_{\Pi_{e_{X/S}}e_{X/S}^*G}
  G.
\]

    \item \textbf{Identify the pulled-back comparison.}
After applying the conservative functor
\[
  e_{X/S}^*,
\]
this comparison is exactly the canonical Cartan map
\[
  h_{i_y}:E\longrightarrow \operatorname{Prol}(i_y)
\]
of Proposition~\ref{prop:split-formal-solution-value-map}.

    \item \textbf{Compare the Cartan section and evaluation.}
Since
\[
  c_{i_y}h_{i_y}\simeq\id_E,
\]
the map
\[
  h_{i_y}
\]
is an equivalence if and only if
\[
  c_{i_y}
\]
is an equivalence.

    \item \textbf{Recognize the quotient-side Cartesian criterion.}
On the quotient side, the comparison
\[
  F\longrightarrow
  \Pi_{e_{X/S}}e_{X/S}^*F
  \times_{\Pi_{e_{X/S}}e_{X/S}^*G}
  G
\]
is an equivalence exactly when the square
\[
\begin{tikzcd}
F \ar[r,"\eta^\Pi_F"] \ar[d,"\bar y"']
  & \Pi_{e_{X/S}}e_{X/S}^*F
      \ar[d,"\Pi_{e_{X/S}}e_{X/S}^*(\bar y)"]\\
G \ar[r,"\eta^\Pi_G"']
  & \Pi_{e_{X/S}}e_{X/S}^*G
\end{tikzcd}
\]
is Cartesian.

    \item \textbf{Transport the criterion back to the Cartan equation.}
Conservativity of
\(e_{X/S}^*\) therefore identifies this condition with
prolongation-completeness of \(i_y\).
\end{enumerate}
\end{proof}

\begin{corollary}[Recognition for embedded ambient equations]
Let
\[
  i:E\hookrightarrow J^\infty Y
\]
be an embedded generalized jet equation. The following are equivalent:
\begin{enumerate}
\item \(i\) is prolongation-complete;
\item \(i\) lies in the essential image of
\[
  V_Y:
  \operatorname{CartEq}_{X/S}(Y)\longrightarrow\operatorname{Eqn}_{X/S}(Y);
\]
equivalently, there exist an object \(F\to Q_{X/S}\), a map
\[
  y:e_{X/S}^*F\to Y,
\]
and a specified equivalence in \(\operatorname{Eqn}_{X/S}(Y)\) from the associated
Cartan equation
\[
  J^\infty(y)\rho_F:e_{X/S}^*F\to J^\infty Y
\]
to \(i\);
\item \(E\) carries a jet-coalgebra structure
\[
  \rho:E\to J^\infty E
\]
for which
\[
  i:(E,\rho)\to\mathbb K(Y)
\]
is a coalgebra morphism.
\end{enumerate}
When these conditions hold, the space of compatible Cartan structures on the
fixed ambient equation \(i\) is contractible.
\end{corollary}

\begin{proof}
\leavevmode
\begin{enumerate}[label=\textup{(\arabic*)}]
    \item \textbf{Construct the intrinsic Cartan structure.}
Prolongation-completeness gives \textup{(2)} and \textup{(3)} by
Theorem~\ref{thm:prolongation-complete-intrinsic-structure}.

    \item \textbf{Translate the Cartan structure through descent.}
Theorem~\ref{thm:effective-jet-descent}
identifies a jet coalgebra with an object over \(Q_{X/S}\), while
Theorem~\ref{thm:cartan-recognition} identifies a coalgebra morphism into
\(\mathbb K(Y)\) with its classifying map \(E\to Y\).

    \item \textbf{Identify the equivalent intrinsic presentations.}
This proves
\textup{(2)}\(\Leftrightarrow\)\textup{(3)} at the level of the corresponding
infinity-categories.

    \item \textbf{Recover the converse.}
Conversely, a compatible coalgebra structure makes \(i\) an embedded Cartan
equation morphism, hence prolongation-complete by the embedded Cartan theorem above.

    \item \textbf{Recover contractible uniqueness.}
Contractible uniqueness is the homotopy-fibre statement of
Theorem~\ref{thm:prolongation-complete-intrinsic-structure} and follows from
prolongation-completeness.
\end{enumerate}
\end{proof}

\begin{corollary}[Cartan prolongation is intrinsically formally integrable]
For every generalized jet equation \(i:E\to J^\infty Y\), its Cartan prolongation
\((\operatorname{Prol}(i),\rho_{\operatorname{Prol}(i)})\) descends canonically to an object of
\(\operatorname{PDE}^{\mathrm{fi}}_{X/S}\). If \(i\) is embedded, the prolonged equation
\[
  p_2:\operatorname{Prol}(i)\to J^\infty Y
\]
is embedded and prolongation-complete.
\end{corollary}

\begin{proof}
\leavevmode
\begin{enumerate}[label=\textup{(\arabic*)}]
    \item \textbf{Descend the prolongation coalgebra.}
Every jet coalgebra descends by Theorem~\ref{thm:effective-jet-descent}.

    \item \textbf{Treat the embedded case.}
If
\(i\) is a monomorphism, then \(J^\infty(i)\) is a monomorphism and its pullback \(p_2\) is
a monomorphism. It is therefore an embedded Cartan equation morphism and hence
prolongation-complete.
\end{enumerate}
\end{proof}

Prolongation-completeness is therefore the bridge between the two uses of the
word ``equation'' in the chapter. An ambient equation
\(E\to J^\infty Y\) becomes intrinsic precisely when its prolongation adds no
new underlying object; then its compatible jet-coalgebra structure has
contractible moduli and descends to the effective quotient.

\subsection{Intrinsic solutions and equation calculus}
\label{sec:intrinsic-equation-calculus}

Once an equation is prolongation-complete, its Cartan prolongation may be
transported back to the original source. Actual solutions can then be read
directly as morphisms from the terminal jet coalgebra, while the quotient-side
presentation makes limits and colimits of intrinsic PDEs transparent.
\begin{theorem}[Solutions intrinsically]
\label{thm:solutions-intrinsically}
If \(i:E\to J^\infty Y\) is prolongation-complete and \((E,\rho_i)\) is the
constructed jet-coalgebra presentation, then
\[
  \operatorname{Sol}(i)
  \simeq
  \Map_{\operatorname{Coalg}_{J^\infty}((\mathcal H_{\Sp})_{/X})}
  \bigl((X,\rho_X),(E,\rho_i)\bigr).
\]
Thus actual solutions are precisely morphisms from the terminal jet coalgebra to
the descended equation presentation.
\end{theorem}

\begin{proof}
\leavevmode
\begin{enumerate}[label=\textup{(\arabic*)}]
    \item \textbf{Represent solutions by Cartan prolongation.}
For every equation, \Ref{thm:actual-solutions-prolongation} gives
\[
  \operatorname{Sol}(i)
  \simeq
  \Map_{\operatorname{Coalg}_{J^\infty}((\mathcal H_{\Sp})_{/X})}
  \bigl((X,\rho_X),\operatorname{Prol}^{\mathrm{Cart}}(i)\bigr).
\]

    \item \textbf{Transport across prolongation evaluation.}
If \(i\) is prolongation-complete,
\Ref{thm:prolongation-complete-intrinsic-structure} transports the prolongation
coalgebra across the equivalence
\[
  c_i:\operatorname{Prol}(i)\xrightarrow{\simeq}E.
\]
Postcomposition with the resulting coalgebra equivalence gives the displayed
mapping-space equivalence.
\end{enumerate}
\end{proof}

\begin{theorem}[The PDE infinity-topos]
\label{thm:PDE-topos-limits}
The category
\[
  \operatorname{PDE}^{\mathrm{fi}}_{X/S}
  =
  (\mathcal H_{\Sp})_{/Q_{X/S}}
\]
is an infinity-topos and admits all small limits and colimits. Under either
the monadic or comonadic presentation, the underlying-object functor to
\((\mathcal H_{\Sp})_{/X}\) preserves all small limits and colimits. Consequently the
underlying object over \(X\) of an intrinsic limit or colimit is canonically the
corresponding limit or colimit of the underlying objects over \(X\).
\end{theorem}

\begin{proof}
\leavevmode
\begin{enumerate}[label=\textup{(\arabic*)}]
    \item \textbf{Use the slice-topos theorem.}
A slice of an infinity-topos is again an infinity-topos
\cite{LurieHTT}.

    \item \textbf{Identify the underlying-object functor.}
Under
Theorem~\ref{thm:effective-jet-descent}, the underlying-object functor is
identified with
\[
  e_{X/S}^*:(\mathcal H_{\Sp})_{/Q_{X/S}}\to(\mathcal H_{\Sp})_{/X}.
\]

    \item \textbf{Use the dependent adjoints.}
Since
\[
  \Sigma_{e_{X/S}}\dashv e_{X/S}^*\dashv\Pi_{e_{X/S}},
\]
this pullback functor is both a left and a right adjoint and hence preserves all
small colimits and limits.

    \item \textbf{Transport preservation through the presentations.}
Transport through the monadic or comonadic equivalence proves the statement.
\end{enumerate}
\end{proof}

\begin{corollary}[Simultaneous formally integrable systems]
Products, fibre products, homotopy equalizers, arbitrary simultaneous systems,
coproducts, and all other small colimits of intrinsic formally integrable PDEs
are again intrinsic formally integrable PDEs and are computed on their
underlying objects over \(X\).
\end{corollary}

\begin{corollary}[Prolonged homotopy equalizers]
Let
\[
  D_0,D_1:E\rightsquigarrow F
\]
be differential operators represented by
\[
  \widehat D_0,\widehat D_1:J^\infty E\rightrightarrows F,
\]
and write
\[
  i_0:
  E_0:=\operatorname{Eq}(\widehat D_0,\widehat D_1)
  \longrightarrow
  J^\infty E
\]
for the derived equation. Its Cartan prolongation satisfies
\[
  \operatorname{Prol}(i_0)
  \simeq
  \operatorname{Eq}
  \bigl(\operatorname{prol}_{\mathrm{op}}^\infty(\widehat D_0),\operatorname{prol}_{\mathrm{op}}^\infty(\widehat D_1)\bigr),
\]
where the equalizer on the right is formed in \((\mathcal H_{\Sp})_{/X}\) over
\(J^\infty E\).
\end{corollary}

\begin{proof}
\leavevmode
\begin{enumerate}[label=\textup{(\arabic*)}]
    \item \textbf{Apply the jet functor to the equalizer.}
Since \(J^\infty\) preserves limits,
\[
  J^\infty E_0
  \simeq
  \operatorname{Eq}
  \bigl(J^\infty(\widehat D_0),J^\infty(\widehat D_1)\bigr)
\]
inside \(J^\infty J^\infty E\).

    \item \textbf{Pull back along comultiplication.}
Pulling this equalizer back along
\(\delta^J_E:J^\infty E\to J^\infty J^\infty E\) gives the equalizer of
\[
  J^\infty(\widehat D_0)\delta^J_E,
  \qquad
  J^\infty(\widehat D_1)\delta^J_E,
\]
which are the two prolonged operators.
\end{enumerate}
\end{proof}

\begin{proposition}[Intrinsic morphisms preserve Cartan structure]
Let \(F,G\to Q_{X/S}\) be intrinsic formally integrable PDEs, with underlying objects \(E,E'\), and let
\(f:F\to G\) be a morphism over \(Q_{X/S}\). Then
\[
  e_{X/S}^*f:E\to E'
\]
is simultaneously a \(\mathsf{Disk}^{\mathrm{dR},\infty}\)-algebra morphism and a
\(J^\infty\)-coalgebra morphism. Consequently it preserves formal-disk
actions, crystalline transport, and every formal-solution family obtained by
precomposition.
\end{proposition}

\begin{proof}
\leavevmode
\begin{enumerate}[label=\textup{(\arabic*)}]
    \item \textbf{Apply the mapping-space descent equivalence.}
The mapping-space statement of Theorem~\ref{thm:effective-jet-descent} gives the
claim.
\end{enumerate}
\end{proof}

The equation calculus is now internal: every generalized equation admits a
Cartan prolongation, and a prolongation-complete equation already carries the
resulting intrinsic descent structure. The remaining task is geometric
functoriality.

\section{Substitution, locality, and geometric models}
\label{sec:jets-locality-models}

So far the calculus has been developed over one relative context. Geometry
requires two further checks. Arbitrary base change must transport the
construction between bases, and effective formally \(\acute{e}\)tale descent
must reconstruct it from local contexts. After these functoriality statements
are fixed, group objects and affine-space charts serve as concrete translation
models of the same intrinsic relation.

\subsection{Base change of prolongations and solution objects}
\label{sec:jet-base-change-consequences}

The substitution theorem already transports the jet functors. We now apply it
to the finite-limit and dependent-product constructions defining
prolongations and solution objects.
\begin{corollary}[Base change of Cartan prolongations]
Let
\[
  i:E\to J^\infty_{X/S}Y
\]
be a generalized jet equation and let
\[
  i':E':=g_X^*E
  \longrightarrow
  J^\infty_{X'/S'}(g_X^*Y)
\]
be its pullback. Then there are canonical equivalences
\[
  g_X^*(\operatorname{Prol}(i))\simeq\operatorname{Prol}(i'),
  \qquad
  g_X^*(c_i)\simeq c_{i'}.
\]
They identify the canonical prolongation coalgebras, both projections, and the
universal formal-solution family. Consequently prolongation-completeness is
preserved by arbitrary base change.
\end{corollary}

\begin{proof}
\leavevmode
\begin{enumerate}[label=\textup{(\arabic*)}]
    \item \textbf{Express Cartan prolongation as a pullback.}
Cartan prolongation is the finite pullback
\[
  \operatorname{Prol}(i)
  =
  J^\infty E
  \times_{J^\infty J^\infty Y}
  J^\infty Y.
\]

    \item \textbf{Base-change the prolongation pullback.}
Pullback along \(g_X\) preserves this limit, and
Theorem~\ref{thm:jet-base-change} identifies every jet term and every
comultiplication map with the corresponding object after base change.

    \item \textbf{Base-change the evaluation map.}
The
prolongation evaluation map is \(\varepsilon^\Pi_Ep_1\), so compatibility with the comonad counit
gives the second equivalence.
\end{enumerate}
\end{proof}

\begin{corollary}[Base change of solution objects]
For a generalized jet equation \(i\), the parameterized solution object satisfies
\[
  \mathbf{Sol}_{X'/S'}(i')
  \simeq
  \mathbf{Sol}_{X/S}(i)\times_SS'.
\]
The same base-change comparison holds for the formal-solution representing
object \(\operatorname{Prol}(i)\).
\end{corollary}

\begin{proof}
\leavevmode
\begin{enumerate}[label=\textup{(\arabic*)}]
    \item \textbf{Base-change the parameterized solution pullback.}
The solution object is a finite pullback of dependent products and the
prolongation map.

    \item \textbf{Apply Beck--Chevalley.}
Apply right Beck--Chevalley and
Theorem~\ref{thm:jet-base-change}.

    \item \textbf{Base-change the formal-solution object.}
The formal-solution statement is the
preceding corollary.
\end{enumerate}
\end{proof}

Cartan prolongation, its evaluation map, and both solution constructions
therefore inherit the substitution law of the jet comonad. Locality asks for
the corresponding reconstruction statement from a cover.

\subsection{Formally \texorpdfstring{$\acute{e}$}{e}tale invariance and effective descent}
\label{sec:jet-formal-etale}

Formally \(\acute{e}\)tale maps are the relevant local changes of context
because the preceding de Rham chapter proves that they preserve the effective
relation itself. The results below first restrict the calculus along one such
map and then reconstruct it from an effective formally \(\acute{e}\)tale
\v{C}ech nerve.
\begin{theorem}[Formally \(\acute{e}\)tale restriction]
\label{thm:formal-etale-jets}
Let \(u:U\to X\) be formally \(\acute{e}\)tale over \(S\). The Cartesian
effective-quotient square of
Theorem~\ref{thm:formal-etale-neighbourhood-invariance} in
Chapter~\ref{chap:spectral-de-Rham} induces canonical equivalences
\[
  u^*\mathsf{Disk}^{\mathrm{dR},\infty}_{X/S}
  \simeq
  \mathsf{Disk}^{\mathrm{dR},\infty}_{U/S}u^*,
  \qquad
  u^*J^\infty_{X/S}
  \simeq
  J^\infty_{U/S}u^*,
\]
compatibly with the monad, comonad, formal-disk--jet adjunction, disk
coalgebras, and mate correspondence.
\end{theorem}

\begin{proof}
\leavevmode
\begin{enumerate}[label=\textup{(\arabic*)}]
    \item \textbf{Apply Beck--Chevalley to the quotient square.}
Apply left and right Beck--Chevalley to the Cartesian effective-quotient square
of Theorem~\ref{thm:formal-etale-neighbourhood-invariance}.

    \item \textbf{Transport the structural maps by mates.}
Compatibility with
the structural maps is the mate coherence of the base-changed adjoint triple.
\end{enumerate}
\end{proof}

\begin{theorem}[Effective formally \(\acute{e}\)tale descent of the jet calculus]
\label{thm:formally-etale-descent-jets}
Let
\[
  u:U\twoheadrightarrow X
\]
be an effective epimorphism over \(S\) which is formally
\(\acute{e}\)tale over \(S\), and let \(U_\bullet\to X\) be its
\v{C}ech nerve. Write
\[
  u_n:U_n\to X
\]
for the structure maps.

Then restriction induces a canonical equivalence
\[
  \operatorname{PDE}^{\mathrm{fi}}_{X/S}
  \xrightarrow{\simeq}
  \operatorname*{Tot}_{[n]\in\Delta}
  \operatorname{PDE}^{\mathrm{fi}}_{U_n/S}.
\]
Moreover, for every \(E\to X\), there are canonical reconstruction
equivalences
\[
  \mathsf{Disk}^{\mathrm{dR},\infty}_{X/S}E
  \simeq
  \operatorname*{Tot}_{[n]\in\Delta}
  \Pi_{u_n}
  \mathsf{Disk}^{\mathrm{dR},\infty}_{U_n/S}(u_n^*E),
\]
and
\[
  J^\infty_{X/S}E
  \simeq
  \operatorname*{Tot}_{[n]\in\Delta}
  \Pi_{u_n}
  J^\infty_{U_n/S}(u_n^*E).
\]
These equivalences reconstruct the monad and comonad structure, the
formal-disk--jet adjunction, the canonical disk coalgebras, and the mate
comparison from their restrictions to the \v{C}ech stages.
\end{theorem}

\begin{proof}
\leavevmode
\begin{enumerate}[label=\textup{(\arabic*)}]
    \item \textbf{Identify the quotient Čech nerve.}
By the effective formally \(\acute{e}\)tale quotient-descent theorem of
Chapter~\ref{chap:spectral-de-Rham}, the simplicial effective quotient
\(Q_{U_\bullet/S}\) is the \v{C}ech nerve of the effective morphism
\[
  Q_{U/S}\longrightarrow Q_{X/S}.
\]

    \item \textbf{Apply effective descent to the quotient slices.}
Effective descent for slices therefore gives
\[
  (\mathcal H_{\Sp})_{/Q_{X/S}}
  \simeq
  \operatorname*{Tot}_{[n]\in\Delta}
  (\mathcal H_{\Sp})_{/Q_{U_n/S}}.
\]

    \item \textbf{Identify the intrinsic PDE totalization.}
Using the definition
\[
  \operatorname{PDE}^{\mathrm{fi}}_{-/S}
  =
  (\mathcal H_{\Sp})_{/Q_{-/S}}
\]
gives the first assertion.

    \item \textbf{Restrict the jet calculus.}
For every \(n\),
Theorem~\ref{thm:formal-etale-jets} gives coherent equivalences
\[
  u_n^*
  \mathsf{Disk}^{\mathrm{dR},\infty}_{X/S}
  \simeq
  \mathsf{Disk}^{\mathrm{dR},\infty}_{U_n/S}u_n^*,
\]
and
\[
  u_n^*J^\infty_{X/S}
  \simeq
  J^\infty_{U_n/S}u_n^*,
\]
compatible with the corresponding adjunction and structural maps.

    \item \textbf{State the slice reconstruction formula.}
Effective \v{C}ech descent for the cover \(u\) gives, for every
\(K\in(\mathcal H_{\Sp})_{/X}\),
\[
  K
  \simeq
  \operatorname*{Tot}_{[n]\in\Delta}
  \Pi_{u_n}u_n^*K.
\]

    \item \textbf{Verify the reconstruction formula on mapping spaces.}
Indeed, for every \(W\to X\),
\[
\begin{aligned}
  \Map_X
  \left(
    W,
    \operatorname*{Tot}_{[n]\in\Delta}
    \Pi_{u_n}u_n^*K
  \right)
  &\simeq
  \operatorname*{Tot}_{[n]\in\Delta}
  \Map_{U_n}(u_n^*W,u_n^*K)\\
  &\simeq
  \Map_X(W,K),
\end{aligned}
\]
and Yoneda gives the reconstruction equivalence.

    \item \textbf{Apply reconstruction to the disk and jet objects.}
Apply this formula to
\[
  K=\mathsf{Disk}^{\mathrm{dR},\infty}_{X/S}E
\]
and to
\[
  K=J^\infty_{X/S}E,
\]
then use the formally \(\acute{e}\)tale restriction equivalences from the
preceding step. This gives the two displayed formulas.

    \item \textbf{Reconstruct the structural maps.}
The units, counits, multiplication, comultiplication, adjunction maps, disk
coalgebras, and mate transformations are natural transformations between the
functors just reconstructed. Effective descent is fully faithful on mapping
spaces, so each such transformation and its higher compatibility are uniquely
reconstructed from their degreewise restrictions.
\end{enumerate}
\end{proof}

\begin{corollary}[Descent of equations and Cartan prolongation]
\label{cor:formally-etale-equation-descent}
Under the hypotheses of
Theorem~\ref{thm:formally-etale-descent-jets}, restriction induces a canonical
equivalence
\[
  \operatorname{Eqn}_{X/S}
  \xrightarrow{\simeq}
  \operatorname*{Tot}_{[n]\in\Delta}
  \operatorname{Eqn}_{U_n/S}.
\]
For fixed \(Y\in(\mathcal H_{\Sp})_{/X}\), this restricts to
\[
  ((\mathcal H_{\Sp})_{/X})_{/J^\infty_{X/S}Y}
  \xrightarrow{\simeq}
  \operatorname*{Tot}_{[n]\in\Delta}
  ((\mathcal H_{\Sp})_{/U_n})_{
    /J^\infty_{U_n/S}(u_n^*Y)}.
\]

If
\[
  i:E\to J^\infty_{X/S}Y
\]
is a generalized jet equation and \(i_n\) its restriction to \(U_n\), then
\[
  u_n^*\operatorname{Prol}(i)
  \simeq
  \operatorname{Prol}(i_n)
\]
for every \(n\), compatibly with the canonical prolongation coalgebras and
evaluation maps. Consequently
\[
  \operatorname{Prol}(i)
  \simeq
  \operatorname*{Tot}_{[n]\in\Delta}
  \Pi_{u_n}\operatorname{Prol}(i_n).
\]
\end{corollary}

\begin{proof}
\leavevmode
\begin{enumerate}[label=\textup{(\arabic*)}]
    \item \textbf{Express the equation category as a finite limit.}
Construction~\ref{con:equation-category} expresses
\(\operatorname{Eqn}_{X/S}\) as a finite limit in \(\operatorname{Cat}_\infty\)
formed from the slice \((\mathcal H_{\Sp})_{/X}\) and the endofunctor
\(J^\infty_{X/S}\).

    \item \textbf{Reconstruct the finite-limit ingredients.}
Effective descent reconstructs the slice, while
Theorem~\ref{thm:formally-etale-descent-jets} reconstructs the jet functor and
its coherent restriction maps.

    \item \textbf{Descend the equation category.}
Since
\(\operatorname{Fun}(\Delta^1,-)\), products, fibres, and pullbacks preserve
limits of infinity-categories, the equation category descends as asserted.

    \item \textbf{Express Cartan prolongation as a pullback.}
Cartan prolongation is the finite pullback
\[
  \operatorname{Prol}(i)
  =
  J^\infty E
  \times_{J^\infty J^\infty Y}
  J^\infty Y.
\]

    \item \textbf{Restrict Cartan prolongation.}
Restriction preserves this pullback and identifies every jet term and the
comultiplication with its corresponding restriction. Hence
\[
  u_n^*\operatorname{Prol}(i)
  \simeq
  \operatorname{Prol}(i_n).
\]

    \item \textbf{Reconstruct the global prolongation.}
The final formula follows from effective reconstruction in the slice:
\[
  K\simeq
  \operatorname*{Tot}_{[n]\in\Delta}\Pi_{u_n}u_n^*K
\]
applied to \(K=\operatorname{Prol}(i)\).
\end{enumerate}
\end{proof}

\begin{corollary}[Descent of solution objects]
\label{cor:formally-etale-solution-descent}
Under the same hypotheses, let
\[
  i:E\to J^\infty_{X/S}Y
\]
and let \(i_n\) denote its restriction to \(U_n\). Then the parameterized
solution object satisfies
\[
  \mathbf{Sol}_{X/S}(i)
  \simeq
  \operatorname*{Tot}_{[n]\in\Delta}
  \mathbf{Sol}_{U_n/S}(i_n)
\]
in \((\mathcal H_{\Sp})_{/S}\). The formal-solution functor is reconstructed
from the corresponding degreewise Cartan prolongations through
\[
  \operatorname{Prol}(i)
  \simeq
  \operatorname*{Tot}_{[n]\in\Delta}
  \Pi_{u_n}\operatorname{Prol}(i_n).
\]
\end{corollary}

\begin{proof}
\leavevmode
\begin{enumerate}[label=\textup{(\arabic*)}]
    \item \textbf{Define the Čech-stage structural maps.}
Put
\[
  p_n:=p\,u_n:U_n\to S.
\]

    \item \textbf{Reconstruct dependent products.}
For every \(Z\to X\), effective reconstruction in the slice and preservation
of limits by \(\Pi_p\) give
\[
\begin{aligned}
  \Pi_pZ
  &\simeq
  \Pi_p
  \operatorname*{Tot}_{[n]\in\Delta}
  \Pi_{u_n}u_n^*Z\\
  &\simeq
  \operatorname*{Tot}_{[n]\in\Delta}
  \Pi_{p_n}u_n^*Z.
\end{aligned}
\]

    \item \textbf{Apply the reconstruction formula to the solution data.}
Apply this to
\[
  Z=E,
  \qquad
  Z=Y,
  \qquad
  Z=J^\infty Y.
\]

    \item \textbf{Compare the parameterized prolongation maps.}
The formally \(\acute{e}\)tale jet comparison identifies the restrictions of
the parameterized prolongation map
\[
  j^\infty_{/S}:\Pi_pY\to\Pi_pJ^\infty Y
\]
with the corresponding maps on the \v{C}ech stages.

    \item \textbf{Reconstruct the solution pullback.}
Since the solution object
is the finite pullback
\[
  \mathbf{Sol}_{X/S}(i)
  =
  \Pi_pE
  \times_{\Pi_pJ^\infty Y}
  \Pi_pY,
\]
the asserted totalization formula follows.

    \item \textbf{Reconstruct the formal-solution object.}
The formal-solution statement is
Corollary~\ref{cor:formally-etale-equation-descent}.
\end{enumerate}
\end{proof}

\begin{corollary}[Locality of prolongation-completeness]
\label{cor:formally-etale-prolongation-local}
Under the hypotheses of
Theorem~\ref{thm:formally-etale-descent-jets}, a generalized jet equation
\[
  i:E\to J^\infty_{X/S}Y
\]
is prolongation-complete if and only if its restriction
\[
  u^*i
\]
to \(U\) is prolongation-complete.
\end{corollary}

\begin{proof}
\leavevmode
\begin{enumerate}[label=\textup{(\arabic*)}]
    \item \textbf{Identify the restricted evaluation map.}
By Corollary~\ref{cor:formally-etale-equation-descent}, restriction identifies
the prolongation evaluation map
\[
  c_i:\operatorname{Prol}(i)\to E
\]
with
\[
  c_{u^*i}:\operatorname{Prol}(u^*i)\to u^*E.
\]

    \item \textbf{Use conservative pullback.}
Pullback along the effective epimorphism \(u\) is conservative. Hence
\(c_i\) is an equivalence if and only if \(c_{u^*i}\) is an equivalence.
\end{enumerate}
\end{proof}

\begin{remark}[Later Nisnevich specialization]
If a later geometric theorem identifies the relevant spectral Nisnevich covers
with effective formally \(\acute{e}\)tale covers, the preceding theorem and its
corollaries yield Nisnevich descent for jets, equations, solution objects,
Cartan prolongations, prolongation-completeness, and the intrinsic PDE
infinity-topos. Nisnevich geometry therefore enters as a later localization theorem for this
calculus.
\end{remark}

Effective formally \(\acute{e}\)tale descent now reconstructs the jet
calculus and detects prolongation-completeness locally. These results establish
locality without choosing coordinates. Coordinates enter only as models of the
already-constructed relation.

\subsection{Formal translation models}
\label{sec:formal-translation}

When the relative context is a commutative group object, subtraction turns the
abstract relation into a familiar action groupoid: two related points differ
by a unique infinitesimal element in the formal neighbourhood of the identity.
This translation description is a checksum for the intrinsic construction,
not a replacement for it.
\begin{proposition}[Formal translation in degree one]
Let \(G\to S\) be a commutative group object in the slice and let
\[
  0:S\to G
\]
be its zero section. The zero section determines the relative de Rham neighbourhood
\[
  \operatorname{Nbh}^{\mathrm{dR}}_{S/S}(G)
  \simeq
  G\times_{G_{\mathrm{dR}/S}}S,
\]
where \(S\to G_{\mathrm{dR}/S}\) is the de Rham image of the zero section.
The zero section also induces \(S\to Q_{G/S}\), and there is a canonical
equivalence
\[
  \operatorname{Nbh}^{\mathrm{dR}}_{S/S}(G)
  \simeq
  G\times_{Q_{G/S}}S.
\]
Moreover \(\operatorname{Nbh}^{\mathrm{dR}}_{S/S}(G)\) carries the canonical commutative group-object structure over
\(S\) induced as the kernel of the group morphism
\[
  \eta_{G/S}^{\mathrm{dR}}:G\to G_{\mathrm{dR}/S}.
\]
Thus this relative de Rham neighbourhood uses the same effective de Rham quotient
as the intrinsic jet calculus. There is a canonical equivalence
\[
  G\times_{G_{\mathrm{dR}/S}}G
  \xrightarrow{\simeq}
  G\times_S\operatorname{Nbh}^{\mathrm{dR}}_{S/S}(G)
\]
which in internal additive notation sends
\[
  (g_0,g_1)\longmapsto(g_0,g_1-g_0)
\]
and whose inverse sends \((g,\xi)\mapsto(g,g+\xi)\). The source projection becomes
\(\operatorname{pr}_G\) and evaluation becomes the formal translation action
\[
  \alpha_G:G\times_S\operatorname{Nbh}^{\mathrm{dR}}_{S/S}(G)\to G,
  \qquad
  (g,\xi)\mapsto g+\xi.
\]
The degree-one relation therefore becomes the action-groupoid square
\[
\begin{tikzcd}[column sep=huge,row sep=large]
G\times_S\operatorname{Nbh}^{\mathrm{dR}}_{S/S}(G)
  \ar[r,"\alpha_G"]
  \ar[d,"\operatorname{pr}_G"']
&
G \ar[d,"e_{G/S}"]\\
G \ar[r,"e_{G/S}"']
&
Q_{G/S},
\end{tikzcd}
\]
and hence
\[
  \mathsf{Disk}^{\mathrm{dR},\infty}_{G/S}
  \simeq
  \Sigma_{\operatorname{pr}_G}\alpha_G^*,
  \qquad
  J^\infty_{G/S}
  \simeq
  \Pi_{\operatorname{pr}_G}\alpha_G^*.
\]
\end{proposition}

\begin{proof}
\leavevmode
\begin{enumerate}[label=\textup{(\arabic*)}]
    \item \textbf{Preserve commutative group objects under reflection.}
Relative de Rham reflection is left exact, hence preserves commutative group
objects, and \(\eta_{G/S}^{\mathrm{dR}}:G\to G_{\mathrm{dR}/S}\) is a group morphism.

    \item \textbf{Identify the relative neighbourhood as the kernel.}
Thus
\(\operatorname{Nbh}^{\mathrm{dR}}_{S/S}(G)\) is its kernel.

    \item \textbf{Induce the kernel group structure.}
Kernels of morphisms of commutative group objects are
finite-limit constructions, so the pullback defining \(\operatorname{Nbh}^{\mathrm{dR}}_{S/S}(G)\) inherits its commutative group structure from \(G\). In particular the expressions \(\xi_1+\xi_2\) and \(-\xi\) used below
are operations of this constructed commutative group object.

    \item \textbf{Pass from the de Rham object to the effective quotient.}
Write
\[
  G\xrightarrow{e_{G/S}}Q_{G/S}
  \xrightarrow{m_{G/S}}G_{\mathrm{dR}/S}
\]
for the effective-image factorization. Both the map
\(G\to G_{\mathrm{dR}/S}\) and the de Rham image of the zero section factor
through \(Q_{G/S}\). Since \(m_{G/S}\) is a monomorphism,
\[
  Q_{G/S}\times_{G_{\mathrm{dR}/S}}Q_{G/S}
  \simeq
  Q_{G/S},
\]
and finite-limit pasting gives
\[
  G\times_{G_{\mathrm{dR}/S}}S
  \simeq
  G\times_{Q_{G/S}}S.
\]

    \item \textbf{Construct the translation equivalence.}
Equality of the de Rham images of \(g_0\) and \(g_1\)
is equivalent to vanishing of the de Rham image of \(g_1-g_0\). Addition and
subtraction in the group object therefore construct the two displayed maps,
and the group laws show that they are inverse equivalences.

    \item \textbf{Transport the jet formulas.}
The jet
formulas are Proposition~\ref{prop:relation-jet-formulas} transported through
this equivalence.
\end{enumerate}
\end{proof}

\begin{theorem}[Formal-translation nerve with explicit simplicial operators]
\label{thm:formal-translation-nerve-explicit}
For every \(n\geq0\), there is a canonical equivalence
\[
  \Phi_{G,n}:
  \mathcal R^{\mathrm{dR}}_{G/S,n}\xrightarrow{\simeq}G\times_S(\operatorname{Nbh}^{\mathrm{dR}}_{S/S}(G))^{\times_S n}
\]
given internally by
\[
  \Phi_{G,n}(g_0,\ldots,g_n)
  =
  (g_0,\xi_1,\ldots,\xi_n),
  \qquad
  \xi_j=g_j-g_{j-1},
\]
with inverse
\[
  (g,\xi_1,\ldots,\xi_n)
  \longmapsto
  (g,g+\xi_1,g+\xi_1+\xi_2,\ldots,g+\xi_1+\cdots+\xi_n).
\]
Under these coordinates the face maps are
\[
\begin{aligned}
  d_0(g,\xi_1,\ldots,\xi_n)
  &=(g+\xi_1,\xi_2,\ldots,\xi_n),\\
  d_i(g,\xi_1,\ldots,\xi_n)
  &=(g,\xi_1,\ldots,\xi_i+\xi_{i+1},\ldots,\xi_n),
    \quad 0<i<n,\\
  d_n(g,\xi_1,\ldots,\xi_n)
  &=(g,\xi_1,\ldots,\xi_{n-1}),
\end{aligned}
\]
and degeneracies insert the zero element of
\(\operatorname{Nbh}^{\mathrm{dR}}_{S/S}(G)\). In low degrees the nerve is
\[
\begin{tikzcd}[column sep=huge]
G\times_S K\times_S K
  \ar[r,shift left=2.2ex,"d_0"]
  \ar[r,"d_1" description]
  \ar[r,shift right=2.2ex,"d_2"']
&
G\times_S K
  \ar[r,shift left=1.2ex,"\operatorname{pr}_G"]
  \ar[r,shift right=1.2ex,"\alpha_G"']
&
G ,
\end{tikzcd}
\]
where \(K=\operatorname{Nbh}^{\mathrm{dR}}_{S/S}(G)\). Hence the entire
relation groupoid is the action-groupoid nerve of formal translation.
\end{theorem}

\begin{proof}
\leavevmode
\begin{enumerate}[label=\textup{(\arabic*)}]
    \item \textbf{Place successive differences in the kernel.}
Every vertex of a relation simplex has the same image in
\(G_{\mathrm{dR}/S}\), so
\[
  \eta_{G/S}^{\mathrm{dR}}(\xi_j)
  =
  \eta_{G/S}^{\mathrm{dR}}(g_j)-\eta_{G/S}^{\mathrm{dR}}(g_{j-1})
  =0.
\]

    \item \textbf{Factor the successive differences through the kernel.}
Thus every difference factors canonically through \(\operatorname{Nbh}^{\mathrm{dR}}_{S/S}(G)\).

    \item \textbf{Construct the inverse coordinates.}
Successive partial
sums construct the inverse.

    \item \textbf{Compute the simplicial operators.}
Deleting the first vertex translates the source point by
\(\xi_1\); deleting an interior vertex adds the adjacent differences; deleting
the last vertex forgets the final difference; repeating a vertex inserts zero.
These are precisely the displayed simplicial operators.
\end{enumerate}
\end{proof}

\begin{remark}[Comparison with synthetic formal disks]
The translation formula is the intrinsic spectral analogue of the formal
translation models used in synthetic differential geometry
\cite{Kock1980,KhavkineSchreiber2017}. The infinitesimal translation object is
\(\operatorname{Nbh}^{\mathrm{dR}}_{S/S}(G)\), the fibre of the
square-zero-generated relative de Rham unit at the identity.
\end{remark}

\subsubsection{Relative linear spaces as translation models}
\label{subsec:relative-linear-translation}

\begin{proposition}[Canonical additive group structure on a relative linear space]
\label{prop:linear-space-additive-group}
Let \(\mathcal E\) be a finite locally free quasi-coherent module on the represented
spectral base \(S\). The relative linear space \(\mathbb L_S(\mathcal E)\) carries a canonical
commutative group-object structure over \(S\).

On the coordinate algebra
\[
  \operatorname{Sym}_{\mathcal O_S}(\mathcal E),
\]
the corresponding Hopf-algebra structure is induced by the maps of
quasi-coherent modules
\[
  \mathcal E\xrightarrow{\Delta_{\mathcal E}}\mathcal E\oplus \mathcal E,
  \qquad
  v\longmapsto(v,v),
\]
\[
  \mathcal E\longrightarrow0,
\]
and
\[
  -\id_{\mathcal E}:\mathcal E\longrightarrow \mathcal E.
\]
Under the canonical equivalence
\[
  \operatorname{Sym}(\mathcal E\oplus \mathcal E)
  \simeq
  \operatorname{Sym}(\mathcal E)\otimes_{\mathcal O_S}\operatorname{Sym}(\mathcal E),
\]
the resulting comultiplication is characterized on generators by
\[
  v\longmapsto v\otimes1+1\otimes v.
\]
The counit sends every generator to zero and the antipode sends \(v\) to
\(-v\).

This group-object construction is contravariantly functorial in \(\mathcal E\): a
morphism
\[
  \varphi:\mathcal E\to \mathcal E'
\]
induces a canonical homomorphism of commutative group objects
\[
  \mathbb L_S(\mathcal E')
  \longrightarrow
  \mathbb L_S(\mathcal E).
\]
It is also compatible with arbitrary base change on \(S\).
\end{proposition}

\begin{proof}
\leavevmode
\begin{enumerate}[label=\textup{(\arabic*)}]
    \item \textbf{Identify the symmetric-monoidal free-algebra comparison.}
The free commutative-algebra functor
\[
  \operatorname{Sym}_{\mathcal O_S}:
  \QCoh(S)\longrightarrow\CAlg(\QCoh(S))
\]
is symmetric monoidal with respect to direct sum in the source and tensor
product in the target \cite{LurieHigherAlgebra}, giving
\[
  \operatorname{Sym}(\mathcal E\oplus \mathcal E)
  \simeq
  \operatorname{Sym}(\mathcal E)\otimes_{\mathcal O_S}\operatorname{Sym}(\mathcal E).
\]

    \item \textbf{Construct the Hopf comultiplication.}
Applying the free-algebra functor to the diagonal
\[
  \Delta_{\mathcal E}:\mathcal E\to \mathcal E\oplus \mathcal E
\]
therefore produces a comultiplication
\[
  \operatorname{Sym}(\mathcal E)
  \longrightarrow
  \operatorname{Sym}(\mathcal E)\otimes_{\mathcal O_S}\operatorname{Sym}(\mathcal E).
\]

    \item \textbf{Identify the comultiplication on generators.}
By the free universal property it is characterized on the generating module
\(\mathcal E\) by
\[
  v\longmapsto v\otimes1+1\otimes v.
\]

    \item \textbf{Construct the remaining Hopf structure maps.}
Likewise, the zero map
\[
  \mathcal E\to0
\]
induces the augmentation
\[
  \operatorname{Sym}(\mathcal E)\longrightarrow\mathcal O_S,
\]
and the negation map
\[
  -\id_{\mathcal E}:\mathcal E\to \mathcal E
\]
induces the antipode
\[
  \operatorname{Sym}(\mathcal E)\longrightarrow\operatorname{Sym}(\mathcal E).
\]

    \item \textbf{Verify the Hopf identities.}
The commutative Hopf identities reduce, by the free-algebra universal property,
to the corresponding identities for diagonal, zero, and negation in the
additive commutative group object \(\mathcal E\) of \(\QCoh(S)\).

    \item \textbf{Pass to relative spectrum.}
Passing contravariantly through relative spectrum gives multiplication, zero
section, and inversion on
\[
  \mathbb L_S(\mathcal E).
\]
Thus \(\mathbb L_S(\mathcal E)\) is canonically a commutative group object over \(S\).

    \item \textbf{Establish functoriality.}
Now let
\[
  \varphi:\mathcal E\to \mathcal E'
\]
be a morphism of finite locally free quasi-coherent modules. Functoriality of
the free commutative-algebra construction gives
\[
  \operatorname{Sym}(\varphi):
  \operatorname{Sym}(\mathcal E)\longrightarrow\operatorname{Sym}(\mathcal E').
\]
The naturality of diagonal, zero, and negation makes this a morphism of the
displayed commutative Hopf algebras. Applying relative spectrum reverses its
variance and therefore gives a homomorphism of commutative group objects
\[
  \mathbb L_S(\mathcal E')
  \longrightarrow
  \mathbb L_S(\mathcal E).
\]
Hence the construction is contravariantly functorial in \(\mathcal E\).

    \item \textbf{Establish base change.}
Finally, pullback of quasi-coherent modules preserves direct sums, diagonals,
zero maps, negation, and the free commutative-algebra construction, while
relative spectrum commutes with base change. Therefore the entire
commutative-group-object structure and its contravariant functoriality are
compatible with arbitrary base change on \(S\).
\end{enumerate}
\end{proof}
\begin{proposition}[Linear translation presentation]
\label{prop:linear-translation-presentation}
Let \(\mathcal E\) be a finite locally free quasi-coherent module on \(S\).
For every \(n\geq0\),
\[
  \mathcal R^{\mathrm{dR}}_{\mathbb L_S(\mathcal E)/S,n}
  \simeq
  \mathbb L_S(\mathcal E)\times_S(\operatorname{Nbh}^{\mathrm{dR}}_{S/S}(\mathbb L_S(\mathcal E)))^{\times_S n},
\]
compatibly with the full simplicial structure. The source map is projection to
\(\mathbb L_S(\mathcal E)\), evaluation is formal addition, and
\[
  \mathsf{Disk}^{\mathrm{dR},\infty}_{\mathbb L_S(\mathcal E)/S}
  \simeq
  \Sigma_{\operatorname{pr}_{\mathbb L}}\alpha_{\mathbb L_S(\mathcal E)}^*,
  \qquad
  J^\infty_{\mathbb L_S(\mathcal E)/S}
  \simeq
  \Pi_{\operatorname{pr}_{\mathbb L}}\alpha_{\mathbb L_S(\mathcal E)}^*.
\]
Writing
\[
  \pi_{\mathcal E}:\mathbb L_S(\mathcal E)\to S
\]
for the structure morphism, the represented free-algebra cotangent calculation gives
\[
  L_{\mathbb L_S(\mathcal E)/S}
  \simeq
  \pi_{\mathcal E}^*\mathcal E.
\]
Thus the first-order linearization has the expected coefficient, while the full
relation retains all finite infinitesimal information.
\end{proposition}

\begin{proof}
\leavevmode
\begin{enumerate}[label=\textup{(\arabic*)}]
    \item \textbf{Apply the translation nerve.}
Theorem~\ref{thm:formal-translation-nerve-explicit} applies to the additive
group object \(\mathbb L_S(\mathcal E)\). The jet formulas follow by transport of the
relation correspondence.

    \item \textbf{Recover the cotangent comparison.}
The cotangent formula is the relative free
commutative-algebra calculation established in the differentiated spectral
doctrine. It supplies the first-order comparison with the structured
neighbourhood, while the de Rham relation supplies the jet comonad.
\end{enumerate}
\end{proof}

The translation nerve therefore gives a concrete checksum: source is projection,
target is formal translation, and higher simplices record successive
infinitesimal increments. Relative linear spaces recover the expected
cotangent coefficient at first order. A formally \(\acute{e}\)tale chart
transports this same model to a represented local context.

\subsection{Represented formally \texorpdfstring{$\acute{e}$}{e}tale local models}
\label{sec:represented-jet-interface}

A formally \(\acute{e}\)tale chart from a represented smooth context into affine
space pulls back the affine translation relation independently of any
group-object structure on the context. The following proposition is conditional only on the existence of
such a chart. Later geometric results may supply charts on larger theorem
domains; they are not inputs to the intrinsic jet construction.
\begin{proposition}[Formally \(\acute{e}\)tale translation action and nerve]
\label{prop:formally-etale-translation-action}
Suppose a later geometric theorem supplies a formally \(\acute{e}\)tale
morphism
\[
  q:U\to\mathbb A^n_V
\]
over a represented base \(V\). Using the relative de Rham neighbourhood
\[
  \operatorname{Nbh}^{\mathrm{dR}}_{V/V}(\mathbb A^n_V)
  \simeq
  \mathbb A^n_V
  \times_{(\mathbb A^n_V)_{\mathrm{dR}/V}}V
\]
of the zero section, there is a canonical action
\[
  \alpha_U:
  U\times_V\operatorname{Nbh}^{\mathrm{dR}}_{V/V}(\mathbb A^n_V)\longrightarrow U
\]
constructed from the evaluation map of the intrinsic relation of \(U/V\).
It satisfies the canonically commutative square
\[
\begin{tikzcd}
U\times_V\operatorname{Nbh}^{\mathrm{dR}}_{V/V}(\mathbb A^n_V)
  \ar[r,"\alpha_U"] \ar[d,"q\times\id"']
  & U \ar[d,"q"]\\
\mathbb A^n_V\times_V\operatorname{Nbh}^{\mathrm{dR}}_{V/V}(\mathbb A^n_V)
  \ar[r,"\alpha_{\mathbb A^n_V}"']
  & \mathbb A^n_V.
\end{tikzcd}
\]  Its unit and associativity equivalences are derived from
the unit and composition of the de Rham relation groupoid:
\[
  \alpha_U(u,0)\simeq u,
\]
\[
  \alpha_U(\alpha_U(u,\xi_1),\xi_2)
  \simeq
  \alpha_U(u,\xi_1+\xi_2).
\]
Relation inversion gives
\[
  \alpha_U(\alpha_U(u,\xi),-\xi)\simeq u.
\]

For every \(k\geq0\), there is a canonical equivalence
\[
  \mathcal R^{\mathrm{dR}}_{U/V,k}
  \simeq
  U\times_V(\operatorname{Nbh}^{\mathrm{dR}}_{V/V}(\mathbb A^n_V))^{\times_V k}
\]
compatible with the full simplicial structure. In these coordinates,
\[
\begin{aligned}
  d_0(u,\xi_1,\ldots,\xi_k)
  &=
  (\alpha_U(u,\xi_1),\xi_2,\ldots,\xi_k),\\
  d_i(u,\xi_1,\ldots,\xi_k)
  &=
  (u,\xi_1,\ldots,\xi_i+\xi_{i+1},\ldots,\xi_k),
  \quad 0<i<k,\\
  d_k(u,\xi_1,\ldots,\xi_k)
  &=
  (u,\xi_1,\ldots,\xi_{k-1}),
\end{aligned}
\]
and degeneracies insert \(0\). Consequently
\[
  \mathsf{Disk}^{\mathrm{dR},\infty}_{U/V}
  \simeq
  \Sigma_{\operatorname{pr}_U}\alpha_U^*,
  \qquad
  J^\infty_{U/V}
  \simeq
  \Pi_{\operatorname{pr}_U}\alpha_U^*.
\]
\end{proposition}

\begin{proof}
\leavevmode
\begin{enumerate}[label=\textup{(\arabic*)}]
    \item \textbf{Apply formally étale invariance.}
Formal-\(\acute{e}\)tale invariance of the relation gives
\[
  \mathcal R^{\mathrm{dR}}_{U/V,k}
  \simeq
  U\times_{\mathbb A^n_V,v_k}\mathcal R^{\mathrm{dR}}_{\mathbb A^n_V/V,k}
\]
simplicially in \(k\).

    \item \textbf{Insert the affine translation nerve.}
By
Theorem~\ref{thm:formal-translation-nerve-explicit},
\[
  \mathcal R^{\mathrm{dR}}_{\mathbb A^n_V/V,k}
  \simeq
  \mathbb A^n_V\times_V(\operatorname{Nbh}^{\mathrm{dR}}_{V/V}(\mathbb A^n_V))^{\times_V k}.
\]

    \item \textbf{Pull back the translation nerve.}
Pullback along \(q:U\to\mathbb A^n_V\) therefore gives
\[
  \mathcal R^{\mathrm{dR}}_{U/V,k}
  \simeq
  U\times_V(\operatorname{Nbh}^{\mathrm{dR}}_{V/V}(\mathbb A^n_V))^{\times_V k}
\]
for every \(k\).

    \item \textbf{Identify the source map in degree one.}
In degree one, under
\[
  \mathcal R^{\mathrm{dR}}_{U/V}
  \simeq
  U\times_V\operatorname{Nbh}^{\mathrm{dR}}_{V/V}(\mathbb A^n_V),
\]
the source map becomes the projection
\[
  \operatorname{pr}_U:U\times_V\operatorname{Nbh}^{\mathrm{dR}}_{V/V}(\mathbb A^n_V)\to U.
\]

    \item \textbf{Construct the induced translation action.}
Transport the evaluation map
\[
  \mathsf t:\mathcal R^{\mathrm{dR}}_{U/V}\to U
\]
through this equivalence and define the resulting map to be
\[
  \alpha_U:
  U\times_V\operatorname{Nbh}^{\mathrm{dR}}_{V/V}(\mathbb A^n_V)\to U.
\]

    \item \textbf{Compare with affine translation.}
The formal-\(\acute{e}\)tale relation square is Cartesian and intertwines
evaluation, so
\[
  q\alpha_U
  \simeq
  \alpha_{\mathbb A^n_V}(q\times\id_{\operatorname{Nbh}^{\mathrm{dR}}_{V/V}(\mathbb A^n_V)}).
\]

    \item \textbf{Identify the intrinsic action.}
Thus \(\alpha_U\) is the action induced by the intrinsic relation.

    \item \textbf{Derive the unit law.}
The relation unit \(U\to \mathcal R^{\mathrm{dR}}_{U/V}\) corresponds to the zero section
\[
  u\longmapsto(u,0),
\]
so the groupoid unit law gives
\[
  \alpha_U(u,0)\simeq u.
\]

    \item \textbf{Derive the composition law.}
In degree two,
\[
  \mathcal R^{\mathrm{dR}}_{U/V,2}
  \simeq
  U\times_V\operatorname{Nbh}^{\mathrm{dR}}_{V/V}(\mathbb A^n_V)\times_V\operatorname{Nbh}^{\mathrm{dR}}_{V/V}(\mathbb A^n_V).
\]
Under this identification relation composition is
\[
  (u,\xi_1,\xi_2)
  \longmapsto
  (u,\xi_1+\xi_2),
\]
while the successive evaluations are
\[
  u
  \longmapsto
  \alpha_U(u,\xi_1)
  \longmapsto
  \alpha_U(\alpha_U(u,\xi_1),\xi_2).
\]
The groupoid composition identity therefore gives
\[
  \alpha_U(\alpha_U(u,\xi_1),\xi_2)
  \simeq
  \alpha_U(u,\xi_1+\xi_2).
\]

    \item \textbf{Derive the inverse law.}
Relation inversion corresponds to
\[
  (u,\xi)
  \longmapsto
  (\alpha_U(u,\xi),-\xi),
\]
and the inverse law gives
\[
  \alpha_U(\alpha_U(u,\xi),-\xi)\simeq u.
\]

    \item \textbf{Recover all simplicial operators.}
Deleting and repeating vertices yields the displayed face and degeneracy maps
in every degree.

    \item \textbf{Read off the jet functors.}
Finally, under the degree-one relation equivalence the source and evaluation
maps are
\[
  \operatorname{pr}_U,
  \qquad
  \alpha_U,
\]
respectively. Proposition~\ref{prop:relation-jet-formulas} therefore gives
the formulas for \(\mathsf{Disk}^{\mathrm{dR},\infty}_{U/V}\) and \(J^\infty_{U/V}\).
\end{enumerate}
\end{proof}

\begin{remark}[Tubular comparison on represented smooth models]
On the represented smooth theorem domain, a coordinate-free tubular theorem
can compare the de Rham relation with a vector bundle constructed from the
cotangent complex. Such a theorem supplies a geometric model for the relation
and a first-order cotangent comparison. The jet comonad continues to arise
from the de Rham relation, so higher finite-infinitesimal information remains
visible in the resulting model.
\end{remark}

Represented coordinates therefore model rather than define the jet calculus.
With locality and the translation checksum in place, only two comparison
questions remain, and they must not be conflated.

\section{Comparison boundaries of the jet calculus}
\label{sec:jets-comparison-boundaries}

The first comparison changes the infinitesimal localization: when does the
reduced-classical relation produce the same infinite jets as the intrinsic
de Rham relation? The second keeps the intrinsic relation fixed but asks for a
functorial filtration by formal order. The first is therefore a convergence
problem. The second requires a new globalization construction.

The following fork records that distinction. Its right-hand branch is dashed
because no global finite-order relation stages have yet been constructed.
\[
\begin{tikzcd}[column sep=huge,row sep=large]
&
\mathcal R^{\mathrm{dR}}_{X/S,\bullet}
  \ar[dl,"{\text{further localization}}"']
  \ar[dr,dashed,"{\text{formal-order filtration}}"]
&
\\
\mathcal R^{\mathrm{red,cl}}_{X/S,\bullet}
  \ar[d]
&&
\{\mathcal R^{(k)}_{X/S,\bullet}\}_{k\geq0}
  \ar[d,dashed]
\\
J^{\infty,\mathrm{red,cl}}_{X/S}
  \ar[dr]
&&
\varprojlim_k J^k_{X/S}
\\
&
J^\infty_{X/S}
  \ar[ur,dashed,"{\text{tower comparison}}"']
&
\end{tikzcd}
\]
\subsection{Reduced-classical comparison}
\label{sec:reduced-jet-comparison}

The reduced-classical boundary is obtained by applying a further localization
to the target of the de Rham unit. The comparison of comonads is therefore
controlled by contraction of the intermediate dependent adjunction, not by a
new definition of jets.

The comparison \(\kappa_{X/S}\) of
Chapter~\ref{chap:spectral-de-Rham} determines the reduced-classical relation
simplicial object
\[
  \mathcal R^{\mathrm{red,cl}}_{X/S,\bullet}
  :=
  \check C_\bullet
  \left(
    X
    \xrightarrow{\ \kappa_{X/S}\eta_{X/S}^{\mathrm{dR}}\ }
    S\times_{\operatorname{Red}_{\mathrm{cl}}(S)}
    \operatorname{Red}_{\mathrm{cl}}(X)
  \right)
\]
and the reduced-classical jet comonad
\[
  J^{\infty,\mathrm{red,cl}}_{X/S}
  :=
  (\kappa_{X/S}\eta_{X/S}^{\mathrm{dR}})^*
  \Pi_{\kappa_{X/S}\eta_{X/S}^{\mathrm{dR}}}.
\]

\begin{remark}[Reduced-classical comparison comonad]
The comonad \(J^{\infty,\mathrm{red,cl}}_{X/S}\) is attached to the composite
\(\kappa_{X/S}\eta_{X/S}^{\mathrm{dR}}\) and serves as the reduced-classical
comparison object. On the theorem domain where \(\kappa_{X/S}\) is an
equivalence, it agrees with the intrinsic jet comonad.
\end{remark}

\begin{lemma}[Contraction of a composite base-change comonad]
\label{lem:composite-base-change-comonad}
Let
\[
  U\xrightarrow{f}V\xrightarrow{g}W
\]
be composable morphisms in an infinity-topos. Put
\[
  J_f:=f^*\Pi_f,
  \qquad
  J_{gf}:=(gf)^*\Pi_{gf}.
\]
Using composition of dependent pullbacks and products, there is a canonical
identification
\[
  J_{gf}
  \simeq
  f^*g^*\Pi_g\Pi_f.
\]
The counit of
\[
  g^*\dashv\Pi_g
\]
therefore induces a canonical natural transformation
\[
  \vartheta_{g,f}:
  J_{gf}
  \simeq
  f^*g^*\Pi_g\Pi_f
  \xrightarrow{\,f^*\varepsilon^g_{\Pi_f}\,}
  f^*\Pi_f
  =
  J_f.
\]
This transformation is a morphism of comonads, is compatible with further
composition, and is an equivalence whenever \(g\) is an equivalence.
\end{lemma}

\begin{proof}
\leavevmode
\begin{enumerate}[label=\textup{(\arabic*)}]
    \item \textbf{Identify the two adjunction comonads.}
The comonad \(J_f\) is induced by the adjunction
\[
  f^*\dashv\Pi_f,
\]
whereas \(J_{gf}\) is induced by the composite adjunction
\[
  f^*g^*\dashv\Pi_g\Pi_f.
\]

    \item \textbf{Contract the inserted adjunction.}
The displayed transformation contracts the inserted adjunction
\(g^*\dashv\Pi_g\) by its counit.

    \item \textbf{Verify comonad compatibility.}
Compatibility with the two comonad counits
and comultiplications is therefore the standard compatibility of units,
counits, and mates under composition of adjunctions. The same pasting
formalism gives compatibility with further composition.

    \item \textbf{Treat the equivalence case.}
If \(g\) is an equivalence, then \(g^*\) and \(\Pi_g\) are inverse
equivalences and the counit
\[
  g^*\Pi_g\to\id
\]
is invertible. Hence \(\vartheta_{g,f}\) is an equivalence of comonads.
\end{enumerate}
\end{proof}

\begin{theorem}[Canonical reduced-classical comparison of jets]
\label{thm:reduced-jet-comparison}
The map \(\kappa_{X/S}\) induces a simplicial map
\[
  \mathcal R^{\mathrm{dR}}_{X/S,\bullet}
  \longrightarrow
  \mathcal R^{\mathrm{red,cl}}_{X/S,\bullet}
\]
which is the identity in degree zero, and a canonical natural transformation
of comonads
\[
  J^{\infty,\mathrm{red,cl}}_{X/S}
  \longrightarrow
  J^\infty_{X/S}.
\]
If \(\kappa_{X/S}\) is an equivalence, this comparison is an equivalence of
comonads. In degree one the relation comparison sits in
\[
\begin{tikzcd}[column sep=huge,row sep=large]
\mathcal R^{\mathrm{dR}}_{X/S}
  \ar[r]
  \ar[d,shift left=.4ex,"\mathsf s"]
  \ar[d,shift right=.4ex,"\mathsf t"']
&
\mathcal R^{\mathrm{red,cl}}_{X/S}
  \ar[d,shift left=.4ex]
  \ar[d,shift right=.4ex]\\
X \ar[r,equal]
&
X .
\end{tikzcd}
\]
\end{theorem}

\begin{proof}
\leavevmode
\begin{enumerate}[label=\textup{(\arabic*)}]
    \item \textbf{Construct the simplicial comparison.}
The comparison is encoded by the commutative triangle
\[
\begin{tikzcd}[column sep=huge,row sep=large]
X \ar[rr,"\kappa_{X/S}\eta_{X/S}^{\mathrm{dR}}"]
  \ar[dr,"\eta_{X/S}^{\mathrm{dR}}"']
&&
S\times_{\operatorname{Red}_{\mathrm{cl}}(S)}
  \operatorname{Red}_{\mathrm{cl}}(X)\\
&
X_{\mathrm{dR}/S}
  \ar[ur,"\kappa_{X/S}"']
&
\end{tikzcd}
\]
and finite-limit functoriality gives the induced simplicial relation map.

    \item \textbf{Apply the composite-comonad lemma.}
Apply Lemma~\ref{lem:composite-base-change-comonad} with
\[
  f=\eta_{X/S}^{\mathrm{dR}},
  \qquad
  g=\kappa_{X/S}.
\]

    \item \textbf{Construct the comparison comonad morphism.}
Since
\[
  \Pi_{\kappa_{X/S}\eta_{X/S}^{\mathrm{dR}}}
  \simeq
  \Pi_{\kappa_{X/S}}
  \Pi_{\eta_{X/S}^{\mathrm{dR}}},
\]
the resulting comonad morphism is
\[
\begin{aligned}
  J^{\infty,\mathrm{red,cl}}_{X/S}
  &=
  (\kappa_{X/S}\eta_{X/S}^{\mathrm{dR}})^*
  \Pi_{\kappa_{X/S}\eta_{X/S}^{\mathrm{dR}}}\\
  &\simeq
  (\eta_{X/S}^{\mathrm{dR}})^*
  \kappa_{X/S}^*
  \Pi_{\kappa_{X/S}}
  \Pi_{\eta_{X/S}^{\mathrm{dR}}}\\
  &\longrightarrow
  (\eta_{X/S}^{\mathrm{dR}})^*
  \Pi_{\eta_{X/S}^{\mathrm{dR}}}\\
  &\simeq
  J^\infty_{X/S},
\end{aligned}
\]
where the displayed arrow is induced by the counit of
\[
  \kappa_{X/S}^*\dashv\Pi_{\kappa_{X/S}}.
\]

    \item \textbf{Verify comonad compatibility.}
Lemma~\ref{lem:composite-base-change-comonad} supplies compatibility with
counit and comultiplication.

    \item \textbf{Treat the equivalence case.}
If \(\kappa_{X/S}\) is an
equivalence, the same lemma makes the comparison an equivalence of comonads.
\end{enumerate}
\end{proof}

\begin{corollary}[Convergence criterion for the reduced-classical jet comparison]
If the absolute comparison maps
\[
  \kappa_X:
  X_{\mathrm{dR}}
  \longrightarrow
  \operatorname{Red}_{\mathrm{cl}}(X),
  \qquad
  \kappa_S:
  S_{\mathrm{dR}}
  \longrightarrow
  \operatorname{Red}_{\mathrm{cl}}(S)
\]
are equivalences, then the relative comparison
\[
  \kappa_{X/S}:X_{\mathrm{dR}/S}\to S\times_{\operatorname{Red}_{\mathrm{cl}}(S)}\operatorname{Red}_{\mathrm{cl}}(X)
\]
and hence
\[
  J^{\infty,\mathrm{red,cl}}_{X/S}
  \xrightarrow{\simeq}
  J^\infty_{X/S}
\]
are equivalences.

Moreover, the convergence theorem of
Chapter~\ref{chap:spectral-de-Rham} identifies invertibility of \(\kappa_X\)
with Postnikov continuity and nilradical continuity of the de Rham-local object
\(X_{\mathrm{dR}}\) on affine tests, and similarly for \(S\).
\end{corollary}

\begin{proof}
\leavevmode
\begin{enumerate}[label=\textup{(\arabic*)}]
    \item \textbf{Use the further localization.}
Chapter~\ref{chap:spectral-de-Rham} proves that
\(\operatorname{Red}_{\mathrm{cl}}\) is the further left-exact localization of
the de Rham-local sector characterized by Postnikov continuity and nilradical
continuity, and that
\[
  \operatorname{Red}_{\mathrm{cl}}\circ\operatorname{dR}
  \simeq
  \operatorname{Red}_{\mathrm{cl}}.
\]

    \item \textbf{Apply the convergence criterion to the de Rham object.}
Apply Theorem~\ref{thm:reduced-convergence} to the de Rham-local object
\[
  X_{\mathrm{dR}}.
\]
It is \(\operatorname{Red}_{\mathrm{cl}}\)-local if and only if it is local for
the Postnikov-continuity generators and nilradical-continuous on affine tests.

    \item \textbf{Identify the absolute comparison map.}
The unit of the further localization is
\[
  X_{\mathrm{dR}}
  \longrightarrow
  \operatorname{Red}_{\mathrm{cl}}(X_{\mathrm{dR}})
  \simeq
  \operatorname{Red}_{\mathrm{cl}}(X),
\]
which is precisely \(\kappa_X\). Hence \(\kappa_X\) is an equivalence exactly
under those two continuity conditions.

    \item \textbf{Apply the argument to the base.}
The same argument applies to \(S\).

    \item \textbf{Pass to the relative comparison.}
When both absolute comparisons are equivalences, substitute them into
\[
  X_{\mathrm{dR}/S}
  =
  S\times_{S_{\mathrm{dR}}}X_{\mathrm{dR}}.
\]
Finite-limit functoriality gives
\[
  \kappa_{X/S}:
  X_{\mathrm{dR}/S}
  \xrightarrow{\simeq}
  S\times_{\operatorname{Red}_{\mathrm{cl}}(S)}\operatorname{Red}_{\mathrm{cl}}(X).
\]
Theorem~\ref{thm:reduced-jet-comparison} then gives
\[
  J^{\infty,\mathrm{red,cl}}_{X/S}
  \xrightarrow{\simeq}
  J^\infty_{X/S}.
\]
\end{enumerate}
\end{proof}

\begin{remark}[Comparison with crystals and nonlinear \(D\)-geometry]
On the theorem domain where intrinsic de Rham reflection agrees with
reduced-classical reflection, the equivalence
\[
  \operatorname{PDE}^{\mathrm{fi}}_{X/S}
  \simeq
  (\mathcal H_{\Sp})_{/Q_{X/S}}
\]
is formally parallel to the description of crystals by descent over a de Rham
prestack and to nonlinear \(D\)-geometry
\cite{GaitsgoryRozenblyumCrystals,BeilinsonDrinfeld2004}. The comparison applies on the stated theorem domain and identifies the two
constructed architectures there. The intrinsic de Rham reflector is generated by the
\textbf{square-zero infinitesimal substitutions}; no
coefficient module, classifying derivation, or square-zero presentation is
retained by the reflector.
\end{remark}

\begin{remark}[Comparison with reduced synthetic jets]
On classical smooth input the preceding theorem gives the comparison with a
reduced de Rham-groupoid jet construction on its convergence domain. The
interaction of reduction with projective limits in synthetic smooth jet geometry
is studied in
\cite{GiotopoulosKhavkineSatiSchreiberReducedJets}.
\end{remark}

Reduced-classical comparison is thus a convergence question about a further
localization of an already constructed infinite-jet comonad. Finite order is a
different problem: it asks for internal stages of the same intrinsic relation.

\subsection{What is required for finite-order jets}
\label{sec:no-finite-order-jets}

Only \(J^\infty\) has been constructed globally. A global finite-jet construction
requires a functorial, base-change-coherent filtration of the de Rham relation
itself. Chosen finite factorizations encode presentation data through their
lengths and intermediate square-zero stages. The reconstruction chapter
supplies the intrinsic affine invariant---the canonical formal-order
filtration---whose coherent extension to the unrestricted simplicial relation
requires an additional theorem.

A finite-order refinement would therefore consist of a tower
\[
  J^0E\longleftarrow J^1E\longleftarrow J^2E\longleftarrow\cdots
\]
together with a comparison of \(J^\infty\) with its inverse limit. The desired
global relationship would have the schematic form
\[
\begin{tikzcd}[column sep=large,row sep=large]
J^\infty E \ar[d,dashed,"\text{comparison}"']\\
\cdots \ar[r]
& J^2E \ar[r]
& J^1E \ar[r]
& J^0E .
\end{tikzcd}
\]
The constructions already available determine what that missing theorem must
do.

First, Chapter~\ref{chap:differentiation-spectral-doctrine} defines finite
spectral infinitesimal substitutions as the equivalence-closed class generated
under finite composition by square-zero substitutions. Membership in that
class is intrinsic, but a chosen factorization and its length are not retained
structure.

Second, Chapter~\ref{chap:formal-reconstruction-spectral-doctrine} constructs
a genuine \textbf{formal-order filtration} on every augmented spectral algebra
from the adjunction
\[
  \operatorname{adic}\dashv\operatorname{aug}.
\]
This filtration supplies an intrinsic affine notion of order. A finite jet
tower would require a base-change-coherent globalization of these affine stages
to the full de Rham relation and, compatibly, to its simplicial structure.
Only after that construction would endofunctors \(J^k\) be available on the
unrestricted slice.

\begin{proposition}[Canonical formal order on an affine diagonal]
\label{prop:affine-diagonal-formal-order}
Let \(R\to A\) be a morphism of connective commutative ring spectra and let
\[
  \mu_A:A\otimes_RA\longrightarrow A
\]
be multiplication, the coordinate augmentation of the affine diagonal. The
formal-order reconstruction of
Chapter~\ref{chap:formal-reconstruction-spectral-doctrine} supplies the
canonical completed adic filtration
\[
  \widehat{\operatorname{adic}}(A\otimes_RA\to A).
\]
The cotangent comparison is displayed by
\[
\begin{tikzcd}[column sep=huge,row sep=large]
L_{A/(A\otimes_RA)}[-1]
  \ar[r]
&
L_{(A\otimes_RA)/R}\otimes_{A\otimes_RA}A
  \ar[r]
&
L_{A/R},
\end{tikzcd}
\]
and in particular there is a canonical equivalence
\[
  L_{A/(A\otimes_RA)}[-1]
  \xrightarrow{\simeq}
  L_{A/R}.
\]
Hence there is a canonical equivalence of graded commutative \(A\)-algebras
\[
  \operatorname{Gr}
  \widehat{\operatorname{adic}}(A\otimes_RA\to A)
  \simeq
  \operatorname{Sym}^{\mathrm{gr}}_A
  \bigl(L_{A/R}\langle1\rangle\bigr).
\]
\end{proposition}

\begin{proof}
\leavevmode
\begin{enumerate}[label=\textup{(\arabic*)}]
    \item \textbf{Apply cotangent transitivity.}
Apply cotangent transitivity
\cite{LurieHigherAlgebra,LurieSAG} to
\[
  R\longrightarrow A\otimes_RA\xrightarrow{\mu_A}A.
\]
It gives a cofibre sequence
\[
  L_{(A\otimes_RA)/R}\otimes_{A\otimes_RA}A
  \longrightarrow
  L_{A/R}
  \longrightarrow
  L_{A/(A\otimes_RA)}.
\]

    \item \textbf{Identify the first cotangent term.}
Cotangent base change for the two tensor factors identifies the first term with
\[
  L_{A/R}\oplus L_{A/R},
\]
and the map induced by multiplication is addition
\[
  +:L_{A/R}\oplus L_{A/R}\longrightarrow L_{A/R}.
\]

    \item \textbf{Compute the fibre of addition.}
Its fibre is canonically \(L_{A/R}\), via
\[
  x\longmapsto(x,-x).
\]

    \item \textbf{Identify the shifted diagonal cotangent complex.}
Stability therefore identifies its cofibre with \(L_{A/R}[1]\), proving
\[
  L_{A/(A\otimes_RA)}[-1]
  \simeq
  L_{A/R}.
\]

    \item \textbf{Apply the associated-graded theorem.}
Now apply the associated-graded theorem for the canonical adic filtration,
\Ref{thm:adic-associated-graded}, to the augmentation
\[
  A\otimes_RA\to A.
\]
It gives
\[
\begin{aligned}
  \operatorname{Gr}
  \widehat{\operatorname{adic}}(A\otimes_RA\to A)
  &\simeq
  \operatorname{Sym}^{\mathrm{gr}}_A
  \bigl(
    L_{A/(A\otimes_RA)}[-1]\langle1\rangle
  \bigr)\\
  &\simeq
  \operatorname{Sym}^{\mathrm{gr}}_A
  \bigl(L_{A/R}\langle1\rangle\bigr).
\end{aligned}
\]
\end{enumerate}
\end{proof}

\begin{remark}[Globalization required for finite-order jets]
The preceding proposition supplies formal-order stages for a represented
affine diagonal augmentation. A global finite-order jet theory requires a
base-change coherent extension of those stages to the topos-level relation
\[
  \mathcal R^{\mathrm{dR}}_{X/S,1}
  =
  X\times_{Q_{X/S}}X,
\]
together with a compatible simplicial filtration of
\(\mathcal R^{\mathrm{dR}}_{X/S,\bullet}\) and endofunctors
\[
  J^k_{X/S}:(\mathcal H_{\Sp})_{/X}\to(\mathcal H_{\Sp})_{/X}
\]
whose inverse limit compares with \(J^\infty_{X/S}\).
The complete totalization filtration of
Chapter~\ref{chap:spectral-de-Rham} remains the descent filtration on
de Rham cochains; its normalized layers supply the intrinsic de Rham forms.
\end{remark}

\begin{remark}[Finite dependence requires a constructed invariant]
\label{prop:finite-dependence-intrinsic}
Let
\[
  D:E\rightsquigarrow F
\]
be a differential operator represented by
\[
  \widehat D:J^\infty E\to F.
\]
The present construction supplies a canonical order-zero class through ordinary
morphisms
\[
  E\to F,
\]
represented in the co-Kleisli category by
\[
  J^\infty E\xrightarrow{\varepsilon^\Pi_E}E\to F.
\]
Beyond order zero, a numerical order function on arbitrary differential
operators requires a functorial invariant that measures finite dependence, for
example a filtration stage, a compactness condition, or a factorization
through a constructed finite-order jet object
\[
  J^kE.
\]
The finite-infinitesimal calculus by itself records membership in the
equivalence-closed finite-composition class generated by square-zero
substitutions; a chosen factorization length remains presentation data rather
than an intrinsic order. The canonical adic filtration of
Chapter~\ref{chap:formal-reconstruction-spectral-doctrine} supplies an
intrinsic affine formal-order invariant, and a numerical jet order on the
unrestricted slice becomes available after that filtration has been
globalized to a functorial tower \(J^k\).
\end{remark}

The boundary is therefore exact: affine diagonals carry canonical formal order
with cotangent-controlled associated graded, while the global simplicial
filtration needed for \(J^k\) remains a further construction. The infinite-jet
calculus is independent of that refinement, and numerical order becomes
intrinsic after the filtration has been globalized.

\section{The crystalline jet interface}
\label{sec:crystalline-jet-package}

The chapter now has a reusable interface. The theorem below is intentionally a
citation index: it collects the precise results that later chapters may invoke.
The retrospective summary that follows has a different job and will compress
these results to the conceptual structure that should remain active.
\begin{theorem}[Crystalline jet calculus]
For every \(X\to S\), the effective de Rham quotient
\[
  X\xrightarrow{e_{X/S}}Q_{X/S}
\]
and its de Rham relation groupoid organize the jet calculus through the
commutative architecture
\[
\begin{tikzcd}[column sep=large,row sep=large]
&
\mathcal R^{\mathrm{dR}}_{X/S,\bullet}
  \ar[dl]
  \ar[dr]
&\\
\mathsf{Disk}^{\mathrm{dR},\infty}_{X/S}
  \ar[rr,"\dashv" description]
&&
J^\infty_{X/S}\\
&
(\mathcal H_{\Sp})_{/Q_{X/S}}
  \ar[ul]
  \ar[ur]
  \ar[d]
&\\
&
\operatorname{PDE}^{\mathrm{fi}}_{X/S}
&
\end{tikzcd}
\]
and canonically determine:
\begin{enumerate}
\item relative de Rham neighbourhoods serving as de Rham formal disks
\[
  \operatorname{Nbh}^{\mathrm{dR}}_{T/S}(T\times_SX)
  \simeq
  T\times_{Q_{X/S}}X
\]
for every generalized point \(a:T\to X\);

\item the adjoint triple
\[
  \Sigma_{e_{X/S}}
  \dashv
  e_{X/S}^*
  \dashv
  \Pi_{e_{X/S}}
\]
and hence the monad/comonad adjunction
\[
  \mathsf{Disk}^{\mathrm{dR},\infty}_{X/S}
  =
  e_{X/S}^*\Sigma_{e_{X/S}}
  \dashv
  J^\infty_{X/S}
  =
  e_{X/S}^*\Pi_{e_{X/S}};
\]

\item the mate equivalence
\[
  \operatorname{Alg}_{\mathsf{Disk}^{\mathrm{dR},\infty}}((\mathcal H_{\Sp})_{/X})
  \xrightarrow{\simeq}
  \operatorname{Coalg}_{J^\infty}((\mathcal H_{\Sp})_{/X}),
\]
together with the effective-descent recognition
\[
  \operatorname{PDE}^{\mathrm{fi}}_{X/S}
  =
  (\mathcal H_{\Sp})_{/Q_{X/S}}
  \simeq
  \operatorname{Alg}_{\mathsf{Disk}^{\mathrm{dR},\infty}}((\mathcal H_{\Sp})_{/X})
  \simeq
  \operatorname{Coalg}_{J^\infty}((\mathcal H_{\Sp})_{/X});
\]

\item crystalline transport as the relation transpose of the jet coaction,
with its unit, inverse, cocycle, and simplicial compatibility induced by the de Rham relation groupoid;

\item the co-Kleisli differential-operator calculus and infinite prolongation;

\item the comma infinity-category
\[
  \operatorname{Eqn}_{X/S}
  =
  (\id_{(\mathcal H_{\Sp})_{/X}}\downarrow J^\infty_{X/S})
\]
of generalized jet equations, together with derived equations, actual and
parameterized solution objects, formal-solution spaces, Cartan prolongation,
the coalgebraic representation of actual solutions, and prolongation-completeness;

\item arbitrary base change and effective formally \(\acute{e}\)tale descent
of the jet calculus and equation infinity-category, together with the
conditional reduced-classical comonad comparison;

\item for a relative commutative group object, the formal-translation model
and full translation nerve, and for a formally \(\acute{e}\)tale affine-space
chart, the induced translation action constructed from relation evaluation.
\end{enumerate}

On represented affine diagonal augmentations, the canonical adic construction
of Chapter~\ref{chap:formal-reconstruction-spectral-doctrine} supplies genuine
formal order with associated graded
\[
  \operatorname{Sym}^{\mathrm{gr}}_A
  \bigl(L_{A/R}\langle1\rangle\bigr),
\]
and a global finite-order refinement would further supply a tower \(J^k\)
and a comparison
\[
  J^\infty\longrightarrow\varprojlim\nolimits_kJ^k.
\]
Such a comparison belongs to the additional globalization problem described
above.
\end{theorem}

\begin{proof}
\leavevmode
\begin{enumerate}[label=\textup{(\arabic*)}]
    \item \textbf{Record the formal-disk theorem.}
The formal-disk statement is
\Ref{thm:formal-disk-graph-neighbourhood}.

    \item \textbf{Record the disk--jet adjunction data.}
The effective-image formulas, jet
comonad, and adjoint-triple mate structure are
\Ref{thm:effective-image-jets},
\Ref{thm:jet-comonad}, and
\Ref{thm:adjoint-triple-mates}.

    \item \textbf{Record the Eilenberg--Moore and descent recognition.}
The Eilenberg--Moore mate equivalence is
\Ref{thm:mate-em-equivalence}, while effective \v{C}ech descent and
Barr--Beck recognition are
\Ref{prop:effective-cech-descent-Q} and
\Ref{thm:effective-jet-descent}.

    \item \textbf{Record crystalline transport.}
Crystalline transport, its cocycle, and its
inverse are constructed in
\Ref{prop:jet-coaction-relation-transport}.

    \item \textbf{Record holonomic and Cartan recognition.}
Holonomic and Cartan recognition are
\Ref{thm:holonomic-recognition} and
\Ref{thm:cartan-recognition}.

    \item \textbf{Record the equation infinity-category.}
The equation infinity-category is
Construction~\ref{con:equation-category}.

    \item \textbf{Record the prolongation and solution representation.}
Cartan prolongation,
formal-solution representability, and actual-solution representation are
\Ref{thm:cartan-prolongation-adjunction},
\Ref{thm:universal-formal-solution}, and
\Ref{thm:actual-solutions-prolongation}.

    \item \textbf{Record the prolongation-completeness criteria.}
Prolongation-completeness is controlled by
\Ref{thm:prolongation-complete-formal-solutions},
\Ref{thm:prolongation-cartesian},
\Ref{thm:tautological-prolongation-completeness}, and
\Ref{thm:prolongation-complete-intrinsic-structure}.

    \item \textbf{Record the functoriality and reduced-comparison results.}
Base change, formally \(\acute{e}\)tale restriction, effective descent of
the jet calculus, equation descent, and the reduced comparison are
\Ref{thm:jet-base-change},
\Ref{thm:formal-etale-jets},
\Ref{thm:formally-etale-descent-jets},
\Ref{cor:formally-etale-equation-descent}, and
\Ref{thm:reduced-jet-comparison}.

    \item \textbf{Record translation and affine formal-order comparisons.}
The group-object translation specialization
is \Ref{thm:formal-translation-nerve-explicit}; the induced translation action
on a formally \(\acute{e}\)tale affine-space chart is
\Ref{prop:formally-etale-translation-action}; and the affine formal-order
comparison is \Ref{prop:affine-diagonal-formal-order}.
\end{enumerate}
\end{proof}

\subsection{Chapter summary}
\label{sec:crystalline-jets-summary}

The chapter has converted one effective infinitesimal relation into three
usable forms. Over \(X\), it acts as the formal-disk monad and infinite-jet
comonad; over \(Q_{X/S}\), the same data are a single descended object.
Crystalline transport is the relation-theoretic expression of that
identification, not an extra field attached to an object.

The jet presentation then supplies the differential-equation calculus.
Differential operators are co-Kleisli maps out of infinite jets, while
generalized equations are maps into them. Cartan prolongation is the universal
passage from an arbitrary ambient equation to one carrying compatible
crystalline transport. Its evaluation
\(c_i:\operatorname{Prol}(i)\to E\) is the recognition map: when it is
invertible, the compatible Cartan structure has contractible moduli and the
equation is already an intrinsic formally integrable PDE.

Two limits of the construction should remain separate. Reduced-classical
comparison asks whether a further infinitesimal localization changes the
already-defined \(J^\infty\), and is controlled by convergence. Finite order
asks for new structure on the intrinsic relation: a global simplicial
formal-order filtration. The affine adic filtration supplies the local
invariant, but no global \(J^k\)-tower is asserted or used before that
globalization is constructed.
